% Le microscope — Physique 1ère E — Cours signé OnBuch+
% Programme officiel MINESEC. Compiler avec : tectonic N09-microscope.tex
\newcommand{\DOCMATIERE}{Physique}
\newcommand{\DOCNIVEAU}{1ère E}
\newcommand{\DOCMODULE}{Module 3 : Optique}
\newcommand{\DOCLECON}{N09}
\newcommand{\DOCTITRE}{Le microscope}
% OnBuch+ — préambule « cours signés » ENRICHI (compilé avec tectonic / XeLaTeX)
% Système visuel « Pop » d'OnBuch+ : fond crème (jamais blanc), bordures encre
% épaisses, ombres « dures » décalées, titres Archivo Black en capitales.
% Version MATHÉMATIQUES : graphes (pgfplots/tikz), boîtes pédagogiques
% (définition, propriété, méthode, attention, le savais-tu, exercice résolu,
% exercice, TP, bilan), page de garde et signature OnBuch+ ancrée en bas.
%
% Macros attendues dans main.tex AVANT \input{preamble} :
%   \DOCMATIERE  \DOCNIVEAU  \DOCMODULE  \DOCLECON (numéro)  \DOCTITRE
\documentclass[11pt]{article}

\usepackage{fontspec}
\usepackage[french]{babel}
\usepackage[a4paper,margin=2cm,top=3.4cm,bottom=2.8cm,headheight=1.4cm,footskip=1.3cm]{geometry}
\usepackage{amsmath,amssymb,amsfonts}
\usepackage{mathtools} % \xmapsto, \coloneqq... souvent utilisés par les modèles
\usepackage{cancel} % simplifications barrées (bilans de forces)
% \abs{z} (module, valeur absolue) et \norm{u} (norme), utilisés par les modèles
\providecommand{\abs}{}\let\abs\relax\DeclarePairedDelimiter{\abs}{\lvert}{\rvert}
\providecommand{\norm}{}\let\norm\relax\DeclarePairedDelimiter{\norm}{\lVert}{\rVert}
\usepackage[table]{xcolor} % [table] : fournit aussi \rowcolors (alternance de lignes)
\usepackage{pagecolor}
\usepackage{tikz}
\usetikzlibrary{arrows.meta,positioning,calc,shapes.geometric,shapes.misc,shapes.arrows,decorations.pathmorphing,decorations.pathreplacing,decorations.markings,patterns,shadows,mindmap,trees,fit,backgrounds,matrix,intersections,angles,quotes}
\usepackage{pgfplots}
\usepackage[version=4]{mhchem} % équations nucléaires \ce{^{235}_{92}U}
\usepackage[european,siunitx]{circuitikz} % schémas électriques
\pgfplotsset{compat=1.18}
\usepgfplotslibrary{fillbetween,groupplots} % aires entre courbes (calcul intégral)
\usepackage{enumitem}
\usepackage{fancyhdr}
% [most] charge toutes les bibliothèques tcolorbox (listings, minted, raster,
% documentation...) et gonfle inutilement la mémoire TeX sur un document long
% (~300 boîtes) : on ne charge que ce qui est réellement utilisé.
\usepackage[skins,breakable]{tcolorbox}
\usepackage{graphicx}
\usepackage{colortbl}
\usepackage{booktabs}
\usepackage{tabularx}
\usepackage{diagbox} % case d'en-tête diagonale des tableaux à double entrée
% Colonnes extensibles centrée / gauche / droite (C, L, R), souvent utilisées par les modèles
\newcolumntype{C}{>{\centering\arraybackslash}X}
\newcolumntype{L}{>{\raggedright\arraybackslash}X}
\newcolumntype{R}{>{\raggedleft\arraybackslash}X}
\usepackage{array}
\usepackage{multicol}
\usepackage{siunitx}
% parse-numbers=false : accepte aussi \SI{2,5 \times 10^{-3}}{...}
\sisetup{locale=FR,output-decimal-marker={,},group-minimum-digits=4,parse-numbers=false}
\DeclareSIUnit\ohm{\ensuremath{\symup{\Omega}}} % U+2126 absent de la police
\DeclareSIUnit\division{div} % réglages d'oscilloscope (ms/div, V/div)
\DeclareSIUnit\tr{tr} % tours (vitesse de rotation, ex. \SI{1500}{\tr\per\minute})
\DeclareSIUnit\ch{ch} % cheval-vapeur (puissance moteur, unité usuelle hors SI)
\DeclareSIUnit\franccfa{FCFA} % franc CFA (coûts, contexte camerounais)
\DeclareSIUnit\dioptre{\ensuremath{\delta}} % vergence d'une lentille (dioptrie)
\usepackage{qrcode}
% \degree : commande fréquemment utilisée par les modèles pour un angle ou une
% température en dehors de \SI{}{} (non fournie par les paquets chargés).
\providecommand{\degree}{^{\circ}}
% \qed (fin de démonstration), utilisé par les modèles sans amsthm
\providecommand{\qed}{\hfill$\square$}
% \barre : double barre des valeurs interdites dans un tableau de variations (tkz-tab)
\providecommand{\barre}{\ensuremath{\|}}
% environnement proof (amsthm non chargé) utilisé par les modèles
\ifdefined\proof\else\newenvironment{proof}[1][Démonstration]{\par\noindent\textit{#1.}\ }{\qed\par}\fi
\usepackage{lastpage}
\usepackage{hyperref}
\hypersetup{hidelinks}

% ============================================================
% Palette « Pop » OnBuch+ — mêmes hex que la classe P (plus_ui.dart)
% ============================================================
\definecolor{poporange}{HTML}{F59321}
\definecolor{popdark}{HTML}{E07A0C}
\definecolor{popcream}{HTML}{FFF6E8}
\definecolor{popink}{HTML}{1C1714}
\definecolor{poppurple}{HTML}{7A5AE0}
\definecolor{popgreen}{HTML}{2E9E6A}
\definecolor{poppink}{HTML}{E0457B}
\definecolor{popblue}{HTML}{2F7DE1}
\definecolor{popgold}{HTML}{C98A12}
\definecolor{popmuted}{HTML}{8A7F76}
\definecolor{popline}{HTML}{EADFD2}
\definecolor{popgray}{HTML}{8A7F76}
\colorlet{popgrey}{popgray}
% Alias tolérants (le modèle nomme parfois les couleurs différemment)
\colorlet{popwhite}{white}
\colorlet{popblack}{popink}
\colorlet{popred}{poppink}
\colorlet{popyellow}{popgold}
% Teintes claires pour fonds de tableaux / zones
\colorlet{popblueL}{popblue!10!white}
\colorlet{poporangeL}{poporange!14!white}
\colorlet{popgreenL}{popgreen!12!white}
\colorlet{poppurpleL}{poppurple!12!white}
\colorlet{poppinkL}{poppink!12!white}
\colorlet{popgoldL}{popgold!14!white}

\pagecolor{popcream}

% ============================================================
% Typographie
% ============================================================
\defaultfontfeatures{Ligatures=TeX}
\setmainfont{PlusJakartaSans-Medium.ttf}[Path=../../../fonts/,BoldFont=PlusJakartaSans-Bold.ttf]
% Maths en sans-serif (Fira Math) avec chiffres et lettres droites de Plus
% Jakarta, pour que les formules se fondent dans le texte.
\usepackage[math-style=ISO,bold-style=ISO]{unicode-math}
\setmathfont{FiraMath-Regular.otf}[Path=../../../fonts/]
\setmathfont{PlusJakartaSans-Medium.ttf}[Path=../../../fonts/,range={up/num,up/{Latin,latin},\mathup}]
\usepackage{icomma} % virgule décimale sans espace en mode math
% Caractères mathématiques Unicode absents des polices de texte (Plus Jakarta,
% Archivo Black) : rendus par leur équivalent mathématique, en texte comme en
% formule — sinon ils s'affichent en carré vide (ex. « Arithmétique dans ℤ »).
\usepackage{newunicodechar}
\newunicodechar{ℝ}{\ensuremath{\mathbb{R}}}
\newunicodechar{ℤ}{\ensuremath{\mathbb{Z}}}
\newunicodechar{ℂ}{\ensuremath{\mathbb{C}}}
\newunicodechar{ℕ}{\ensuremath{\mathbb{N}}}
\newunicodechar{ℚ}{\ensuremath{\mathbb{Q}}}
\newunicodechar{𝔻}{\ensuremath{\mathbb{D}}}
\newunicodechar{α}{\ensuremath{\alpha}}
\newunicodechar{β}{\ensuremath{\beta}}
\newunicodechar{γ}{\ensuremath{\gamma}}
\newunicodechar{δ}{\ensuremath{\delta}}
\newunicodechar{ε}{\ensuremath{\varepsilon}}
\newunicodechar{θ}{\ensuremath{\theta}}
\newunicodechar{λ}{\ensuremath{\lambda}}
\newunicodechar{μ}{\ensuremath{\mu}}
\newunicodechar{σ}{\ensuremath{\sigma}}
\newunicodechar{τ}{\ensuremath{\tau}}
\newunicodechar{φ}{\ensuremath{\varphi}}
\newunicodechar{ψ}{\ensuremath{\psi}}
\newunicodechar{ω}{\ensuremath{\omega}}
\newunicodechar{Σ}{\ensuremath{\Sigma}}
% majuscules produites par \MakeUppercase dans les titres (α-aminés -> Α)
\newunicodechar{Α}{\ensuremath{\alpha}}
\newunicodechar{Β}{\ensuremath{\beta}}
\newunicodechar{Γ}{\ensuremath{\Gamma}}
\newunicodechar{Δ}{\ensuremath{\Delta}}
\newunicodechar{Ε}{\ensuremath{\varepsilon}}
\newunicodechar{Θ}{\ensuremath{\Theta}}
\newunicodechar{Λ}{\ensuremath{\Lambda}}
\newunicodechar{Μ}{\ensuremath{\mu}}
\newunicodechar{Τ}{\ensuremath{\tau}}
\newunicodechar{Φ}{\ensuremath{\Phi}}
\newunicodechar{Ψ}{\ensuremath{\Psi}}
\newunicodechar{Ω}{\ensuremath{\Omega}}
\newunicodechar{↦}{\ensuremath{\mapsto}}
\newunicodechar{∈}{\ensuremath{\in}}
\newunicodechar{∉}{\ensuremath{\notin}}
\newunicodechar{∧}{\ensuremath{\wedge}}
\newunicodechar{⇔}{\ensuremath{\Leftrightarrow}}
\newunicodechar{⟺}{\ensuremath{\Longleftrightarrow}}
\newunicodechar{⇒}{\ensuremath{\Rightarrow}}
\newunicodechar{⟹}{\ensuremath{\Longrightarrow}}
\newunicodechar{∘}{\ensuremath{\circ}}
\newunicodechar{∪}{\ensuremath{\cup}}
\newunicodechar{∩}{\ensuremath{\cap}}
\newunicodechar{≡}{\ensuremath{\equiv}}
\newunicodechar{⊂}{\ensuremath{\subset}}
\newunicodechar{⊥}{\ensuremath{\perp}}
\newunicodechar{∃}{\ensuremath{\exists}}
\newunicodechar{∀}{\ensuremath{\forall}}
\newunicodechar{∛}{\ensuremath{\sqrt[3]{\,}}}
\newunicodechar{∜}{\ensuremath{\sqrt[4]{\,}}}
\newunicodechar{⃗}{\ensuremath{{}^{\to}}}
\newunicodechar{ⁿ}{\ifmmode^{n}\else\textsuperscript{\ensuremath{n}}\fi}
\newunicodechar{ˣ}{\ifmmode^{x}\else\textsuperscript{\ensuremath{x}}\fi}
\newunicodechar{ᵇ}{\ifmmode^{b}\else\textsuperscript{\ensuremath{b}}\fi}
\newunicodechar{ʳ}{\ifmmode^{r}\else\textsuperscript{\ensuremath{r}}\fi}
\newunicodechar{ᵘ}{\ifmmode^{u}\else\textsuperscript{\ensuremath{u}}\fi}
\newunicodechar{ʸ}{\ifmmode^{y}\else\textsuperscript{\ensuremath{y}}\fi}
\newunicodechar{ᵃ}{\ifmmode^{a}\else\textsuperscript{\ensuremath{a}}\fi}
\newunicodechar{ᵉ}{\ifmmode^{e}\else\textsuperscript{\ensuremath{e}}\fi}
\newunicodechar{ᵏ}{\ifmmode^{k}\else\textsuperscript{\ensuremath{k}}\fi}
\newunicodechar{ᵖ}{\ifmmode^{p}\else\textsuperscript{\ensuremath{p}}\fi}
\newunicodechar{ᴺ}{\ifmmode^{N}\else\textsuperscript{\ensuremath{N}}\fi}
\newunicodechar{ᶜ}{\ifmmode^{c}\else\textsuperscript{\ensuremath{c}}\fi}
\newunicodechar{ʰ}{\ifmmode^{h}\else\textsuperscript{\ensuremath{h}}\fi}
\newunicodechar{ᵠ}{\ifmmode^{q}\else\textsuperscript{\ensuremath{q}}\fi}
\newunicodechar{ₙ}{\ifmmode_{n}\else\textsubscript{\ensuremath{n}}\fi}
\newunicodechar{ₐ}{\ifmmode_{a}\else\textsubscript{\ensuremath{a}}\fi}
\newunicodechar{ᵢ}{\ifmmode_{i}\else\textsubscript{\ensuremath{i}}\fi}
\newunicodechar{ᵣ}{\ifmmode_{r}\else\textsubscript{\ensuremath{r}}\fi}
\newunicodechar{ₖ}{\ifmmode_{k}\else\textsubscript{\ensuremath{k}}\fi}
\newunicodechar{ᵦ}{\ifmmode_{\beta}\else\textsubscript{\ensuremath{\beta}}\fi}
\newunicodechar{ₓ}{\ifmmode_{x}\else\textsubscript{\ensuremath{x}}\fi}
\newfontfamily\popdisplay{ArchivoBlack-Regular.ttf}[Path=../../../fonts/]
\newfontfamily\poplogo{SpaceGrotesk-Bold.ttf}[Path=../../../fonts/]
\newfontfamily\popsemi{PlusJakartaSans-SemiBold.ttf}[Path=../../../fonts/]
\newfontfamily\popxbold{PlusJakartaSans-ExtraBold.ttf}[Path=../../../fonts/]
\newfontfamily\popxboldls{PlusJakartaSans-ExtraBold.ttf}[Path=../../../fonts/,LetterSpace=3.0]

\setlist[enumerate]{leftmargin=1.7em,itemsep=5pt,topsep=4pt}
\setlist[itemize]{leftmargin=1.7em,itemsep=4pt,topsep=4pt,label={\color{poporange}\textbullet}}
\setlength{\parindent}{0pt}
\setlength{\parskip}{5pt}
\renewcommand{\arraystretch}{1.25}

% pgfplots : style Pop par défaut
\pgfplotsset{
  every axis/.append style={axis line style={popink,line width=1pt},tick style={popink},
    grid style={popline,line width=0.5pt},label style={font=\small},tick label style={font=\footnotesize}},
  popcurve/.style={poporange,line width=1.8pt,no markers,smooth},
}
% Même style disponible dans un \draw TikZ simple (hors pgfplots)
\tikzset{popcurve/.style={poporange,line width=1.8pt}}
\tikzset{popgrid/.style={popline,very thin}}
\colorlet{popcurve}{poporange} % le nom du style est parfois utilisé comme couleur

% ============================================================
% Pill / squiggle
% ============================================================
\newcommand{\poppill}[3]{%
  \tikz[baseline=-0.55ex]\node[draw=popink,line width=1.2pt,rounded corners=2.6pt,%
    fill=#2,inner xsep=6.5pt,inner ysep=3pt]{%
    \popxboldls\fontsize{7}{7}\selectfont\color{#3}\MakeUppercase{#1}};%
}
\newcommand{\popsquiggle}[1]{%
  \tikz[baseline=-0.4ex]{
    \draw[#1,line width=2.6pt,line cap=round]
      plot [smooth] coordinates {(0,0.09) (0.34,0.01) (0.68,0.21) (1.02,0.01) (1.36,0.10)};
  }%
}
% Mot-clé surligné (marqueur orange sous le texte)
\newcommand{\cle}[1]{{\popxbold\color{popdark}#1}}

% ============================================================
% En-tête / pied
% ============================================================
\pagestyle{fancy}
\fancyhf{}
\renewcommand{\headrulewidth}{0pt}
\renewcommand{\footrulewidth}{0pt}
\lhead{\raisebox{-0.15ex}{\poplogo\fontsize{13}{13}\selectfont\color{popink}OnBuch\color{poporange}+}}
\rhead{\poppill{\DOCMATIERE\ \textbullet\ \DOCNIVEAU\ \textbullet\ Leçon \DOCLECON}{white}{popink}}
\lfoot{\popxbold\fontsize{7.6}{7.6}\selectfont\color{popmuted}\MakeUppercase{Cours signé OnBuch+}\ \textbullet\ {\popsemi\DOCTITRE}}
\rfoot{\tikz[baseline=-0.6ex]\node[rounded corners=8.5pt,fill=popink,minimum height=17pt,minimum width=17pt,inner xsep=5pt,inner sep=0]{\popxbold\fontsize{8}{8}\selectfont\color{popcream}\thepage\,/\,\pageref*{LastPage}};}

% ============================================================
% Titres
% ============================================================
\newcommand{\chaptitre}[2]{%
  \begin{tcolorbox}[enhanced,colback=poporange,colframe=popink,boxrule=2.6pt,arc=11pt,
      drop shadow=popink,left=15pt,right=15pt,top=12pt,bottom=11pt,before skip=3pt,after skip=18pt]
    \poppill{#2}{popink}{popcream}\par\vspace{9pt}
    {\raggedright\hyphenpenalty=10000\exhyphenpenalty=10000\popdisplay\fontsize{22}{24}\selectfont\color{popcream}\MakeUppercase{#1}\par}
  \end{tcolorbox}
}
% Section numérotée : pastille orange avec le numéro + titre Archivo
\newcounter{popsec}
\newcommand{\coursec}[1]{%
  \stepcounter{popsec}\setcounter{popsub}{0}%
  \par\vspace{16pt}\noindent
  \begin{minipage}{\linewidth}
  \tikz[baseline=-0.7ex]\node[draw=popink,line width=1.6pt,fill=poporange,rounded corners=4pt,minimum size=22pt,inner sep=0]{\popdisplay\fontsize{12}{12}\selectfont\color{popcream}\thepopsec};%
  \hspace{8pt}{\popdisplay\fontsize{15}{17}\selectfont\color{popink}\MakeUppercase{#1}}\par
  \vspace{4pt}\noindent{\color{poporange}\rule{36pt}{3.6pt}}
  \end{minipage}\par\nopagebreak\vspace{8pt}
}
\newcounter{popsub}
\newcommand{\courssub}[1]{%
  \stepcounter{popsub}%
  \par\vspace{9pt}\noindent{\popxbold\fontsize{11.5}{13}\selectfont\color{popink}{\color{poporange}\thepopsec.\thepopsub}\hspace{6pt}#1}\par\nopagebreak\vspace{4pt}
}

% ============================================================
% Boîtes pédagogiques — même recette : fond blanc, bordure épaisse colorée,
% ombre dure encre, badge plein. Titre optionnel [#1].
% ============================================================
\tcbset{popbase/.style={breakable,enhanced,colback=white,boxrule=2.3pt,arc=9pt,
  shadow={3pt}{-3pt}{0pt}{popink},left=14pt,right=14pt,top=11pt,bottom=11pt,before skip=11pt,after skip=15pt,
  fontupper=\popsemi\fontsize{10.6}{14.5}\selectfont\color{popink},
  fontlower=\popsemi\fontsize{10.6}{14.5}\selectfont\color{popink}}}
% Coche dessinée (le glyphe ✓ n'existe pas dans les polices du document)
\newcommand{\popcheck}[1]{\tikz[baseline=-0.5ex]\draw[#1,line width=1.6pt,line cap=round,line join=round](-0.13,0.0)--(-0.04,-0.09)--(0.13,0.1);}
\providecommand{\coche}{\popcheck{popgreen}}
\providecommand{\resultat}[1]{\par\smallskip\noindent\fcolorbox{popgreen}{popgreenL}{\parbox{\dimexpr\linewidth-2\fboxsep-2\fboxrule\relax}{\textbf{Résultat.} #1}}\par\smallskip}
\newcommand{\popboxhead}[3]{\poppill{#1}{#2}{white}\ifx\relax#3\relax\else\hspace{6pt}{\popxbold\fontsize{10.6}{12}\selectfont\color{popink}#3}\fi\par\vspace{7pt}}

\newtcolorbox{aretenir}[1][]{popbase,colframe=popgreen,before upper={\popboxhead{À retenir}{popgreen}{#1}}}
\newtcolorbox{exemplebox}[1][]{popbase,colframe=poporange,before upper={\popboxhead{Exemple}{poporange}{#1}}}
\newtcolorbox{definition}[1][]{popbase,colframe=popblue,before upper={\popboxhead{Définition}{popblue}{#1}}}
\newtcolorbox{propriete}[1][]{popbase,colframe=poppurple,before upper={\popboxhead{Propriété}{poppurple}{#1}}}
\newtcolorbox{methode}[1][]{popbase,colframe=popgold,before upper={\popboxhead{Méthode}{popgold}{#1}}}
\newtcolorbox{attention}[1][]{popbase,colframe=poppink,before upper={\popboxhead{Attention}{poppink}{#1}}}
\newtcolorbox{experience}[1][]{popbase,colframe=popblue,colback=popblueL,before upper={\popboxhead{Expérience}{popblue}{#1}}}
% « Le savais-tu ? » : carte violette pleine (contexte, culture, Cameroun)
\newtcolorbox{savaistu}[1][]{popbase,colback=poppurple,colframe=popink,
  fontupper=\popsemi\fontsize{10.4}{14.2}\selectfont\color{white},
  before upper={\poppill{Le savais-tu\,?}{white}{poppurple}\ifx\relax#1\relax\else\hspace{6pt}{\popxbold\color{white}#1}\fi\par\vspace{7pt}}}
% Exercice résolu : énoncé en haut, \tcblower puis la solution sur fond crème
\newcounter{popexo}\newcounter{popexr}
\newtcolorbox{exoresolu}[1][]{popbase,colframe=poporange,colbacklower=popcream,
  segmentation style={popink,line width=1.2pt,dashed},
  before upper={\stepcounter{popexr}\popboxhead{Exercice résolu \thepopexr}{poporange}{#1}},
  before lower={\poppill{Solution}{popink}{popcream}\par\vspace{6pt}}}
% Exercice (non résolu en place) — la correction est dans la partie Corrigés
\newtcolorbox{exercice}[1][]{popbase,colframe=popink,
  before upper={\stepcounter{popexo}\popboxhead{Exercice \thepopexo}{popink}{#1}}}
% Corrigé
\newtcolorbox{corrige}[1][]{popbase,colframe=popgreen,colback=popgreenL,
  before upper={\popboxhead{Corrigé}{popgreen}{#1}}}
% Objectifs / prérequis / bilan
\newtcolorbox{objectifs}{popbase,colframe=popink,colback=poporangeL,
  before upper={\popboxhead{Objectifs de la leçon}{popink}{}}}
\newtcolorbox{prerequis}{popbase,colframe=popink,colback=popblueL,
  before upper={\popboxhead{Prérequis}{popblue}{}}}
\newtcolorbox{bilan}{popbase,colframe=popink,colback=white,boxrule=2.8pt,
  before upper={\popboxhead{Fiche bilan}{poporange}{L'essentiel à connaître pour l'examen}}}
% Figure encadrée avec légende
\newtcolorbox{popfigure}[1][]{enhanced,breakable=false,colback=white,colframe=popline,boxrule=1.4pt,arc=8pt,
  left=10pt,right=10pt,top=10pt,bottom=8pt,before skip=10pt,after skip=12pt,halign=center,
  lower separated=false,halign lower=center,
  fontlower=\popsemi\fontsize{9}{11.5}\selectfont\color{popmuted},
  #1}
\newcommand{\legende}[1]{\par\vspace{6pt}{\popsemi\fontsize{9}{11.5}\selectfont\color{popmuted}#1}}

% ============================================================
% Signature OnBuch+
% ============================================================
\newcommand{\poplogoinline}[1]{{\poplogo\fontsize{#1}{#1}\selectfont\color{popink}OnBuch\color{poporange}+}}
\newcommand{\signatureonbuch}{%
  \begin{tcolorbox}[enhanced,colback=popink,colframe=popink,boxrule=2.2pt,arc=10pt,
      drop shadow=poporange,left=14pt,right=14pt,top=10pt,bottom=10pt,before skip=0pt,after skip=0pt]
    \tikz[baseline=-0.7ex]\node[circle,fill=popgreen,draw=popcream,line width=1.3pt,minimum size=22pt,inner sep=0]{\popcheck{white}};%
    \hspace{9pt}\begin{minipage}[c]{0.62\linewidth}
      {\popxbold\fontsize{12}{13}\selectfont\color{popcream}Signé\ \textbullet\ {\poplogo OnBuch\color{poporange}+}}\par\vspace{2pt}
      {\raggedright\popsemi\fontsize{8}{10.5}\selectfont\color{popline}\MakeUppercase{Équipe pédagogique OnBuch+}\\\MakeUppercase{Conforme au programme officiel MINESEC}\par}
    \end{minipage}\hfill
    {\poplogo\fontsize{22}{22}\selectfont\color{popcream}OnBuch\color{poporange}+}
  \end{tcolorbox}%
}

% ============================================================
% Page de garde
% #1 titre · #2 matière · #3 niveau · #4 module · #5 n° leçon · #6 durée ·
% #7 description / objectifs (texte ou itemize)
% ============================================================
\newcommand{\pagegarde}[7]{%
  \thispagestyle{empty}
  \newgeometry{a4paper,margin=1.7cm,top=1.6cm,bottom=1.4cm}%
  \begin{tikzpicture}[remember picture,overlay]
    \fill[poporange!18] ([xshift=-3.2cm,yshift=-3.6cm]current page.north east) circle (4.6cm);
    \fill[poppurple!14] ([xshift=2.4cm,yshift=7.6cm]current page.south west) circle (3.8cm);
  \end{tikzpicture}%
  {\poplogo\fontsize{30}{30}\selectfont\color{popink}OnBuch\color{poporange}+}\hfill
  \raisebox{0.6ex}{\poppill{Cours signé \textbullet\ Premium}{popink}{popcream}}\par
  \vspace{1pt}\popsquiggle{poppurple}\par
  \vspace{8pt}
  \begin{tcolorbox}[enhanced,colback=poporange,colframe=popink,boxrule=2.8pt,arc=13pt,
      drop shadow=popink,left=18pt,right=18pt,top=12pt,bottom=13pt,after skip=10pt]
    \poppill{#2 \textbullet\ #3}{popink}{popcream}\hspace{5pt}\poppill{#4}{white}{popink}\par\vspace{8pt}
    {\popdisplay\fontsize{14}{14}\selectfont\color{popink}LEÇON #5}\par\vspace{4pt}
    {\raggedright\hyphenpenalty=10000\exhyphenpenalty=10000\popdisplay\fontsize{25}{27}\selectfont\color{popcream}\MakeUppercase{#1}\par}
    \vspace{8pt}
    {\popxbold\fontsize{9.5}{9.5}\selectfont\color{popink}\MakeUppercase{Durée conseillée : #6}}
  \end{tcolorbox}
  \begin{tcolorbox}[enhanced,colback=white,colframe=popink,boxrule=2.2pt,arc=10pt,
      drop shadow=popink,left=16pt,right=16pt,top=10pt,bottom=10pt,after skip=8pt]
    \poppill{Dans cette leçon}{poppurple}{white}\par\vspace{5pt}
    \popsemi\fontsize{9.3}{12.6}\selectfont\color{popink}#7
  \end{tcolorbox}
  \noindent\begin{minipage}[t]{0.487\linewidth}
    \begin{tcolorbox}[enhanced,colback=popgreen,colframe=popink,boxrule=2.2pt,arc=10pt,
        drop shadow=popink,left=11pt,right=11pt,top=10pt,bottom=10pt,height=3.1cm,valign=center]
      \tcbox[colback=white,colframe=popink,boxrule=1.3pt,arc=5pt,boxsep=0pt,nobeforeafter,
        left=3pt,right=3pt,top=3pt,bottom=3pt,valign=center]{\qrcode[height=1.9cm]{https://play.google.com/store/apps/details?id=com.luvvix.onbuch&hl=fr}}%
      \hspace{8pt}\begin{minipage}[c]{0.5\linewidth}
        {\popdisplay\fontsize{11}{12.5}\selectfont\color{white}\MakeUppercase{Télécharge OnBuch}}\par\vspace{4pt}
        {\raggedright\popsemi\fontsize{8.4}{11}\selectfont\color{white}Gratuit \textbullet\ Google Play\\Tuteur IA, annales, concours, résultats\par}
      \end{minipage}
    \end{tcolorbox}
  \end{minipage}\hfill
  \begin{minipage}[t]{0.487\linewidth}
    \begin{tcolorbox}[enhanced,colback=poppurple,colframe=popink,boxrule=2.2pt,arc=10pt,
        drop shadow=popink,left=11pt,right=11pt,top=10pt,bottom=10pt,height=3.1cm,valign=center]
      \tikz[baseline=-0.7ex]\node[circle,fill=white,draw=popink,line width=1.3pt,minimum size=1.6cm,inner sep=0]{\popdisplay\fontsize{24}{24}\selectfont\color{poppurple}+};%
      \hspace{8pt}\begin{minipage}[c]{0.6\linewidth}
        {\popdisplay\fontsize{11}{12.5}\selectfont\color{white}\MakeUppercase{Passe à OnBuch+}}\par\vspace{4pt}
        {\raggedright\popsemi\fontsize{8.4}{11}\selectfont\color{white}Tous les cours signés, Tuteur IA illimité, fiches PDF — onglet OnBuch+ de l'app\par}
      \end{minipage}
    \end{tcolorbox}
  \end{minipage}
  \vspace{8pt}
  \popcontenu
  \vfill
  \signatureonbuch
  \restoregeometry
  \newpage
}

% Rangée « Ce cours contient » (page de garde)
\newcommand{\poptile}[3]{% #1 couleur #2 gros texte #3 légende
  \begin{tcolorbox}[enhanced,colback=#1,colframe=popink,boxrule=2pt,arc=9pt,shadow={2.5pt}{-2.5pt}{0pt}{popink},
      left=6pt,right=6pt,top=9pt,bottom=9pt,halign=center,valign=center,height=1.75cm,width=\linewidth,nobeforeafter]
    {\popdisplay\fontsize{12.5}{13}\selectfont\color{white}\MakeUppercase{#2}}\par\vspace{3pt}
    {\centering\popsemi\fontsize{7.8}{9.5}\selectfont\color{white}#3\par}
  \end{tcolorbox}}
\newcommand{\popcontenu}{%
  \par\noindent{\popxboldls\fontsize{7.5}{7.5}\selectfont\color{popmuted}CE COURS SIGNÉ CONTIENT}\par\vspace{7pt}
  \noindent\begin{minipage}[t]{0.235\linewidth}\poptile{popblue}{Cours}{complet, illustré, pas à pas}\end{minipage}\hfill
  \begin{minipage}[t]{0.235\linewidth}\poptile{poporange}{Méthodes}{et exercices résolus}\end{minipage}\hfill
  \begin{minipage}[t]{0.235\linewidth}\poptile{poppink}{Exercices}{type examen + corrigés}\end{minipage}\hfill
  \begin{minipage}[t]{0.235\linewidth}\poptile{popgreen}{Fiche bilan}{l'essentiel à retenir}\end{minipage}\par}

% Sommaire
\newtcolorbox{sommaire}{enhanced,colback=white,colframe=popink,boxrule=2.2pt,arc=10pt,
  drop shadow=popink,left=18pt,right=18pt,top=13pt,bottom=13pt,before skip=6pt,after skip=20pt,
  before upper={\poppill{Sommaire}{poppurple}{white}\par\vspace{9pt}\popsemi\fontsize{10.6}{14.5}\selectfont\color{popink}}}

% Signature de fin de cours — poussée tout en bas de la dernière page
\newcommand{\findecours}{%
  \par\vfill\nopagebreak
  \begin{center}\popsquiggle{poporange}\end{center}\vspace{4pt}
  \signatureonbuch
}
\AtBeginDocument{\def\llbracket{\mathopen{[\mkern-3.5mu[}}\def\rrbracket{\mathclose{]\mkern-3.5mu]}}} % intervalles d'entiers (glyphe ⟦ absent de la police)
% Glyphes absents de Fira Math : redéfinis à partir de glyphes disponibles
\AtBeginDocument{%
  \def\setminus{\mathbin{\backslash}}\let\smallsetminus\setminus
  \def\vdots{\mathord{\vbox{\baselineskip3.2pt\lineskiplimit0pt\kern4pt\hbox{.}\hbox{.}\hbox{.}}}}%
  \def\ddots{\mathinner{\mkern1mu\raise6.5pt\hbox{.}\mkern2mu\raise3.6pt\hbox{.}\mkern2mu\raise0.7pt\hbox{.}\mkern1mu}}%
  \def\triangle{\mathord{\scalebox{0.9}{$\symup{\Delta}$}}}%
  \def\nabla{\mathord{\raisebox{\depth}{\scalebox{1}[-1]{$\symup{\Delta}$}}}}%
  \def\checkmark{\popcheck{popdark}}%
  \def\star{\mathbin{\ast}}\def\bigstar{\mathord{\ast}}%
  \def\diamond{\mathbin{\scalebox{0.6}{$\Diamond$}}}%
  \def\bigcup{\mathop{\mathchoice{\raisebox{-0.3em}{\scalebox{1.7}{$\cup$}}}{\scalebox{1.25}{$\cup$}}{\cup}{\cup}}\displaylimits}%
  \def\bigcap{\mathop{\mathchoice{\raisebox{-0.3em}{\scalebox{1.7}{$\cap$}}}{\scalebox{1.25}{$\cap$}}{\cap}{\cap}}\displaylimits}%
  \def\lesseqqgtr{\mathrel{\lessgtr}}\def\bigoplus{\mathop{\oplus}\displaylimits}%
}
\providecommand{\ssi}{si et seulement si} % abréviation fréquente
\AtBeginDocument{\providecommand{\wideparen}{\overparen}} % arc de cercle

%% FIGFIX-BEGIN v3 — correctifs globaux des figures (docs/onbuch-generation/tools/figfix)
\usepackage{adjustbox}\usepackage{etoolbox}
% toute figure TikZ/pgfplots plus large que la ligne est réduite à la largeur disponible
\BeforeBeginEnvironment{tikzpicture}{\begin{adjustbox}{max width=\linewidth}}
\AfterEndEnvironment{tikzpicture}{\end{adjustbox}}
% légendes pgfplots lisibles par-dessus les courbes
\pgfplotsset{every axis/.append style={legend style={fill=white,fill opacity=0.92,text opacity=1,draw=black!15,inner sep=3pt}}}
% étiquettes d'arêtes (syntaxe edge["..."]) avec fond blanc
\tikzset{every edge quotes/.append style={fill=white,inner sep=1.2pt}}
% bulles de cartes mentales : texte à la ligne dans la bulle, bulle un peu plus grande
\tikzset{every concept/.append style={text width=2.0cm,align=center,minimum size=2.3cm,inner sep=2pt}}
%% FIGFIX-END
\begin{document}

\pagegarde{Le microscope}{Physique}{1ère E}{Module 3 : Optique}{N09}{5 h}{Cette leçon détaille le fonctionnement du microscope optique composé, de sa description géométrique au calcul de son grossissement et de sa puissance commerciale. Elle guide l'élève vers la maîtrise de la mise au point et l'analyse des aberrations, compétences clés pour le Probatoire.
\vspace{4pt}
\begin{multicols}{2}\small
\begin{itemize}
\item À la fin de cette leçon, je sais décrire le microscope (objectif, oculaire, tube, platine, vis de mise au point) et expliquer le rôle de chaque élément.
\item À la fin de cette leçon, je sais tracer le schéma de principe (faisceaux parallèles, foyers, plans principaux) et justifier la formation de l'image finale virtuelle, inversée et grossie.
\item À la fin de cette leçon, je sais définir le grossissement commercial G = G\_obj × G\_oc et la puissance commerciale P = P\_obj + P\_oc (en dioptries), et les calculer à partir des focales.
\item À la fin de cette leçon, je sais expliquer la condition de fonctionnement normal (image finale à l'infini) et la condition de confort (image au point proche, 25 cm).
\item À la fin de cette leçon, je sais réaliser la mise au point sur un schéma (déplacer l'objectif pour que l'objet soit au foyer image de l'objectif) et décrire la manipulation réelle (macrométrique puis micrométrique).
\item À la fin de cette leçon, je sais calculer la longueur du tube optique L = Δ + f'\_obj + f'\_oc et en déduire l'encombrement de l'instrument.
\end{itemize}\end{multicols}}

\begin{sommaire}
\begin{multicols}{2}
\begin{enumerate}
\item 1. Description et rôle du microscope
\item 2. Schéma de principe et construction géométrique
\item 3. Grossissement commercial et Puissance commerciale
\item 4. Mise au point : théorie et pratique
\item 5. Longueur du tube, encombrement et limites (Complément Probatoire)
\item 6. Résolution de problèmes types (Probatoire)
\item Activité d'intégration
\item Méthodes et exercices résolus
\item Exercices
\item Corrigés des exercices
\item Fiche bilan
\end{enumerate}
\end{multicols}
\end{sommaire}

\begin{objectifs}
\begin{itemize}
\item À la fin de cette leçon, je sais décrire le microscope (objectif, oculaire, tube, platine, vis de mise au point) et expliquer le rôle de chaque élément.
\item À la fin de cette leçon, je sais tracer le schéma de principe (faisceaux parallèles, foyers, plans principaux) et justifier la formation de l'image finale virtuelle, inversée et grossie.
\item À la fin de cette leçon, je sais définir le grossissement commercial G = G\_obj × G\_oc et la puissance commerciale P = P\_obj + P\_oc (en dioptries), et les calculer à partir des focales.
\item À la fin de cette leçon, je sais expliquer la condition de fonctionnement normal (image finale à l'infini) et la condition de confort (image au point proche, 25 cm).
\item À la fin de cette leçon, je sais réaliser la mise au point sur un schéma (déplacer l'objectif pour que l'objet soit au foyer image de l'objectif) et décrire la manipulation réelle (macrométrique puis micrométrique).
\item À la fin de cette leçon, je sais calculer la longueur du tube optique L = Δ + f'\_obj + f'\_oc et en déduire l'encombrement de l'instrument.
\item À la fin de cette leçon, je sais identifier les limites du microscope optique (pouvoir séparateur, diffraction, aberrations) — complément Probatoire.
\item À la fin de cette leçon, je sais résoudre un problème complet : dimensionner un microscope pour un grossissement donné ou vérifier la conformité d'un instrument du commerce.
\end{itemize}
\end{objectifs}

\begin{prerequis}
\begin{itemize}
\item Lentille mince : conjugaison, grossissement, formule de Newton et de Descartes (1ère C, Module 3, Leçons 1-4).
\item Modèle de la lentille épaisse : plans principaux, foyers principaux, distance inter-foculaire Δ (Début Module 3).
\item Loupe : grossissement commercial G = 25 cm / f' + 1 (confort) ou 25 cm / f' (normal), puissance P = 1/f' (m⁻¹).
\item Notion d'image réelle et virtuelle, inversée/redressée, grandissement algébrique γ.
\item Unités : mètre, centimètre, dioptrie (δ), conversion mm ↔ m.
\end{itemize}
\end{prerequis}

\newpage

\coursec{Description et rôle du microscope}

Au laboratoire de biologie du Lycée de Nkolbisson, un élève observe une lame de sang à la recherche de parasites du paludisme : l'image est floue, puis soudain nette et grossie 400 fois. Comment cet assemblage de deux lentilles simples permet-il de voir l'invisible ? Pourquoi doit-on tourner la vis macrométrique puis la micrométrique ? Et si l'on change l'objectif par un autre de puissance différente, que se passe-t-il pour la taille finale de l'image ? C'est à ces questions que cette section répond en décrivant l'instrument, ses éléments et son principe de fonctionnement.

\courssub{Définition et rôle général}

\begin{definition}[Microscope composé]
Instrument d'optique à deux lentilles convergentes (objectif et oculaire) séparées par une distance $\Delta$ (intervalle inter-foculaire) servant à observer de très petits objets placés près du foyer objet de l'objectif.
\end{definition}

Le microscope composé a pour \textbf{rôle} de produire une \textbf{image finale virtuelle}, \textbf{inversée} par rapport à l'objet initial, \textbf{très grossie}, située soit au \textbf{point proche} de l'œil ($D = \SI{25}{cm}$) pour un œil accommodé (fonctionnement au confort), soit à \textbf{l'infini} pour un œil au repos (fonctionnement normal). Contrairement à la loupe qui grossit un objet déjà visible, le microscope révèle des détails invisibles à l'œil nu (cellules, bactéries, globules rouges du paludisme).

\courssub{Éléments mécaniques et optiques}

\begin{popfigure}
\begin{tikzpicture}[scale=0.9, >=Stealth]
% --- Dimensions ---
\def\Ltube{10} % longueur tube optique
\def\fobj{1.2} % focale objectif
\def\foc{2.5}  % focale oculaire
\def\Delta{6}  % intervalle inter-foculaire
\def\hobj{0.8} % hauteur objet
\def\himg{2.5} % hauteur image intermédiaire

% --- Plans principaux et foyers ---
% Objectif
\coordinate (Hobj) at (0,0);
\coordinate (Fobj) at (-\fobj,0);
\coordinate (Fpobj) at (\fobj,0);
% Oculaire
\coordinate (Hoc) at (\Ltube,0);
\coordinate (Foc) at (\Ltube-\foc,0);
\coordinate (Fpoc) at (\Ltube+\foc,0);

% --- Tube mécanique ---
\draw[popdark, thick] (-0.5,-2.5) rectangle (\Ltube+0.5,2.5);
\draw[popmuted, dashed] (0,-2.5) -- (0,2.5) node[midway, left] {Objectif};
\draw[popmuted, dashed] (\Ltube,-2.5) -- (\Ltube,2.5) node[midway, right] {Oculaire};

% --- Lentilles (représentation symbolique) ---
\draw[popblue, very thick] (0,-1.5) -- (0,1.5) node[above] {$L_{obj}$};
\draw[popblue, very thick] (\Ltube,-1.2) -- (\Ltube,1.2) node[above] {$L_{oc}$};

% --- Foyers et plans principaux ---
\draw[poporange, thick] (Fobj) -- ++(0,-0.3) node[below] {$F_{obj}$};
\draw[poporange, thick] (Fpobj) -- ++(0,-0.3) node[below] {$F'_{obj}$};
\draw[poporange, thick] (Foc) -- ++(0,0.3) node[above] {$F_{oc}$};
\draw[poporange, thick] (Fpoc) -- ++(0,0.3) node[above] {$F'_{oc}$};
\draw[poppurple, thick] (Hobj) -- ++(0,0.3) node[above] {$H_{obj}$};
\draw[poppurple, thick] (Hoc) -- ++(0,0.3) node[above] {$H_{oc}$};

% --- Intervalle inter-foculaire Δ ---
\draw[<->, popgreen, thick] (Fpobj) ++(0,1.2) -- (Foc) ++(0,1.2) node[midway, above] {$\Delta$};

% --- Objet AB ---
\coordinate (A) at (-\fobj-0.15, \hobj);
\coordinate (B) at (-\fobj-0.15, 0);
\draw[popdark, very thick] (A) -- (B) node[midway, left] {$A$} node[below] {$B$};
\draw[<->, popmuted] (Fobj) -- (A) node[midway, below] {$p \approx f'_{obj}+\varepsilon$};

% --- Image intermédiaire A'B' (réelle, inversée, grossie) ---
\coordinate (Ap) at (\Ltube-\foc, -\himg);
\coordinate (Bp) at (\Ltube-\foc, 0);
\draw[popred, very thick] (Ap) -- (Bp) node[midway, right] {$A'$} node[above] {$B'$};
\draw[popmuted, dashed] (Foc) -- (Ap);

% --- Rayons principaux pour l'objectif ---
% Rayon parallèle -> passe par F'obj
\draw[popblue, ->] (A) -- (0,\hobj) -- (Fpobj);
% Rayon centre optique -> non dévié
\draw[popblue, ->] (A) -- (Hobj) -- (Ap);
% Rayon passant par Fobj -> sort parallèle
\draw[popblue, ->] (A) -- (Fobj) -- (0,\hobj) -- ++(\Ltube,0);

% --- Rayons principaux pour l'oculaire (image finale à l'infini) ---
% Depuis A' (au foyer Foc) -> rayons sortent parallèles
\draw[popred, ->] (Ap) -- (Hoc) -- ++(3,0) node[above, pos=0.8] {$\parallel$};
\draw[popred, ->] (Ap) -- (\Ltube, -\himg) -- ++(3,0) node[above, pos=0.8] {$\parallel$};

% --- Œil ---
\draw[popdark, fill=popcream] (13.5,0) ellipse (0.4 and 0.25);
\draw[popdark] (13.5,0) circle (0.15);
\node[below] at (13.5,-0.4) {Œil au repos};

% --- Légendes ---
\node[align=center, popdark] at (5, -3.2) {Fonctionnement normal : image finale à l'infini};
\end{tikzpicture}
\legende{Schéma de principe du microscope composé en fonctionnement normal (image finale à l'infini). L'objet $AB$ est placé juste au-delà de $F_{obj}$. L'image intermédiaire $A'B'$ est réelle, inversée et grossie, formée au foyer objet $F_{oc}$ de l'oculaire. L'oculaire agit comme une loupe donnant une image finale virtuelle à l'infini.}
\end{popfigure}

\begin{center}
\begin{tabularx}{\linewidth}{|l|X|X|}\hline
\rowcolor{popblueL}
\textbf{Élément} & \textbf{Description} & \textbf{Rôle optique} \\ \hline
\textbf{Tube} & Longueur mécanique normalisée $\approx \SI{160}{mm}$ (standard DIN/ISO). Supporte l'objectif et l'oculaire. & Définit l'encombrement ; $\Delta = \overline{F'_{obj}F_{oc}}$ est la grandeur optique utile. \\ \hline
\textbf{Revolver porte-objectifs} & Tourelle rotative portant 3 à 4 objectifs de puissances croissantes (ex. $\times 4$, $\times 10$, $\times 40$, $\times 100$). & Permet de changer le grossissement sans déplacer la mise au point (parfocalité). \\ \hline
\textbf{Platine mobile} & Translation $X/Y$ micrométrique (vis à pas fin). & Déplace l'échantillon dans le plan objet pour centrer la zone d'intérêt. \\ \hline
\textbf{Vis macrométrique} & Approche rapide (pas $\sim \SI{1}{mm/tour}$). & Rapproche grossièrement l'objectif de la platine (ou inverse) pour la mise au point approximative. \\ \hline
\textbf{Vis micrométrique} & Précision $\sim \SI{0,002}{mm}$ ($\SI{2}{\mu m}$). & Affine la mise au point pour obtenir une image nette (profondeur de champ $\sim \mu m$). \\ \hline
\textbf{Condenseur + diaphragme iris} & Sous la platine, lentilles + diaphragme variable. & Réalise l'éclairage de Köhler : illumination uniforme, contrôle de l'ouverture numérique (résolution, contraste). \\ \hline
\end{tabularx}
\end{center}

\begin{definition}[Objectif]
Lentille convergente (souvent système lenticulaire corrigé) de \textbf{courte focale} $f'_{obj}$ (\SIrange{2}{20}{mm}), de \textbf{forte puissance} ($P_{obj} = 1/f'_{obj}$ jusqu'à $\SI{500}{\dioptre}$), à \textbf{grand diamètre numérique} (ouverture angulaire grande) pour collecter un maximum de lumière diffractée par l'objet.
\end{definition}

\begin{definition}[Oculaire]
Lentille convergente (souvent doublet ou triplet pour corriger les aberrations) de \textbf{focale plus longue} $f'_{oc}$ (\SIrange{10}{50}{mm}), de \textbf{puissance moindre} ($P_{oc} \approx \SIrange{20}{100}{\dioptre}$), définissant le \textbf{champ de vision} apparent (diamètre du diaphragme de champ).
\end{definition}

\courssub{Position de l'objet et formation des images}

L'objet $AB$ (ex. un globule rouge de $\sim \SI{7}{\mu m}$) est placé \textbf{juste au-delà du foyer objet de l'objectif} :
\[
p = \overline{H_{obj}A} \approx f'_{obj} + \varepsilon \qquad (\varepsilon \ll f'_{obj})
\]
L'objectif forme alors une \textbf{image intermédiaire $A'B'$} :
\begin{itemize}
    \item \textbf{Réelle} (formée par convergence effective des rayons),
    \item \textbf{Inversée} par rapport à $AB$,
    \item \textbf{Très grossie} ($|\gamma_{obj}| \gg 1$),
    \item Située dans le \textbf{plan focal objet de l'oculaire} ($F_{oc}$) pour le fonctionnement normal, ou \textbf{légèrement en avant} de $F_{oc}$ pour le fonctionnement au confort.
\end{itemize}

\begin{definition}[Fonctionnement normal]
L'image intermédiaire se forme exactement au foyer objet de l'oculaire ($F_{oc}$). L'image finale est à l'infini. L'œil est au repos (accommodation nulle). C'est le mode recommandé pour les observations prolongées (fatigue oculaire minimale).
\end{definition}

\begin{definition}[Fonctionnement au confort (ou vision distincte)]
L'image intermédiaire se forme légèrement en avant de $F_{oc}$. L'image finale se forme au point proche de l'œil ($D = \SI{25}{cm}$). L'œil accommode au maximum. Grossissement légèrement plus grand mais fatigant.
\end{definition}

L'oculaire agit alors comme une \textbf{loupe} : il prend l'image intermédiaire comme objet et produit une \textbf{image finale $A''B''$ virtuelle}, \textbf{inversée par rapport à l'objet initial} (donc de même orientation que $A'B'$), et \textbf{encore grossie}.

\begin{attention}[Erreur fréquente : longueur mécanique vs distance inter-foculaire]
Ne confondez \textbf{jamais} la \textbf{longueur mécanique du tube} (souvent $\SI{160}{mm}$ normalisée, distance entre les épaulements de l'objectif et de l'oculaire) et la \textbf{distance inter-foculaire} $\Delta = \overline{F'_{obj}F_{oc}}$. C'est $\Delta$ qui entre dans les formules de grossissement et de puissance commerciale. Sur un microscope standard, $\Delta \approx \SI{130}{mm}$ à $\SI{140}{mm}$ selon les marques.
\end{attention}

\courssub{Éclairage et résolution (aperçu)}

L'éclairage se fait par \textbf{lumière transmise} (échantillon transparent sur lame) ou \textbf{réfléchie} (échantillon opaque, métallographie). Le \textbf{condenseur} (sous la platine) et son \textbf{diaphragme iris} réalisent l'éclairage de \textbf{Köhler} : la source (LED ou lampe halogène) est imagée au plan pupillaire de l'objectif, assurant une illumination uniforme et contrôlant l'\textbf{ouverture numérique} :
\[
\text{NA} = n \sin u'_{max}
\]
où $n$ est l'indice du milieu entre l'objet et l'objectif (air $n=1$, huile $n=1,51$) et $u'_{max}$ le demi-angle d'ouverture. La \textbf{puissance de résolution} (distance minimale entre deux points discernables) est donnée par le critère d'Abbe-Rayleigh :
\[
\delta_{min} \approx \frac{0,61 \lambda}{\text{NA}}
\]
C'est pourquoi les objectifs à immersion ($\times 100$, NA $\approx 1,3-1,4$) atteignent $\delta_{min} \approx \SI{200}{nm}$ (lumière visible), révélant les bactéries du paludisme.

\begin{methode}[Repérer les plans principaux sur un schéma]
\begin{enumerate}
    \item Pour chaque lentille (objectif, oculaire), tracer les plans principaux $H$ et $H'$ (confondus au centre pour lentilles minces).
    \item Placer les foyers $F, F'$ à distance $f, f'$ de $H, H'$ ($f' > 0$ pour lentille convergente).
    \item Mesurer $\Delta = \overline{F'_{obj}F_{oc}}$ (distance entre le foyer image de l'objectif et le foyer objet de l'oculaire).
    \item Positionner l'objet $AB$ juste au-delà de $F_{obj}$ ($p \gtrsim f'_{obj}$).
    \item Construire l'image intermédiaire $A'B'$ par l'objectif (réelle, inversée, grossie).
    \item Vérifier que $A'B'$ se trouve en $F_{oc}$ (normal) ou légèrement avant (confort).
    \item Construire l'image finale $A''B''$ par l'oculaire (virtuelle, à l'infini ou en $D$).
\end{enumerate}
\end{methode}

\begin{exoresolu}[Analyse d'un microscope standard]
Un microscope de laboratoire (type école) possède un tube mécanique de $\SI{160}{mm}$ (standard DIN). L'objectif est gravé $\times 40 / 0,65$ (grossissement commercial $\times 40$, ouverture numérique $0,65$). L'oculaire est gravé $\times 10$ (champ $18$ mm). On prend $f'_{oc} = \SI{25}{mm}$.
\begin{enumerate}
    \item Calculer la focale de l'objectif $f'_{obj}$.
    \item Déterminer l'intervalle inter-foculaire $\Delta$.
    \item L'objet est un globule rouge de diamètre $d = \SI{7,5}{\mu m}$. Quelle est la taille de l'image intermédiaire ? De l'image finale (fonctionnement normal) ?
\end{enumerate}
\textbf{Solution détaillée}
\begin{enumerate}
    \item Le grossissement commercial de l'objectif est défini par rapport à la longueur mécanique du tube $L = \SI{160}{mm}$ (convention constructeur) : $G_{obj} = \frac{L}{f'_{obj}(\text{mm})}$. Avec $G_{obj} = 40$, on obtient $f'_{obj} = \frac{\SI{160}{mm}}{40} = \SI{4,0}{mm}$.
    \item L'intervalle inter-foculaire est $\Delta = \overline{F'_{obj}F_{oc}} = L - f'_{obj} - f'_{oc} = \SI{160}{mm} - \SI{4,0}{mm} - \SI{25}{mm} = \SI{131}{mm}$.
    \item Grossissement transverse de l'objectif (relation de Newton) : $\gamma_{obj} \approx -\frac{\Delta}{f'_{obj}} = -\frac{131}{4,0} \approx -32,8$. \textit{(La valeur commerciale $40$ est une valeur normalisée pour $L=160$ mm ; le grossissement réel transverse dépend de $\Delta$ effectif).} Taille image intermédiaire : $d' = |\gamma_{obj}| \times d \approx 32,8 \times \SI{7,5}{\mu m} \approx \SI{246}{\mu m} = \SI{0,246}{mm}$.
    \item Grossissement angulaire de l'oculaire (fonctionnement normal) : $G_{oc} = \frac{250}{f'_{oc}(\text{mm})} = \frac{250}{25} = 10$. Grossissement commercial total : $G = G_{obj} \times G_{oc} = 40 \times 10 = 400$. L'image finale virtuelle à l'infini sous-tend un angle $\alpha' = G \times \frac{d}{250\,\text{mm}} = 400 \times \frac{\SI{7,5}{\mu m}}{\SI{250}{mm}} = \SI{12}{mrad}$.
\end{enumerate}
\end{exoresolu}

\begin{savaistu}[Le microscope au Cameroun : de la recherche au diagnostic]
Au Centre Pasteur du Cameroun (Yaoundé) et dans les laboratoires des hôpitaux de district, le microscope optique reste l'outil \textbf{de référence} pour le diagnostic du paludisme (frottis sanguin épais et mince, coloration de Giemsa). Les objectifs $\times 100$ à immersion d'huile (NA $1,3$) permettent de distinguer les stades du \emph{Plasmodium falciparum} (anneaux, trophozoïtes, gamétocytes). L'électrification rurale (projets ENERGIE, SONATREL) a permis d'équiper les centres de santé en microscopes à LED (faible consommation, durée de vie $> \SI{20000}{h}$, température de couleur stable), remplaçant les lampes halogènes coûteuses. En enseignement, les lycées (Nkolbisson, Biyem-Assi, etc.) utilisent des microscopes « école » robustes (objectifs achromatiques $\times 4, \times 10, \times 40$) pour l'observation de cellules végétales (oignon, elodea) et animales (frottis sanguin, épithélium buccal), illustrant concrètement le programme de SVT et de Physique (optique géométrique).
\end{savaistu}

\begin{exemplebox}[Choix de l'objectif pour une observation]
Un élève doit observer une coupe de tige de maïs (épaisseur $\sim \SI{50}{\mu m}$) puis une lame de sang (globules $\sim \SI{7}{\mu m}$).
\begin{itemize}
    \item \textbf{Coupe de maïs :} objectif $\times 4$ ($f' \approx \SI{30}{mm}$, NA $\approx 0,10$). Grande distance de travail ($\sim \SI{10}{mm}$), champ large ($\sim \SI{4}{mm}$), faible grossissement suffisant pour voir les tissus.
    \item \textbf{Lame de sang :} objectif $\times 40$ ($f' \approx \SI{4}{mm}$, NA $\approx 0,65$). Distance de travail courte ($\sim \SI{0,5}{mm}$), fort grossissement, bonne résolution pour distinguer les parasites.
    \item \textbf{Objectif $\times 100$ (huile) :} indispensable pour la morphologie fine des bactéries (tuberculose, méningite) ou les stades du paludisme. Nécessite une goutte d'huile d'immersion ($n=1,51$) entre la lame et la lentille frontale.
\end{itemize}
\textbf{Règle pratique :} On commence toujours par le plus petit grossissement ($\times 4$) pour localiser la zone, puis on monte en puissance ($\times 10$, $\times 40$) sans toucher à la vis macrométrique (parfocalité), en n'utilisant que la vis micrométrique.
\end{exemplebox}

\begin{attention}[Pièges de rédaction au Probatoire]
\begin{itemize}
    \item Dire « l'image finale est droite » \textbf{FAUX}. Elle est \textbf{inversée} par rapport à l'objet (renversée haut/bas et gauche/droite). L'image intermédiaire est inversée, l'oculaire ne la renverse pas.
    \item Confondre « puissance de l'objectif » (en dioptries, $P = 1/f'$) et « grossissement commercial de l'objectif » (sans unité, $G_{obj} = L/f'_{obj}$).
    \item Oublier que $\Delta$ est la distance \textbf{inter-foculaire} ($F'_{obj}F_{oc}$), pas la distance entre lentilles ($H_{obj}H_{oc}$) ni la longueur mécanique du tube.
    \item Ne pas préciser le mode de fonctionnement (normal ou confort) lors du calcul du grossissement de l'oculaire.
\end{itemize}
\end{attention}

\coursec{Schéma de principe et construction géométrique}

\courssub{Construction rigoureuse de l'image intermédiaire (Objectif)}

L'objectif fonctionne comme une lentille convergente unique (modèle de Gauss) formant l'image réelle d'un objet réel. On utilise les relations de \textbf{Newton} (préférables car l'objet est près du foyer) :
\[
x \cdot x' = f'_{obj}{}^2 \quad \text{avec} \quad x = \overline{F_{obj}A},\; x' = \overline{F'_{obj}A'}
\]
L'objet étant placé à $p = f'_{obj} + \varepsilon$ ($\varepsilon > 0$ petit), on a $x = -\varepsilon$ (objet entre $F_{obj}$ et $H_{obj}$). Donc :
\[
x' = \frac{f'_{obj}{}^2}{-\varepsilon} \ll -f'_{obj}
\]
L'image $A'B'$ est réelle ($x' < 0$), inversée, située loin de $F'_{obj}$ du côté sortant. Le grossissement transverse de l'objectif est :
\[
\gamma_{obj} = \frac{A'B'}{AB} = -\frac{f'_{obj}}{x} = \frac{f'_{obj}}{\varepsilon} \quad \text{(valeur algébrique négative, image inversée)}
\]
En valeur absolue, $|\gamma_{obj}| = \frac{f'_{obj}}{\varepsilon} \gg 1$. Comme $A'B'$ doit se former au voisinage de $F_{oc}$, on a $\overline{F'_{obj}A'} \approx \Delta$, d'où l'approximation fondamentale :
\[
|\gamma_{obj}| \approx \frac{\Delta}{f'_{obj}}
\]

\courssub{Construction de l'image finale (Oculaire)}

L'oculaire fonctionne comme une loupe observant l'image intermédiaire $A'B'$.

\begin{popfigure}
\begin{tikzpicture}[scale=0.9, >=Stealth]
% --- Dimensions ---
\def\Ltube{10}
\def\fobj{1.2}
\def\foc{2.5}
\def\Delta{6}
\def\hobj{0.8}
\def\himg{2.5}
\def\eps{0.3} % epsilon pour confort

% --- Plans principaux et foyers ---
\coordinate (Hobj) at (0,0);
\coordinate (Fobj) at (-\fobj,0);
\coordinate (Fpobj) at (\fobj,0);
\coordinate (Hoc) at (\Ltube,0);
\coordinate (Foc) at (\Ltube-\foc,0);
\coordinate (Fpoc) at (\Ltube+\foc,0);

% --- Tube ---
\draw[popdark, thick] (-0.5,-2.5) rectangle (\Ltube+0.5,2.5);
\draw[popblue, very thick] (0,-1.5) -- (0,1.5) node[above] {$L_{obj}$};
\draw[popblue, very thick] (\Ltube,-1.2) -- (\Ltube,1.2) node[above] {$L_{oc}$};

% --- Foyers ---
\draw[poporange, thick] (Fobj) -- ++(0,-0.3) node[below] {$F_{obj}$};
\draw[poporange, thick] (Fpobj) -- ++(0,-0.3) node[below] {$F'_{obj}$};
\draw[poporange, thick] (Foc) -- ++(0,0.3) node[above] {$F_{oc}$};
\draw[poporange, thick] (Fpoc) -- ++(0,0.3) node[above] {$F'_{oc}$};
\draw[poppurple, thick] (Hobj) -- ++(0,0.3) node[above] {$H_{obj}$};
\draw[poppurple, thick] (Hoc) -- ++(0,0.3) node[above] {$H_{oc}$};

% --- Δ ---
\draw[<->, popgreen, thick] (Fpobj) ++(0,1.2) -- (Foc) ++(0,1.2) node[midway, above] {$\Delta$};

% --- Objet AB ---
\coordinate (A) at (-\fobj-\eps, \hobj);
\coordinate (B) at (-\fobj-\eps, 0);
\draw[popdark, very thick] (A) -- (B) node[midway, left] {$A$} node[below] {$B$};

% --- Image intermédiaire A'B' (légèrement AVANT Foc pour confort) ---
\coordinate (Ap) at (\Ltube-\foc-0.5, -\himg);
\coordinate (Bp) at (\Ltube-\foc-0.5, 0);
\draw[popred, very thick] (Ap) -- (Bp) node[midway, right] {$A'$} node[above] {$B'$};

% --- Rayons Objectif ---
\draw[popblue, ->] (A) -- (0,\hobj) -- (Fpobj);
\draw[popblue, ->] (A) -- (Hobj) -- (Ap);
\draw[popblue, ->] (A) -- (Fobj) -- (0,\hobj) -- ++(\Ltube,0);

% --- Rayons Oculaire (Image finale en D = -250 mm à gauche de Hoc) ---
% On trace les rayons émergents divergents comme s'ils venaient de A''B'' virtuel à gauche de l'oculaire
% Pour la construction, on trace les rayons incidents sur l'oculaire depuis A'
% Rayon passant par Hoc -> non dévié, prolongé en pointillés vers la gauche
\draw[popred, ->] (Ap) -- (Hoc) -- ++(-3,0);
\draw[popred, dashed] (Hoc) -- ++(-3,0);
% Rayon incident parallèle à l'axe -> passe par F'oc
\draw[popred, ->] (Ap) -- (\Ltube, -\himg) -- (Fpoc);
\draw[popred, dashed] (Fpoc) -- ++(-3,0);

% --- Image finale virtuelle A''B'' (intersection des prolongements) ---
% Calcul approximatif position A'' : p' = distance A'Hoc = f'oc + 0.5 = 3.0 cm. f'oc = 2.5 cm.
% 1/p' + 1/p'' = 1/f' -> 1/3 + 1/p'' = 1/2.5 -> 1/p'' = 0.4 - 0.333 = 0.0667 -> p'' = 15 cm = 150 mm.
% Donc A'' est à 15 cm à gauche de Hoc. (D=250mm standard, ici p'=30mm donne p''=150mm).
% On place A''B'' à gauche de l'oculaire.
\coordinate (App) at (\Ltube-1.5, 1.2); % Position approximative pour le dessin
\coordinate (Bpp) at (\Ltube-1.5, 0);
\draw[poppurple, very thick] (App) -- (Bpp) node[midway, left] {$A''$} node[below] {$B''$};
\draw[<->, poppurple] (Hoc) -- (App) node[midway, above] {$p'' \approx -D$};

% --- Œil ---
\draw[popdark, fill=popcream] (\Ltube-2.5,0) ellipse (0.4 and 0.25);
\draw[popdark] (\Ltube-2.5,0) circle (0.15);
\node[below] at (\Ltube-2.5,-0.4) {Œil accommodé};

\node[align=center, popdark] at (5, -3.2) {Fonctionnement au confort : image finale au point proche $D$};
\end{tikzpicture}
\legende{Construction géométrique en fonctionnement au confort. L'image intermédiaire $A'B'$ est placée légèrement en avant de $F_{oc}$ ($p'_{oc} < f'_{oc}$). L'oculaire donne une image finale virtuelle $A''B''$ au point proche $D = \SI{25}{cm}$ (distance algébrique $p''_{oc} \approx -D$). L'image finale est inversée par rapport à l'objet initial $AB$.}
\end{popfigure}

\courssub{Fonctionnement normal (image à l'infini)}
L'image intermédiaire est \textbf{exactement} au foyer objet de l'oculaire : $p'_{oc} = f'_{oc}$.
Les rayons émergents sont parallèles. L'œil au repos reçoit un faisceau parallèle. Le grossissement angulaire de l'oculaire est :
\[
G_{oc}^{(\text{norm})} = \frac{\alpha'}{\alpha} = \frac{250\,\text{mm}}{f'_{oc}(\text{mm})} = \frac{D}{f'_{oc}}
\]
où $\alpha = \frac{A'B'}{f'_{oc}}$ est l'angle sous-tendu par l'image intermédiaire au centre de l'oculaire, et $\alpha'$ l'angle sous-tendu par l'image finale à l'infini.

\courssub{Fonctionnement au confort (image au point proche)}
L'image intermédiaire est entre $F_{oc}$ et $H_{oc}$ : $p'_{oc} < f'_{oc}$.
La relation de conjugaison $\frac{1}{p'_{oc}} + \frac{1}{p''_{oc}} = \frac{1}{f'_{oc}}$ avec $p''_{oc} = -D$ (image virtuelle à gauche) donne :
\[
p'_{oc} = \frac{f'_{oc} D}{D + f'_{oc}}
\]
Le grossissement angulaire (ou commercial) de l'oculaire devient :
\[
G_{oc}^{(\text{confort})} = \frac{\alpha'}{\alpha} = \frac{D}{p'_{oc}} = \frac{D}{f'_{oc}} + 1 = G_{oc}^{(\text{norm})} + 1
\]
\begin{attention}[Convention de grossissement de l'oculaire]
Au Probatoire, on utilise presque toujours la \textbf{valeur commerciale gravée sur l'oculaire} (ex. $\times 10$), qui correspond au fonctionnement normal : $G_{oc} = \frac{250}{f'_{oc}(\text{mm})}$. Si le sujet impose le fonctionnement au confort, on utilise $G_{oc} = \frac{250}{f'_{oc}(\text{mm})} + 1$. Précisez toujours le mode utilisé.
\end{attention}

\coursec{Grossissement commercial et Puissance commerciale}

\courssub{Grossissement commercial total $G$}

Le grossissement commercial (ou pouvoir grossissant) d'un microscope est le produit du grossissement commercial de l'objectif par celui de l'oculaire. C'est un nombre \textbf{sans unité}.

\begin{propriete}[Grossissement commercial du microscope]
\[
\boxed{G = G_{obj} \times G_{oc}}
\]
\begin{itemize}
    \item \textbf{Objectif :} $G_{obj} = \frac{\Delta}{f'_{obj}}$ (valeur absolue). Par convention constructeur (standard DIN/ISO, tube mécanique $L = \SI{160}{mm}$), on utilise souvent $G_{obj} = \frac{160}{f'_{obj}(\text{mm})}$. La valeur gravée sur l'objectif (ex. $40\times$) est cette valeur normalisée.
    \item \textbf{Oculaire :} $G_{oc} = \frac{250}{f'_{oc}(\text{mm})}$ (fonctionnement normal) ou $G_{oc} = \frac{250}{f'_{oc}(\text{mm})} + 1$ (fonctionnement au confort). La valeur gravée (ex. $10\times$) correspond au fonctionnement normal.
\end{itemize}
\end{propriete}

\begin{aretenir}[Formules clés à connaître par cœur]
\begin{align*}
\text{Grossissement objectif (normalisé)} &\quad G_{obj} = \frac{160}{f'_{obj}(\text{mm})} \\
\text{Grossissement oculaire (normal)} &\quad G_{oc} = \frac{250}{f'_{oc}(\text{mm})} \\
\text{Grossissement total (normal)} &\quad G = \frac{160 \times 250}{f'_{obj}(\text{mm}) \times f'_{oc}(\text{mm})} = \frac{40000}{f'_{obj}(\text{mm}) \, f'_{oc}(\text{mm})} \\
\text{Grossissement total (confort)} &\quad G = G_{obj} \times \left( \frac{250}{f'_{oc}(\text{mm})} + 1 \right)
\end{align*}
\textbf{Analyse dimensionnelle :} $f'$ en mm, $160$ et $250$ en mm $\rightarrow$ $G$ sans unité. \textbf{Vérifié.}
\end{aretenir}

\courssub{Puissance commerciale et puissance optique (dioptries)}

Le programme distingue le grossissement (adimensionnel) de la puissance (en dioptries).

\begin{definition}[Puissance optique d'un système centré]
Pour un système de deux lentilles minces (objectif + oculaire) séparées par une distance $d = \overline{H_{obj}H_{oc}}$ (distance entre plans principaux), la puissance équivalente (vergérence) est :
\[
P = P_{obj} + P_{oc} - d \, P_{obj} P_{oc}
\]
L'unité est le \textbf{dioptre} ($\si{\dioptre} = \text{m}^{-1}$).
\end{definition}

\begin{attention}[Distinction cruciale Probatoire]
\begin{itemize}
    \item \textbf{Grossissement commercial $G$} : nombre pur (ex. $\times 400$). Caractérise l'agrandissement angulaire.
    \item \textbf{Puissance optique $P$} : en dioptries (ex. $\SI{250}{\dioptre}$). Caractérise la vergérence du système équivalent.
    \item \textbf{Puissance de l'objectif $P_{obj}$} : en dioptries ($1/f'_{obj}$). Ne pas confondre avec $G_{obj}$.
\end{itemize}
Ne mélangez jamais ces grandeurs dans une même formule.
\end{attention}

\begin{exemplebox}[Calcul de la puissance optique du microscope]
Reprenons le microscope de l'exemple résolu : $f'_{obj} = \SI{4,0}{mm}$, $f'_{oc} = \SI{25}{mm}$, $L = \SI{160}{mm}$.
$d = \overline{H_{obj}H_{oc}} = L = \SI{0,160}{m}$.
$P_{obj} = \frac{1}{0,004} = \SI{250}{\dioptre}$.
$P_{oc} = \frac{1}{0,025} = \SI{40}{\dioptre}$.
$P = 250 + 40 - 0,160 \times 250 \times 40 = 290 - 1600 = \SI{-1310}{\dioptre}$.
La puissance est négative : le système microscope est \textbf{divergent} globalement (il produit une image virtuelle). C'est normal pour un instrument de vision.
\end{exemplebox}

\coursec{Mise au point : théorie et pratique}

\courssub{Principe de la mise au point}

La netteté de l'image finale exige que l'image intermédiaire $A'B'$ soit située \textbf{exactement} dans le plan focal objet de l'oculaire ($F_{oc}$) pour le fonctionnement normal, ou à la distance $p'_{oc}$ calculée pour le confort.
Comme la position de l'oculaire est fixe (souvent), c'est la position de l'objet (ou de l'objectif) qui doit être ajustée pour que l'image intermédiaire tombe au bon endroit.
La relation de conjugaison pour l'objectif : $\frac{1}{p} + \frac{1}{p'} = \frac{1}{f'_{obj}}$.
Avec $p' \approx \Delta + f'_{obj}$ (distance $H_{obj}A' \approx \Delta + f'_{obj}$), on en déduit la distance objet $p$ requise.
Comme $f'_{obj}$ est très petit (quelques mm), $p \approx f'_{obj}$. La \textbf{profondeur de champ} (intervalle de positions de l'objet donnant une image acceptablement nette) est de l'ordre du $\mu m$.

\courssub{Les deux vis de mise au point}

\begin{center}
\begin{tabularx}{\linewidth}{|l|X|X|}\hline
\rowcolor{popblueL}
\textbf{Vis} & \textbf{Action mécanique} & \textbf{Rôle optique} \\ \hline
\textbf{Macrométrique} & Déplace le tube (ou la platine) de $\sim \SI{1}{mm}$ par tour. & \textbf{Approche rapide} : amène l'objet dans la zone $f'_{obj} \pm \SI{1}{mm}$. Image floue mais visible. \\ \hline
\textbf{Micrométrique} & Déplace de $\sim \SI{0,01}{mm}$ à $\SI{0,002}{mm}$ par tour. & \textbf{Réglage fin} : compense la faible profondeur de champ pour obtenir la netteté parfaite. \\ \hline
\end{tabularx}
\end{center}

\begin{methode}[Procédure standard de mise au point (Probatoire)]
\begin{enumerate}
    \item Placer la préparation sur la platine, centrer grossièrement sur la zone d'intérêt.
    \item Sélectionner l'objectif de plus faible grossissement ($\times 4$ ou $\times 10$).
    \item \textbf{Vis macrométrique} : En regardant \textbf{de côté} (pas dans l'oculaire !), descendre l'objectif (ou monter la platine) au maximum sans toucher la lame (distance de travail).
    \item Regarder dans l'oculaire. Tourner la vis macrométrique en \textbf{remontant} l'objectif (ou descendant la platine) jusqu'à l'apparition d'une image (même floue).
    \item \textbf{Vis micrométrique} : Affiner lentement jusqu'à la netteté maximale.
    \item Changer d'objectif (revolver) pour grossir. \textbf{Ne toucher qu'à la vis micrométrique} (parfocalité des objectifs corrigés).
    \item Ajuster le diaphragme du condenseur (ouverture numérique) pour le contraste/résolution.
\end{enumerate}
\end{methode}

\begin{attention}[Erreurs pratiques fréquentes]
\begin{itemize}
    \item Descendre l'objectif \textbf{en regardant dans l'oculaire} : risque de casser la lame ou rayer la lentille frontale (distance de travail $\sim \SI{0,1}{mm}$ pour $\times 100$).
    \item Utiliser la macrométrique après avoir changé d'objectif pour un fort grossissement : l'image « saute » hors du champ, on perd la mise au point, risque de casse.
    \item Oublier de refermer le diaphragme du condenseur pour les faibles grossissements (éblouissement, perte de contraste).
\end{itemize}
\end{attention}

\coursec{Longueur du tube, encombrement et limites (Complément Probatoire)}

\courssub{Longueur mécanique vs Intervalle inter-foculaire}

\begin{propriete}[Relations géométriques]
Soit $L = \overline{H_{obj}H_{oc}}$ la distance entre plans principaux (longueur optique du tube, $\approx$ longueur mécanique normalisée $\SI{160}{mm}$).
\[
\Delta = \overline{F'_{obj}F_{oc}} = L - f'_{obj} - f'_{oc}
\]
\begin{proof}
$F'_{obj}$ est à $f'_{obj}$ à droite de $H_{obj}$. $F_{oc}$ est à $f'_{oc}$ à gauche de $H_{oc}$.
$\Delta = L - f'_{obj} - f'_{oc}$.
\end{proof}
\end{propriete}

\begin{aretenir}[Valeurs types (Standard DIN 160 mm)]
\begin{center}
\begin{tabular}{|c|c|c|c|c|}\hline
\textbf{Objectif} & $f'_{obj}$ (mm) & $G_{obj}$ (normé) & $\Delta$ (mm) & Distance travail (mm) \\ \hline
$\times 4$ & 40 & 4 & $\approx 116$ & $\sim 18$ \\ \hline
$\times 10$ & 16 & 10 & $\approx 119$ & $\sim 5$ \\ \hline
$\times 40$ & 4 & 40 & $\approx 131$ & $\sim 0,5$ \\ \hline
$\times 100$ (huile) & 1,6 & 100 & $\approx 133$ & $\sim 0,13$ \\ \hline
\end{tabular}
\end{center}
\textbf{Remarque :} Les microscopes modernes « à tube infini » (infinity-corrected) ont un objectif qui donne une image à l'infini ; un tube-lentille (tube lens) forme l'image intermédiaire. $\Delta$ n'est plus la grandeur de référence, mais le grossissement commercial reste gravé selon la convention $G_{obj} = 160/f'_{obj}$.
\end{aretenir}

\courssub{Limites : Résolution et Aberrations}

\begin{propriete}[Pouvoir de résolution (Critère d'Abbe - Rayleigh)]
La distance minimale $\delta_{min}$ entre deux points distincts sur l'objet est :
\[
\delta_{min} = \frac{0,61 \lambda}{\text{NA}} = \frac{0,61 \lambda}{n \sin u'_{max}}
\]
\begin{itemize}
    \item $\lambda$ : longueur d'onde dans le vide (ex. $\SI{550}{nm}$ lumière verte).
    \item $n$ : indice du milieu (air $1$, huile $\num{1,51}$).
    \item $u'_{max}$ : demi-angle d'ouverture de l'objectif.
\end{itemize}
\end{propriete}

\begin{exemplebox}[Calcul de résolution]
Objectif $\times 40$, NA $= 0,65$ (air). $\lambda = \SI{550}{nm}$.
$\delta_{min} = \frac{0,61 \times 550}{0,65} \approx \SI{516}{nm} = \SI{0,52}{\mu m}$.
Objectif $\times 100$, NA $= 1,30$ (huile). $\delta_{min} = \frac{0,61 \times 550}{1,30} \approx \SI{258}{nm} = \SI{0,26}{\mu m}$.
\textbf{Conclusion :} L'immersion d'huile ($n>1$) augmente NA et améliore la résolution (divise $\delta_{min}$ par $\sim 1,5$).
\end{exemplebox}

\begin{attention}[Grossissement utile vs Grossissement vide]
Augmenter $G$ au-delà de la limite de résolution ne révèle pas plus de détails (grossissement vide).
Règle empirique : $G_{max\,utile} \approx 500 \times \text{NA}$ à $1000 \times \text{NA}$.
Pour NA $= 0,65$ ($\times 40$), $G_{max} \approx 300$ à $650$. Un oculaire $\times 10$ donne $G=400$ (correct). Un oculaire $\times 15$ donnerait $G=600$ (limite). Un oculaire $\times 20$ serait du grossissement vide.
\end{attention}

\coursec{Résolution de problèmes types (Probatoire)}

\begin{exoresolu}[Type 1 : Détermination des caractéristiques à partir des gravures]
Un microscope possède un tube de $\SI{160}{mm}$. L'objectif porte la gravure $\times 40 / 0,65 / 160 / 0,17$. L'oculaire porte $\times 10 / 18$.
\begin{enumerate}
    \item Interpréter les gravures de l'objectif.
    \item Calculer les focales $f'_{obj}$ et $f'_{oc}$.
    \item Calculer l'intervalle inter-foculaire $\Delta$.
    \item Déterminer le grossissement commercial total en fonctionnement normal.
    \item Calculer la puissance optique $P$ du système (en dioptries).
\end{enumerate}
\textbf{Solution détaillée}
\begin{enumerate}
    \item $\times 40$ : grossissement commercial normalisé. $0,65$ : ouverture numérique (NA). $160$ : longueur de tube de référence (mm). $0,17$ : épaisseur de lame de couvre-objet prévue pour la correction (mm).
    \item $G_{obj} = 40 = \frac{160}{f'_{obj}(\text{mm})} \Rightarrow f'_{obj} = \SI{4,0}{mm}$.
        $G_{oc} = 10 = \frac{250}{f'_{oc}(\text{mm})} \Rightarrow f'_{oc} = \SI{25,0}{mm}$.
    \item $\Delta = L - f'_{obj} - f'_{oc} = 160 - 4 - 25 = \SI{131}{mm}$.
    \item $G = G_{obj} \times G_{oc} = 40 \times 10 = \textbf{400}$ (sans unité).
    \item $d = L = \SI{0,160}{m}$. $P_{obj} = 1/0,004 = \SI{250}{\dioptre}$. $P_{oc} = 1/0,025 = \SI{40}{\dioptre}$.
        $P = 250 + 40 - 0,160 \times 250 \times 40 = 290 - 1600 = \textbf{\SI{-1310}{\dioptre}}$.
\end{enumerate}
\end{exoresolu}

\begin{exoresolu}[Type 2 : Mise au point et distance objet]
Un objectif de focale $f'_{obj} = \SI{4,0}{mm}$ est utilisé sur un tube $L = \SI{160}{mm}$ avec un oculaire $f'_{oc} = \SI{25}{mm}$. On veut observer en fonctionnement normal.
\begin{enumerate}
    \item Calculer la distance $\overline{H_{obj}A}$ (distance objet-objectif) pour obtenir l'image intermédiaire en $F_{oc}$.
    \item L'élève tourne la vis micrométrique d'un demi-tour (pas $\SI{0,01}{mm}$/tour). De combien la position de l'image intermédiaire se décale-t-elle ? L'image finale reste-t-elle nette ?
\end{enumerate}
\textbf{Solution détaillée}
\begin{enumerate}
    \item Fonctionnement normal $\Rightarrow$ Image intermédiaire en $F_{oc}$.
        $p' = \overline{H_{obj}A'} = L - f'_{oc} = 160 - 25 = \SI{135}{mm}$.
        Relation de Descartes : $\frac{1}{p} + \frac{1}{p'} = \frac{1}{f'_{obj}}$.
        $\frac{1}{p} = \frac{1}{4} - \frac{1}{135} = \frac{135 - 4}{540} = \frac{131}{540}$.
        $p = \frac{540}{131} \approx \SI{4,12}{mm}$.
        Vérification : $p \approx f'_{obj} + \varepsilon$ avec $\varepsilon \approx \SI{0,12}{mm}$. Correct.
    \item Déplacement objet $\delta p = \pm \SI{0,005}{mm}$ (demi-tour).
        Relation de Newton : $x x' = f'^2$. $x = p - f' \approx \SI{0,12}{mm}$. $x' = \frac{16}{0,12} \approx \SI{133}{mm}$.
        Différenciation : $\frac{\delta x'}{x'} = -\frac{\delta x}{x}$.
        $\delta x = \delta p = \SI{0,005}{mm}$.
        $|\delta x'| = x' \frac{|\delta x|}{x} = 133 \times \frac{0,005}{0,12} \approx \SI{5,5}{mm}$.
        L'image intermédiaire se déplace de $\sim \SI{5,5}{mm}$ ! Elle sort largement du plan $F_{oc}$ (profondeur de champ oculaire $\sim \mu m$). L'image finale devient \textbf{floue}.
        \textbf{Conclusion :} La micrométrie est indispensable ; la sensibilité est extrême.
\end{enumerate}
\end{exoresolu}

\begin{exoresolu}[Type 3 : Choix d'instrument pour une application]
On souhaite observer des bactéries de taille $\sim \SI{1}{\mu m}$ (ex. \emph{Mycobacterium tuberculosis}) au laboratoire de l'Hôpital Central de Yaoundé. Deux microscopes sont disponibles :
\begin{itemize}
    \item Micro A : Objectif $\times 40$ (NA $0,65$), Oculaire $\times 10$.
    \item Micro B : Objectif $\times 100$ (NA $1,30$, huile), Oculaire $\times 10$.
\end{itemize}
Lequel choisir ? Justifier par le calcul de la résolution et du grossissement utile.
\textbf{Solution détaillée}
\begin{itemize}
    \item \textbf{Micro A :} $\delta_{min} = \frac{0,61 \times 550}{0,65} \approx \SI{516}{nm} = \SI{0,52}{\mu m}$. $G = 400$.
        Grossissement utile max $\approx 1000 \times 0,65 = 650$. $G=400$ est dans la plage utile.
        Mais $\delta_{min} = 0,52 \mu m < 1 \mu m$ : \textbf{théoriquement résolvable}.
    \item \textbf{Micro B :} $\delta_{min} = \frac{0,61 \times 550}{1,30} \approx \SI{258}{nm} = \SI{0,26}{\mu m}$. $G = 1000$.
        Grossissement utile max $\approx 1000 \times 1,3 = 1300$. $G=1000$ est optimal.
        $\delta_{min} = 0,26 \mu m \ll 1 \mu m$ : \textbf{résolution excellente}, détails morphologiques visibles.
\end{itemize}
\textbf{Choix :} **Microscope B (Objectif $\times 100$ huile)**. C'est la référence pour la bactériologie (diagnostic TB, méningite). L'objectif $\times 40$ permet le dépistage, mais la confirmation morphologique exige le $\times 100$ huile.
\end{exoresolu}

\begin{exercice}[Entraînement Probatoire - Microscope]
\begin{enumerate}
    \item Un microscope a un tube de $\SI{160}{mm}$. L'objectif est un $\times 10$ (NA $0,25$), l'oculaire un $\times 15$.
        \begin{enumerate}
            \item Calculer $f'_{obj}$ et $f'_{oc}$.
            \item Calculer $\Delta$.
            \item Calculer $G$ (normal) et $G$ (confort).
            \item Calculer la puissance optique $P$ du système.
        \end{enumerate}
    \item On observe un cheveu (diamètre $\SI{70}{\mu m}$) avec un objectif $\times 4$ ($f'_{obj}=\SI{40}{mm}$) et oculaire $\times 10$ ($f'_{oc}=\SI{25}{mm}$), tube $\SI{160}{mm}$.
        \begin{enumerate}
            \item Taille de l'image intermédiaire ?
            \item Angle sous-tendu par l'image finale à l'infini ?
        \end{enumerate}
    \item Un objectif $\times 40$ (NA $0,65$) est utilisé à sec (air) puis avec huile ($n=1,51$, NA max $1,4$).
        Calculer le gain de résolution (rapport $\delta_{air}/\delta_{huile}$) pour $\lambda = \SI{500}{nm}$.
\end{enumerate}
\end{exercice}

\begin{corrige}[Exercice Entraînement]
\begin{enumerate}
    \item \textbf{Données :} $L=160$ mm. $G_{obj}=10$, $G_{oc}=15$.
        \textbf{(a)} $f'_{obj} = 160/10 = \SI{16}{mm}$. $f'_{oc} = 250/15 \approx \SI{16,7}{mm}$.
        \textbf{(b)} $\Delta = 160 - 16 - 16,7 = \SI{127,3}{mm}$.
        \textbf{(c)} $G_{norm} = 10 \times 15 = 150$. $G_{confort} = 10 \times (15+1) = 160$.
        \textbf{(d)} $P_{obj}=1/0,016=\SI{62,5}{\dioptre}$. $P_{oc}=1/0,0167=\SI{60}{\dioptre}$. $d=0,160$ m.
        $P = 62,5 + 60 - 0,160 \times 62,5 \times 60 = 122,5 - 600 = \SI{-477,5}{\dioptre}$.
    \item \textbf{Données :} $d_{cheveu}=70 \mu m$. $f'_{obj}=40$ mm ($G_{obj}=4$). $f'_{oc}=25$ mm ($G_{oc}=10$).
        \textbf{(a)} $\gamma_{obj} \approx -\Delta/f'_{obj}$. $\Delta = 160 - 40 - 25 = 95$ mm. $|\gamma_{obj}| \approx 95/40 = 2,375$. (Valeur normalisée $G_{obj}=4$).
        Utilisons la valeur normalisée pour la taille image intermédiaire standard : $d' = 4 \times 70 = \SI{280}{\mu m} = \SI{0,28}{mm}$.
        \textbf{(b)} $G = 40$. $\alpha' = G \times \frac{d}{250} = 40 \times \frac{70 \times 10^{-3}}{250} = 40 \times 2,8 \times 10^{-4} = \SI{11,2}{mrad}$.
    \item \textbf{Données :} $\lambda=500$ nm. Air : NA $=0,65$. Huile : NA $=1,40$ (valeur max théorique pour $\times 100$, pour $\times 40$ huile rare, supposons NA passe à $1,4$ pour l'exercice ou restons sur $\times 100$). L'énoncé dit "Objectif $\times 40$... puis avec huile". Un objectif $\times 40$ n'est généralement pas à immersion. Supposons un objectif $\times 100$ huile pour la comparaison standard, ou calculons avec NA huile $=1,4$ pour le même grossissement.
        $\delta_{air} = 0,61 \times 500 / 0,65 \approx 469$ nm.
        $\delta_{huile} = 0,61 \times 500 / 1,40 \approx 218$ nm.
        Gain = $469 / 218 \approx 2,15$. La résolution est \textbf{doublée}.
\end{enumerate}
\end{corrige}

\coursec{Schéma de principe et construction géométrique}

\noindent
Après avoir identifié les deux lentilles constitutives du microscope (l'objectif et l'oculaire) et leur rôle respectif dans la section précédente, nous allons maintenant comprendre \emph{comment} ces deux lentilles coopèrent pour produire une image finale grandie et inversée. La clé réside dans la position relative de l'image intermédiaire formée par l'objectif par rapport au foyer de l'oculaire. C'est cette géométrie précise qui définit le « fonctionnement normal » (image à l'infini, repos de l'œil) et le « fonctionnement confort » (image à \SI{25}{cm}, point proche).

\courssub{Choix de l'objet et action de l'objectif : formation de l'image intermédiaire}

\noindent
Rappelons la condition fondamentale : l'objet \textbf{AB} est placé \emph{très légèrement} au-delà du foyer objet de l'objectif $F_{\text{obj}}$. En pratique, la distance $p_{\text{obj}} = \overline{F_{\text{obj}}A}$ est de l'ordre du dixième de millimètre, alors que la focale de l'objectif $f'_{\text{obj}}$ vaut quelques millimètres (\SI{4}{mm}, \SI{10}{mm}, \SI{40}{mm}...). On a donc $p_{\text{obj}} \approx f'_{\text{obj}}$ et l'objet est \emph{quasi} au foyer.

\begin{popfigure}
\begin{tikzpicture}[scale=0.9, >=Stealth]
    % Axes
    \draw[popline, thick] (-5,0) -- (7,0) node[right] {Axe optique};
    % Objectif
    \draw[popblue, very thick] (0,-1.2) -- (0,1.2) node[above] {$L_{\text{obj}}$};
    \node[popblue] at (0,1.5) {$O_{\text{obj}}$};
    % Foyers objectif
    \draw[popblue!70] (-2.5,-0.1) -- (-2.5,0.1) node[below] {$F_{\text{obj}}$};
    \draw[popblue!70] (2.5,-0.1) -- (2.5,0.1) node[below] {$F'_{\text{obj}}$};
    % Objet AB
    \draw[popdark, very thick] (-2.7,0) -- (-2.7,1.5) node[above] {$B$};
    \node[popdark] at (-2.7,-0.3) {$A$};
    % Rayon 1 : parallèle à l'axe -> passe par F'obj
    \draw[poporange, ->, thick] (-2.7,1.5) -- (0,1.5) node[midway, above] {1};
    \draw[poporange, ->, thick] (0,1.5) -- (2.5,0);
    % Rayon 2 : passe par centre optique -> non dévié
    \draw[poporange, ->, thick] (-2.7,1.5) -- (0,0) node[midway, below left] {2};
    \draw[poporange, ->, thick] (0,0) -- (6,3.5);
    % Image intermédiaire A'B'
    \draw[popgreen, very thick] (5.5,-0.1) -- (5.5,-3.2) node[below] {$B'$};
    \node[popgreen] at (5.5,0.2) {$A'$};
    % Flèches
    \draw[<->, popmuted] (-2.7, -1.8) -- (0, -1.8) node[midway, below] {$p_{\text{obj}} \approx f'_{\text{obj}}$};
    \draw[<->, popmuted] (0, -2.3) -- (5.5, -2.3) node[midway, below] {$p'_{\text{obj}} \approx \Delta + f'_{\text{obj}}$};
    % Delta
    \draw[<->, poppurple, thick] (2.5, -1.2) -- (6.5, -1.2) node[midway, below] {$\Delta$ (\cle{intervalle optique})};
    % Foyer oculaire (position future)
    \draw[poppurple!70] (6.5,-0.1) -- (6.5,0.1) node[above] {$F_{\text{oc}}$};
\end{tikzpicture}
\legende{Construction géométrique de l'image intermédiaire $A'B'$ par l'objectif. L'objet $AB$ est placé juste après $F_{\text{obj}}$. L'image $A'B'$ est réelle, inversée et grandie. Elle se forme au voisinage de $F_{\text{oc}}$ (foyer objet de l'oculaire).}
\end{popfigure}

\begin{propriete}[Relation de Newton pour l'objectif]
Pour l'objectif, notons $x_{\text{obj}} = \overline{F_{\text{obj}}A}$ (distance objet/foyer objet, positive car $A$ est après $F_{\text{obj}}$) et $x'_{\text{obj}} = \overline{F'_{\text{obj}}A'}$ (distance image/foyer image).
La relation de Newton s'écrit :
\[
x_{\text{obj}} \times x'_{\text{obj}} = {f'_{\text{obj}}}^2
\]
Comme l'objet est placé \emph{très près} du foyer ($x_{\text{obj}} \ll f'_{\text{obj}}$), l'image se forme \emph{très loin} du foyer image ($x'_{\text{obj}} \gg f'_{\text{obj}}$). C'est le principe de l'agrandissement : un petit déplacement de l'objet près du foyer produit un grand déplacement de l'image.
\end{propriete}

\begin{propriete}[Grandissement de l'objectif — Démonstration guidée]
Le grandissement transverse (algébrique) de l'objectif est $\gamma_{\text{obj}} = -\dfrac{A'B'}{AB} = -\dfrac{p'_{\text{obj}}}{p_{\text{obj}}}$.
\begin{enumerate}
    \item On a $p_{\text{obj}} = f'_{\text{obj}} + x_{\text{obj}} \approx f'_{\text{obj}}$ (car $x_{\text{obj}}$ est très petit).
    \item L'image intermédiaire $A'B'$ doit se former \emph{au plan focal objet de l'oculaire} $F_{\text{oc}}$ pour le fonctionnement normal. La distance entre les foyers image de l'objectif et objet de l'oculaire est l'\cle{intervalle optique} $\Delta = \overline{F'_{\text{obj}}F_{\text{oc}}}$.
    \item Donc $p'_{\text{obj}} = \overline{O_{\text{obj}}A'} \approx \overline{F'_{\text{obj}}F_{\text{oc}}} + \overline{F'_{\text{obj}}O_{\text{obj}}} = \Delta + f'_{\text{obj}}$.
    \item Comme $\Delta$ (de l'ordre de \SI{16}{cm} = \SI{160}{mm}) est très grand devant $f'_{\text{obj}}$ (quelques mm), on a $p'_{\text{obj}} \approx \Delta$.
\end{enumerate}
Finalement :
\[
\boxed{\gamma_{\text{obj}} \approx -\frac{\Delta}{f'_{\text{obj}}}}
\]
Le signe négatif confirme que l'image intermédiaire $A'B'$ est \textbf{inversée} par rapport à l'objet $AB$.
\end{propriete}

\begin{attention}[Piège classique : Confondre $\Delta$ et $L$]
L'intervalle optique $\Delta = \overline{F'_{\text{obj}}F_{\text{oc}}}$ est une constante \textbf{géométrique} de l'instrument (standardisée à \SI{160}{mm} pour les microscopes « tube fini »).
La longueur mécanique du tube $L = \overline{O_{\text{obj}}O_{\text{oc}}}$ vaut $L = \Delta + f'_{\text{obj}} + f'_{\text{oc}}$.
N'utilisez \textbf{jamais} $L$ dans la formule du grandissement de l'objectif, mais bien $\Delta$.
\end{attention}

\courssub{Action de l'oculaire : de l'image intermédiaire à l'image finale}

\noindent
L'image intermédiaire $A'B'$ devient l'objet de l'oculaire. Sa position par rapport au foyer objet de l'oculaire $F_{\text{oc}}$ dicte le mode de fonctionnement.

\begin{popfigure}
\begin{tikzpicture}[scale=0.9, >=Stealth]
    % Axes
    \draw[popline, thick] (0,0) -- (10,0) node[right] {Axe optique};
    % Oculaire
    \draw[popblue, very thick] (7,-1.2) -- (7,1.2) node[above] {$L_{\text{oc}}$};
    \node[popblue] at (7,1.5) {$O_{\text{oc}}$};
    % Foyers oculaire
    \draw[popblue!70] (4.5,-0.1) -- (4.5,0.1) node[below] {$F_{\text{oc}}$};
    \draw[popblue!70] (9.5,-0.1) -- (9.5,0.1) node[below] {$F'_{\text{oc}}$};
    % Image intermédiaire A'B' (Objet pour l'oculaire)
    \draw[popgreen, very thick] (4.5,-0.1) -- (4.5,-2.5) node[below] {$B'$};
    \node[popgreen] at (4.5,0.2) {$A'$};
    % Rayons sortants PARALLELES (Fonctionnement Normal)
    \draw[poporange, ->, thick] (4.5,-2.5) -- (7,-2.5) node[midway, below] {1};
    \draw[poporange, ->, thick] (7,-2.5) -- (10,-2.5) node[midway, above] {Sortie $\parallel$ axe};
    \draw[poporange, ->, thick] (4.5,-2.5) -- (7,0) node[midway, above right] {2};
    \draw[poporange, ->, thick] (7,0) -- (10,0);
    % Oeil
    \draw[popdark] (10.5,0) ellipse (0.5 and 0.3);
    \node at (10.5,0) {$\text{Œil}$};
    \node at (10.5,-0.7) {Relâché (infini)};
    % Legende
    \node[poporange] at (5.5, -3.2) {Fonctionnement NORMAL : $A' \equiv F_{\text{oc}}$};
\end{tikzpicture}
\legende{Fonctionnement normal : L'image intermédiaire $A'B'$ est \textbf{confondue} avec le foyer objet $F_{\text{oc}}$. Les rayons émergents de l'oculaire sont parallèles. L'image finale $A''B''$ est à l'infini. L'œil est au repos (accommodation nulle).}
\end{popfigure}

\begin{propriete}[Grandissement angulaire de l'oculaire (Rappel Loupe)]
L'oculaire fonctionne comme une loupe observant l'image intermédiaire $A'B'$.
\begin{itemize}
    \item \textbf{Fonctionnement normal (image à l'infini) :} Le \cle{grandissement commercial} (angulaire) est
    \[
    G_{\text{oc}}^{\text{(normal)}} = \frac{\SI{25}{cm}}{f'_{\text{oc}}} = \frac{\SI{250}{mm}}{f'_{\text{oc}}(\text{mm})}
    \]
    \item \textbf{Fonctionnement confort (image au point proche $D = \SI{25}{cm}$) :} On décale légèrement l'objectif (vis micrométrique) pour que $A'B'$ soit \emph{entre} $F_{\text{oc}}$ et l'oculaire ($p_{\text{oc}} < f'_{\text{oc}}$). L'oculaire donne une image virtuelle $A''B''$ à $p'_{\text{oc}} = -\SI{25}{cm}$. Le grandissement vaut alors
    \[
    G_{\text{oc}}^{\text{(confort)}} = \frac{\SI{25}{cm}}{f'_{\text{oc}}} + 1
    \]
\end{itemize}
Dans les deux cas, $\gamma_{\text{oc}} > 0$ : l'oculaire \textbf{ne ré-inverse pas} l'image. L'image finale $A''B''$ a le même sens que $A'B'$ (donc inversée par rapport à $AB$ initial).
\end{propriete}

\begin{attention}[Erreur fatale au Probatoire : Le « +1 » du microscope]
Beaucoup d'élèves appliquent systématiquement $G = \frac{25}{f'} + 1$ (formule de la loupe au point proche) au microscope.
\begin{itemize}
    \item Si l'énoncé précise « \textbf{fonctionnement normal} » ou « \textbf{œil au repos} » ou « \textbf{image à l'infini} » $\rightarrow$ $G_{\text{oc}} = \frac{25}{f'_{\text{oc}}}$ \textbf{SANS le +1}.
    \item Si l'énoncé précise « \textbf{fonctionnement confort} » ou « \textbf{image au point proche} » ou « \textbf{œil accommodé} » $\rightarrow$ $G_{\text{oc}} = \frac{25}{f'_{\text{oc}}} + 1$ \textbf{AVEC le +1}.
\end{itemize}
Lisez attentivement l'énoncé !
\end{attention}

\begin{popfigure}
\begin{tikzpicture}[scale=0.9, >=Stealth]
    % Axes
    \draw[popline, thick] (0,0) -- (10,0) node[right] {Axe optique};
    % Oculaire
    \draw[popblue, very thick] (7,-1.2) -- (7,1.2) node[above] {$L_{\text{oc}}$};
    \node[popblue] at (7,1.5) {$O_{\text{oc}}$};
    % Foyers oculaire
    \draw[popblue!70] (4.5,-0.1) -- (4.5,0.1) node[below] {$F_{\text{oc}}$};
    \draw[popblue!70] (9.5,-0.1) -- (9.5,0.1) node[below] {$F'_{\text{oc}}$};
    % Image intermédiaire A'B' (Objet pour l'oculaire) -> DECALEE vers l'oculaire
    \draw[popgreen, very thick] (5.5,-0.1) -- (5.5,-2.0) node[below] {$B'$};
    \node[popgreen] at (5.5,0.2) {$A'$};
    % Rayons sortants DIVERGENTS (Fonctionnement Confort)
    \draw[poporange, ->, thick] (5.5,-2.0) -- (7,-1.5) node[midway, below] {1};
    \draw[poporange, ->, thick] (7,-1.5) -- (10,-0.5);
    \draw[dashed, poporange] (10,-0.5) -- (4.5, -3.5); % prolongation arrière
    \draw[poporange, ->, thick] (5.5,-2.0) -- (7,0) node[midway, above right] {2};
    \draw[poporange, ->, thick] (7,0) -- (10,0);
    \draw[dashed, poporange] (10,0) -- (4.5, -3.5);
    % Image finale virtuelle A''B''
    \draw[poppurple, very thick] (4.5,-3.5) -- (4.5,-0.1) node[above] {$B''$};
    \node[poppurple] at (4.5,-3.8) {$A''$};
    % Fleche p'oc
    \draw[<->, popmuted] (5.5, -2.5) -- (7, -2.5) node[midway, below] {$p_{\text{oc}} < f'_{\text{oc}}$};
    % Fleche p'oc'
    \draw[<->, popmuted] (4.5, -4.3) -- (7, -4.3) node[midway, below] {$p'_{\text{oc}} = -\SI{25}{cm}$};
    % Oeil
    \draw[popdark] (10.5,0) ellipse (0.5 and 0.3);
    \node at (10.5,0) {$\text{Œil}$};
    \node at (10.5,-0.7) {Accommodé};
    % Legende
    \node[poporange] at (5.5, -5.0) {Fonctionnement CONFORT : $A'B'$ entre $F_{\text{oc}}$ et $L_{\text{oc}}$};
\end{tikzpicture}
\legende{Fonctionnement confort : L'image intermédiaire $A'B'$ est placée entre $F_{\text{oc}}$ et l'oculaire. Les rayons émergents sont divergents ; leurs prolongements se coupent en $A''B''$ (image virtuelle) à \SI{25}{cm}. L'œil doit accommoder.}
\end{popfigure}

\courssub{Schéma global complet et bilan des sens}

\noindent
Assemblons les deux étapes. C'est la figure de référence pour tout problème de microscope.

\begin{popfigure}
\begin{tikzpicture}[scale=0.75, >=Stealth]
    % Echelle : 1cm = 10mm environ. Delta = 16cm = 16 units. fobj=4mm=0.4u. foc=25mm=2.5u.
    % Placement : O_obj à x=0. F'_obj à x=0.4. F_oc à x=16.4. O_oc à x=18.9.
    \def\fobj{0.4}
    \def\foc{2.5}
    \def\Delta{16}
    \def\Oobj{0}
    \def\Fpobj{\fobj}
    \def\Foc{\Fpobj+\Delta}
    \def\Ooc{\Foc+\foc}
    
    % Axe optique
    \draw[popline, thick] (-2,0) -- (\Ooc+3,0) node[right] {Axe optique principal};
    
    % Objectif
    \draw[popblue, very thick] (\Oobj,-2) -- (\Oobj,2) node[above] {$L_{\text{obj}}$};
    \node[popblue] at (\Oobj,2.3) {$O_{\text{obj}}$};
    % Foyers Objectif
    \draw[popblue!70] (\Oobj-\fobj,-0.15) -- (\Oobj-\fobj,0.15) node[below] {$F_{\text{obj}}$};
    \draw[popblue!70] (\Fpobj,-0.15) -- (\Fpobj,0.15) node[below] {$F'_{\text{obj}}$};
    
    % Oculaire
    \draw[popblue, very thick] (\Ooc,-2) -- (\Ooc,2) node[above] {$L_{\text{oc}}$};
    \node[popblue] at (\Ooc,2.3) {$O_{\text{oc}}$};
    % Foyers Oculaire
    \draw[popblue!70] (\Foc,-0.15) -- (\Foc,0.15) node[below] {$F_{\text{oc}}$};
    \draw[popblue!70] (\Ooc+\foc,-0.15) -- (\Ooc+\foc,0.15) node[below] {$F'_{\text{oc}}$};
    
    % Objet AB (juste après F_obj)
    \pgfmathsetmacro{\pobj}{\fobj+0.15}
    \draw[popdark, very thick] (\Oobj-\pobj,0) -- (\Oobj-\pobj,1.2) node[above] {$B$};
    \node[popdark] at (\Oobj-\pobj,-0.3) {$A$};
    
    % Image intermédiaire A'B' (en F_oc pour normal)
    \draw[popgreen, very thick] (\Foc,0) -- (\Foc,-3.5) node[below] {$B'$};
    \node[popgreen] at (\Foc,0.3) {$A'$};
    
    % Rayons Objectif
    % Rayon 1 : Parallele -> F'obj
    \draw[poporange, thick, ->] (\Oobj-\pobj,1.2) -- (\Oobj,1.2);
    \draw[poporange, thick, ->] (\Oobj,1.2) -- (\Fpobj,0);
    \draw[poporange, thick, ->] (\Fpobj,0) -- (\Foc,-3.5);
    % Rayon 2 : Centre optique
    \draw[poporange, thick, ->] (\Oobj-\pobj,1.2) -- (\Oobj,0);
    \draw[poporange, thick, ->] (\Oobj,0) -- (\Foc,-3.5);
    
    % Rayons Oculaire (Sortie parallele = Normal)
    % Rayon 1 (etait parallele a l'axe intermediaire -> passe par F'oc) -> Sort parallèle
    \draw[poporange, thick, ->] (\Foc,-3.5) -- (\Ooc,-3.5);
    \draw[poporange, thick, ->] (\Ooc,-3.5) -- (\Ooc+3,-3.5) node[midway, above] {Sortie $\parallel$};
    % Rayon 2 (passe par O_oc)
    \draw[poporange, thick, ->] (\Foc,-3.5) -- (\Ooc,0);
    \draw[poporange, thick, ->] (\Ooc,0) -- (\Ooc+3,0);
    
    % Cotes
    \draw[<->, poppurple, thick] (\Fpobj,-1.5) -- (\Foc,-1.5) node[midway, below] {$\Delta$ (Intervalle optique)};
    \draw[<->, popmuted] (\Oobj,-2.5) -- (\Ooc,-2.5) node[midway, below] {$L = \Delta + f'_{\text{obj}} + f'_{\text{oc}}$ (Longueur tube)};
    \draw[<->, popmuted] (\Oobj-\pobj, -3.0) -- (\Oobj, -3.0) node[midway, below] {$p_{\text{obj}} \approx f'_{\text{obj}}$};
    \draw[<->, popmuted] (\Oobj, -3.5) -- (\Foc, -3.5) node[midway, below] {$p'_{\text{obj}} \approx \Delta + f'_{\text{obj}}$};
    
    % Oeil
    \draw[popdark] (\Ooc+3.5,0) ellipse (0.6 and 0.4);
    \node at (\Ooc+3.5,0) {$\text{Œil}$};
    \node at (\Ooc+3.5,-0.9) {Relâché};
    
    % Legende sens
    \node[align=center, popdark] at (\Ooc+4, -3.5) {Image finale $A''B''$\\à l'infini\\Inversée par rapport à $AB$};
\end{tikzpicture}
\legende{Schéma global complet du microscope en fonctionnement normal. $\Delta = \overline{F'_{\text{obj}}F_{\text{oc}}}$ est l'intervalle optique. $L = \overline{O_{\text{obj}}O_{\text{oc}}}$ est la longueur mécanique du tube. L'image finale est à l'infini, inversée par rapport à l'objet initial.}
\end{popfigure}

\begin{aretenir}[Bilan des sens et grandissements]
\begin{itemize}
    \item \textbf{Objectif :} $\gamma_{\text{obj}} = -\dfrac{p'_{\text{obj}}}{p_{\text{obj}}} \approx -\dfrac{\Delta}{f'_{\text{obj}}} < 0$. Image intermédiaire $A'B'$ \textbf{réelle, inversée, grandie}.
    \item \textbf{Oculaire :} $G_{\text{oc}} > 0$ (loupe). Image finale $A''B''$ \textbf{virtuelle, même sens que $A'B'$}.
    \item \textbf{Total :} $\gamma_{\text{total}} = \gamma_{\text{obj}} \times G_{\text{oc}} < 0$. Image finale \textbf{inversée} par rapport à l'objet initial $AB$.
    \item \textbf{Grandissement commercial (Puissance) du microscope :}
    \begin{align*}
    \text{Normal : } & G_{\text{mic}} = |\gamma_{\text{obj}}| \times G_{\text{oc}}^{\text{(normal)}} = \frac{\Delta}{f'_{\text{obj}}} \times \frac{\SI{25}{cm}}{f'_{\text{oc}}} \\
    \text{Confort : } & G_{\text{mic}} = |\gamma_{\text{obj}}| \times G_{\text{oc}}^{\text{(confort)}} = \frac{\Delta}{f'_{\text{obj}}} \times \left( \frac{\SI{25}{cm}}{f'_{\text{oc}}} + 1 \right)
    \end{align*}
\end{itemize}
\end{aretenir}

\courssub{Exemple chiffré complet : Microscope standard}

\begin{exemplebox}[Calcul des grandissements — Microscope « standard »]
Un microscope est constitué d'un objectif de focale $f'_{\text{obj}} = \SI{4,0}{mm}$ et d'un oculaire de focale $f'_{\text{oc}} = \SI{25}{mm}$. L'intervalle optique est normalisé à $\Delta = \SI{160}{mm}$.
\begin{enumerate}
    \item Calculer le grandissement de l'objectif $\gamma_{\text{obj}}$.
    \item Calculer le grandissement commercial de l'oculaire en fonctionnement normal et en fonctionnement confort.
    \item Déduire la puissance commerciale (grandissement total) du microscope dans les deux cas.
\end{enumerate}
\textbf{Solution détaillée}
\begin{enumerate}
    \item \textbf{Objectif :}
    \[
    \gamma_{\text{obj}} \approx -\frac{\Delta}{f'_{\text{obj}}} = -\frac{\SI{160}{mm}}{\SI{4,0}{mm}} = -40
    \]
    L'image intermédiaire est réelle, inversée et 40 fois plus grande que l'objet.
    
    \item \textbf{Oculaire :}
    \begin{itemize}
        \item Normal (image à l'infini) : $G_{\text{oc}} = \dfrac{\SI{250}{mm}}{\SI{25}{mm}} = 10$.
        \item Confort (image à \SI{25}{cm}) : $G_{\text{oc}} = \dfrac{\SI{250}{mm}}{\SI{25}{mm}} + 1 = 10 + 1 = 11$.
    \end{itemize}
    
    \item \textbf{Microscope total :}
    \begin{itemize}
        \item Normal : $G_{\text{mic}} = |\gamma_{\text{obj}}| \times G_{\text{oc}} = 40 \times 10 = \mathbf{400}$ (souvent noté $\times 400$).
        \item Confort : $G_{\text{mic}} = 40 \times 11 = \mathbf{440}$.
    \end{itemize}
    \textbf{Interprétation :} En fonctionnement normal, l'œil observe une image à l'infini 400 fois plus grande (en angle) que l'objet vu à l'œil nu à \SI{25}{cm}. En confort, le grossissement est légèrement plus fort (440) mais l'œil fatigue.
\end{enumerate}
\end{exemplebox}

\begin{exoresolu}[Analyse d'un schéma de microscope — Type Probatoire]
\textbf{Énoncé.} On observe la figure ci-dessous représentant un microscope en fonctionnement normal. L'intervalle optique $\Delta = \SI{16}{cm}$. La focale de l'objectif est $f'_{\text{obj}} = \SI{5}{mm}$. La focale de l'oculaire est $f'_{\text{oc}} = \SI{2,5}{cm}$.
\begin{enumerate}
    \item Sur la figure, l'image intermédiaire $A'B'$ est tracée. Justifier sa position par rapport au foyer $F_{\text{oc}}$.
    \item Calculer le grandissement $\gamma_{\text{obj}}$ de l'objectif.
    \item Déterminer le grandissement commercial $G_{\text{mic}}$ du microscope.
    \item L'image finale est-elle droite ou inversée par rapport à l'objet $AB$ ? Justifier.
\end{enumerate}

\begin{tikzpicture}[scale=0.7, >=Stealth]
    % Reproduction simplifiée pour l'exo
    \def\fobj{0.5} \def\foc{2.5} \def\Delta{16}
    \def\Oobj{0} \def\Fpobj{\fobj} \def\Foc{\Fpobj+\Delta} \def\Ooc{\Foc+\foc}
    \draw[popline, thick] (-1,0) -- (\Ooc+2,0);
    \draw[popblue, very thick] (\Oobj,-1.5) -- (\Oobj,1.5) node[above] {$L_{\text{obj}}$};
    \draw[popblue!70] (\Oobj-\fobj,-0.1) -- (\Oobj-\fobj,0.1) node[below] {$F_{\text{obj}}$};
    \draw[popblue!70] (\Fpobj,-0.1) -- (\Fpobj,0.1) node[below] {$F'_{\text{obj}}$};
    \draw[popblue, very thick] (\Ooc,-1.5) -- (\Ooc,1.5) node[above] {$L_{\text{oc}}$};
    \draw[popblue!70] (\Foc,-0.1) -- (\Foc,0.1) node[below] {$F_{\text{oc}}$};
    \draw[popblue!70] (\Ooc+\foc,-0.1) -- (\Ooc+\foc,0.1) node[below] {$F'_{\text{oc}}$};
    \draw[popdark, very thick] (\Oobj-0.6,0) -- (\Oobj-0.6,0.8) node[above] {$B$};
    \draw[popgreen, very thick] (\Foc,0) -- (\Foc,-2.5) node[below] {$B'$};
    \draw[poporange, thick, ->] (\Oobj-0.6,0.8) -- (\Oobj,0.8) -- (\Fpobj,0) -- (\Foc,-2.5);
    \draw[poporange, thick, ->] (\Oobj-0.6,0.8) -- (\Oobj,0) -- (\Foc,-2.5);
    \draw[poporange, thick, ->] (\Foc,-2.5) -- (\Ooc,-2.5) -- (\Ooc+2,-2.5);
    \draw[poporange, thick, ->] (\Foc,-2.5) -- (\Ooc,0) -- (\Ooc+2,0);
\end{tikzpicture}

\tcblower
\textbf{Solution détaillée.}
\begin{enumerate}
    \item \textbf{Position de $A'B'$ :} Le microscope fonctionne en \textbf{fonctionnement normal} (rayons émergents de l'oculaire parallèles entre eux et à l'axe). Pour qu'une lentille convergente (l'oculaire) donne des rayons émergents parallèles, son objet doit être placé \textbf{au foyer objet} $F_{\text{oc}}$. Donc $A'$ est confondu avec $F_{\text{oc}}$. L'image intermédiaire $A'B'$ est dans le plan focal objet de l'oculaire.
    
    \item \textbf{Grandissement de l'objectif :}
    \[
    \gamma_{\text{obj}} \approx -\frac{\Delta}{f'_{\text{obj}}} = -\frac{\SI{160}{mm}}{\SI{5}{mm}} = -32
    \]
    (On a converti $\Delta = \SI{16}{cm} = \SI{160}{mm}$ pour avoir des unités cohérentes).
    
    \item \textbf{Grandissement commercial du microscope (fonctionnement normal) :}
    Grandissement oculaire : $G_{\text{oc}} = \dfrac{\SI{25}{cm}}{f'_{\text{oc}}} = \dfrac{\SI{25}{cm}}{\SI{2,5}{cm}} = 10$.
    Grandissement total : $G_{\text{mic}} = |\gamma_{\text{obj}}| \times G_{\text{oc}} = 32 \times 10 = \mathbf{320}$.
    
    \item \textbf{Sens de l'image finale :}
    \begin{itemize}
        \item Objectif : $\gamma_{\text{obj}} = -32 < 0 \Rightarrow$ Image intermédiaire $A'B'$ \textbf{inversée} par rapport à $AB$.
        \item Oculaire : Fonctionne en loupe sur un objet réel ($A'B'$ entre $F_{\text{oc}}$ et lentille pour le confort, \emph{sur} $F_{\text{oc}}$ pour le normal). Dans les deux cas, le grandissement angulaire $G_{\text{oc}} > 0$. L'image finale $A''B''$ a le \textbf{même sens} que $A'B'$.
    \end{itemize}
    Conclusion : L'image finale $A''B''$ est \textbf{inversée} par rapport à l'objet initial $AB$.
\end{enumerate}
\end{exoresolu}

\begin{methode}[Construction géométrique rigoureuse au Probatoire]
Si l'énoncé demande « Faire la construction géométrique de l'image intermédiaire » ou « Construire l'image finale » :
\begin{enumerate}
    \item \textbf{Repérage :} Tracer l'axe optique, placer $O_{\text{obj}}$, $F_{\text{obj}}$, $F'_{\text{obj}}$, $O_{\text{oc}}$, $F_{\text{oc}}$, $F'_{\text{oc}}$. Respecter $\Delta = \overline{F'_{\text{obj}}F_{\text{oc}}}$.
    \item \textbf{Objet :} Placer $AB$ perpendiculaire à l'axe, $A$ sur l'axe, $B$ légèrement après $F_{\text{obj}}$ ($p_{\text{obj}} > f'_{\text{obj}}$).
    \item \textbf{Objectif (2 rayons suffisent) :}
        \begin{itemize}
            \item Rayon 1 : Partant de $B$, \textbf{parallèle à l'axe} $\rightarrow$ réfracté par $L_{\text{obj}}$ en passant par $F'_{\text{obj}}$.
            \item Rayon 2 : Partant de $B$, passant par le \textbf{centre optique $O_{\text{obj}}$} $\rightarrow$ non dévié.
        \end{itemize}
        Intersection $\rightarrow B'$. $A'$ est sur l'axe (pied de la perpendiculaire).
    \item \textbf{Vérification cruciale :} $A'$ doit tomber \textbf{sur $F_{\text{oc}}$} (fonctionnement normal) ou \textbf{entre $F_{\text{oc}}$ et $O_{\text{oc}}$} (confort). Si ce n'est pas le cas, votre construction de l'objectif est fausse (objet mal placé).
    \item \textbf{Oculaire (si image finale demandée) :} Traiter $A'B'$ comme objet.
        \begin{itemize}
            \item Normal : Rayons sortants \textbf{parallèles} (un passant par $O_{\text{oc}}$, l'autre $\parallel$ axe incident $\rightarrow$ passe par $F'_{\text{oc}}$ puis $\parallel$ axe).
            \item Confort : Construire l'image virtuelle $A''B''$ par intersection des prolongements des rayons émergents divergents.
        \end{itemize}
    \item \textbf{Légender :} Nommer tous les points ($A, B, A', B', F, F', O$), indiquer les sens des flèches (rayons incidents $\rightarrow$ réfractés), noter $\Delta$, $f'_{\text{obj}}$, $f'_{\text{oc}}$.
\end{enumerate}
\end{methode}

\begin{savaistu}[Le standard RMS (Royal Microscopical Society) et le « tube fini »]
La valeur $\Delta = \SI{160}{mm}$ n'est pas un hasard. Elle correspond à une norme historique (RMS, 19e siècle) pour les microscopes à « tube fini » (longueur mécanique $L \approx \SI{160}{mm} + f'_{\text{obj}} + f'_{\text{oc}}$). Cela garantissait l'interchangeabilité des objectifs et oculaires de marques différentes (Zeiss, Leitz, Olympus, Nikon).
Aujourd'hui, les microscopes de recherche utilisent l'optique « à l'infini » (objectif donne image à l'infini, tube lens avant l'oculaire). $\Delta$ n'est plus pertinent, on parle de distance tube (souvent \SI{200}{mm}). Mais au Probatoire, \textbf{le modèle à tube fini avec $\Delta = \SI{160}{mm}$ reste la référence exclusive}. Ne soyez pas déstabilisé si vous voyez d'autres valeurs dans des sujets d'annales récentes : l'énoncé donne toujours $\Delta$.
\end{savaistu}

\coursec{Grossissement commercial et Puissance commerciale}

Dans la section précédente, nous avons construit le schéma de principe du microscope et déterminé la position et la nature des images intermédiaires. Nous savons maintenant que l'objectif forme une image réelle, inversée et grossie $A_1B_1$ de l'objet $AB$, et que l'oculaire agit comme une loupe sur cette image intermédiaire pour fournir une image finale virtuelle $A_2B_2$ à l'infini (fonctionnement normal) ou au point proche (fonctionnement confort). Il nous reste à quantifier ce grossissement global et à introduire la notion de puissance commerciale, langage standard des fabricants (Zeiss, Olympus, Leica, ou les équipements des laboratoires du CHU de Yaoundé ou de l'Université de Douala).

\courssub{Définition du grossissement commercial $G$}

Le grossissement d'un microscope ne se définit pas comme le rapport de tailles linéaires (l'image finale est à l'infini, sa taille linéaire n'a pas de sens), mais comme un \textbf{grossissement angulaire}. C'est le rapport de l'angle sous-tendu par l'image finale à l'œil à l'angle sous-tendu par l'objet vu à l'œil nu dans les meilleures conditions de netteté (au point proche, $D = 25\,\text{cm}$).

\begin{definition}[Grossissement commercial $G$]
Soit un microscope utilisé par un œil emmétrope (sans défaut de vision). On note :
\begin{itemize}
    \item $\alpha$ : l'angle sous-tendu par l'objet $AB$ placé au point proche $D = 25\,\text{cm}$ de l'œil nu (vision de référence).
    \item $\alpha'$ : l'angle sous-tendu par l'image finale $A_2B_2$ à l'œil placé au niveau de l'oculaire.
\end{itemize}
Le \textbf{grossissement commercial} $G$ est défini par :
\[
G = \left| \frac{\alpha'}{\alpha} \right|
\]
$G$ est un nombre sans unité. Il indique combien de fois l'instrument agrandit l'angle de vision par rapport à la vision directe au point proche.
\end{definition}

\courssub{Expression du grossissement commercial}

Rappelons les résultats de la section 2 (schéma de principe) :
\begin{itemize}
    \item L'objectif (focale $f'_{\text{obj}}$, puissance $P_{\text{obj}}$) donne un grossissement transverse (linéaire) $\gamma_{\text{obj}} = \dfrac{\overline{A_1B_1}}{\overline{AB}}$. Comme l'image $A_1B_1$ est réelle et inversée, $\gamma_{\text{obj}} < 0$. Sa valeur absolue est $|\gamma_{\text{obj}}| = \dfrac{\Delta}{f'_{\text{obj}}}$, où $\Delta = \overline{F'_{\text{obj}}F_{\text{oc}}}$ est l'intervalle optique (distance entre le foyer image de l'objectif et le foyer objet de l'oculaire). C'est la \textbf{longueur du tube optique} (souvent normalisée à $\Delta = 160\,\text{mm}$ sur les microscopes standards).
    \item L'oculaire (focale $f'_{\text{oc}}$, puissance $P_{\text{oc}}$) fonctionne comme une loupe. Son grossissement commercial (angulaire) est noté $G_{\text{oc}}$.
\end{itemize}

Le grossissement total est le produit du grossissement linéaire de l'objectif (en valeur absolue) par le grossissement angulaire de l'oculaire :
\[
G = |\gamma_{\text{obj}}| \times G_{\text{oc}}
\]

\paragraph{Cas 1 : Fonctionnement normal (image finale à l'infini, œil au repos).}
L'œil est détendu, l'image intermédiaire $A_1B_1$ est au foyer objet de l'oculaire ($F_{\text{oc}}$).
Le grossissement de l'oculaire vaut $G_{\text{oc}} = \dfrac{D}{f'_{\text{oc}}}$ avec $D = 25\,\text{cm} = 250\,\text{mm}$.
\[
\boxed{G_{\text{normal}} = \frac{\Delta}{f'_{\text{obj}}} \times \frac{25\,\text{cm}}{f'_{\text{oc}}}}
\]

\paragraph{Cas 2 : Fonctionnement confort (image finale au point proche $D = 25\,\text{cm}$).}
L'œil accommode au maximum. L'image intermédiaire $A_1B_1$ est légèrement en avant de $F_{\text{oc}}$.
Le grossissement de l'oculaire vaut $G_{\text{oc}} = \dfrac{D}{f'_{\text{oc}}} + 1$.
\[
\boxed{G_{\text{confort}} = \frac{\Delta}{f'_{\text{obj}}} \times \left( \frac{25\,\text{cm}}{f'_{\text{oc}}} + 1 \right)}
\]

\begin{propriete}[Formules de calcul du grossissement commercial]
Pour un microscope de longueur de tube optique $\Delta$ (en mm), objectif de focale $f'_{\text{obj}}$ (en mm) et oculaire de focale $f'_{\text{oc}}$ (en mm) :
\begin{align*}
\text{Fonctionnement normal :}\quad & G = \frac{\Delta}{f'_{\text{obj}}} \times \frac{250}{f'_{\text{oc}}} \\
\text{Fonctionnement confort :}\quad & G = \frac{\Delta}{f'_{\text{obj}}} \times \left( \frac{250}{f'_{\text{oc}}} + 1 \right)
\end{align*}
\textbf{Attention à l'homogénéité des unités :} $\Delta$ et $f'_{\text{obj}}$ doivent être dans la \textbf{même unité} (mm ou cm). La constante $25\,\text{cm}$ s'exprime en $250\,\text{mm}$ si $\Delta$ est en mm.
\end{propriete}

\begin{exemplebox}[Calcul de $G$ : Microscope standard de lycée]
Un microscope de laboratoire (type équipant les lycées de Douala ou Yaoundé) a les caractéristiques suivantes :
\begin{itemize}
    \item Longueur de tube optique normalisée : $\Delta = 160\,\text{mm}$.
    \item Objectif : $f'_{\text{obj}} = 4,0\,\text{mm}$ (marquage $40\times$).
    \item Oculaire : $f'_{\text{oc}} = 25\,\text{mm}$ (marquage $10\times$).
\end{itemize}
\textbf{1. Fonctionnement normal (image à l'infini) :}
\[
G_{\text{normal}} = \frac{160}{4,0} \times \frac{250}{25} = 40 \times 10 = \mathbf{400}
\]
\textbf{2. Fonctionnement confort (image à 25 cm) :}
\[
G_{\text{confort}} = \frac{160}{4,0} \times \left( \frac{250}{25} + 1 \right) = 40 \times (10 + 1) = \mathbf{440}
\]
L'instrument est donc un microscope $400\times$ (normal) / $440\times$ (confort).
\end{exemplebox}

\courssub{Puissance commerciale $P$ (en dioptries)}

La puissance commerciale est une grandeur additive qui caractérise la convergence globale de l'instrument. Elle est définie par convention comme la somme des puissances des deux lentilles (objectif + oculaire), exprimée en dioptries ($\delta = \text{m}^{-1}$).

\begin{definition}[Puissance commerciale $P$]
Soit $P_{\text{obj}} = \dfrac{1}{f'_{\text{obj}}}$ et $P_{\text{oc}} = \dfrac{1}{f'_{\text{oc}}}$ les puissances de l'objectif et de l'oculaire (focales en \textbf{mètres}).
La \textbf{puissance commerciale} du microscope est :
\[
\boxed{P = P_{\text{obj}} + P_{\text{oc}} \quad (\text{en } \delta \text{ ou m}^{-1})}
\]
\end{definition}

\begin{attention}[Puissance vs Grossissement : Ne pas confondre !]
\begin{itemize}
    \item Les \textbf{puissances s'additionnent} : $P = P_{\text{obj}} + P_{\text{oc}}$. C'est une définition conventionnelle commode pour classer les instruments.
    \item Les \textbf{grossissements se multiplient} : $G = |\gamma_{\text{obj}}| \times G_{\text{oc}}$.
    \item \textbf{Erreur fatale au Probatoire :} Écrire $G = P$ ou $G = k \times P$ sans justification, ou additionner les grossissements ($G = G_{\text{obj}} + G_{\text{oc}}$).
\end{itemize}
\end{attention}

\paragraph{Lien approximatif entre $G$ et $P$ (fonctionnement normal).}
En remplaçant $f'_{\text{obj}} = 1/P_{\text{obj}}$ et $f'_{\text{oc}} = 1/P_{\text{oc}}$ (focales en mètres) dans la formule de $G_{\text{normal}}$ :
\[
G_{\text{normal}} = \Delta \times P_{\text{obj}} \times \frac{0,25}{1} \times P_{\text{oc}} = 0,25 \times \Delta \times P_{\text{obj}} \times P_{\text{oc}}
\]
avec $\Delta$ en \textbf{mètres}.
Si on utilise la normale $\Delta = 160\,\text{mm} = 0,16\,\text{m}$ :
\[
G_{\text{normal}} \approx 0,25 \times 0,16 \times P_{\text{obj}} \times P_{\text{oc}} = 0,04 \times P_{\text{obj}} \times P_{\text{oc}}
\]
On admet souvent la relation approchée (valable pour $\Delta = 16\,\text{cm}$) :
\[
\boxed{G \approx \frac{P_{\text{obj}} \times P_{\text{oc}}}{25} \quad \text{ou} \quad G \approx \frac{P}{4} \text{ (très approximatif, si } P_{\text{obj}} \gg P_{\text{oc}})}
\]
\textbf{Prudence :} Cette relation n'est qu'une estimation rapide. Au Probatoire, on exige le calcul rigoureux via les focales.

\begin{exemplebox}[Calcul de la puissance commerciale]
Reprenons le microscope de l'exemple précédent ($f'_{\text{obj}} = 4\,\text{mm}$, $f'_{\text{oc}} = 25\,\text{mm}$).
\begin{enumerate}
    \item Conversion en mètres : $f'_{\text{obj}} = 4 \times 10^{-3}\,\text{m}$ ; $f'_{\text{oc}} = 25 \times 10^{-3}\,\text{m}$.
    \item Puissances :
    \[
    P_{\text{obj}} = \frac{1}{4 \times 10^{-3}} = 250\,\delta \quad ; \quad P_{\text{oc}} = \frac{1}{25 \times 10^{-3}} = 40\,\delta
    \]
    \item Puissance commerciale :
    \[
    P = 250 + 40 = \mathbf{290\,\delta}
    \]
    \item Vérification du lien $G \approx 0,04 P_{\text{obj}} P_{\text{oc}}$ :
    $0,04 \times 250 \times 40 = 400 = G_{\text{normal}}$. La relation est exacte ici car $\Delta = 0,16\,\text{m}$ exactement.
\end{enumerate}
\end{exemplebox}

\courssub{Marquage normalisé des objectifs et oculaires (Norme ISO 8038)}

Sur le corps de l'objectif et de l'oculaire, on ne trouve pas les focales, mais des codes normalisés. Savoir les déchiffrer est une compétence clé.

\begin{aretenir}[Lecture du marquage normalisé]
\begin{center}
\begin{tabularx}{\linewidth}{|l|X|X|}\hline
\textbf{Élément} & \textbf{Marquage typique} & \textbf{Signification et déduction} \\ \hline
\textbf{Objectif} & \texttt{40 / 0,65} / \texttt{160 / 0,17} & \begin{itemize}
    \item \texttt{40} : Grossissement linéaire de l'objectif $|\gamma_{\text{obj}}| = 40\times$ (pour $\Delta = 160\,\text{mm}$).
    \item \texttt{0,65} : Ouverture numérique $NA = n \sin i_{\max}$ (détermine la résolution, voir Section 5).
    \item \texttt{160} : Longueur de tube mécanique de référence (mm).
    \item \texttt{0,17} : Épaisseur de la lame de couvre-objet prévue (mm).
\end{itemize}
\textbf{Focale déduite :} $f'_{\text{obj}} = \dfrac{\Delta}{|\gamma_{\text{obj}}|} = \dfrac{160\,\text{mm}}{40} = 4,0\,\text{mm}$. \\ \hline
\textbf{Oculaire} & \texttt{10$\times$} ou \texttt{WF 10$\times$ / 18} & \begin{itemize}
    \item \texttt{10$\times$} : Grossissement commercial de l'oculaire $G_{\text{oc}} = 10$ (fonctionnement normal).
    \item \texttt{WF} : \emph{Wide Field} (champ large).
    \item \texttt{18} : Diamètre du champ de vision (mm) au plan de l'image intermédiaire.
\end{itemize}
\textbf{Focale déduite :} $f'_{\text{oc}} = \dfrac{250\,\text{mm}}{G_{\text{oc}}} = \dfrac{250}{10} = 25\,\text{mm}$. \\ \hline
\textbf{Microscope complet} & \textbf{Non marqué globalement.} & Le grossissement total $G$ et la puissance $P$ se \textbf{calculent} à partir des composants. \\ \hline
\end{tabularx}
\end{center}
\end{aretenir}

\begin{methode}[Déchiffrer un marquage et calculer $G$ et $P$]
\begin{enumerate}
    \item \textbf{Objectif :} Lire le grossissement gravé (ex: $40\times$). C'est $|\gamma_{\text{obj}}|$. En déduire $f'_{\text{obj}} = \Delta / |\gamma_{\text{obj}}|$ (avec $\Delta = 160\,\text{mm}$ sauf indication contraire).
    \item \textbf{Oculaire :} Lire le grossissement gravé (ex: $10\times$). C'est $G_{\text{oc}}$. En déduire $f'_{\text{oc}} = 250\,\text{mm} / G_{\text{oc}}$.
    \item \textbf{Grossissement total :} $G = |\gamma_{\text{obj}}| \times G_{\text{oc}}$ (normal) ou $|\gamma_{\text{obj}}| \times (G_{\text{oc}}+1)$ (confort).
    \item \textbf{Puissance :} Convertir $f'_{\text{obj}}$ et $f'_{\text{oc}}$ en mètres. Calculer $P_{\text{obj}} = 1/f'_{\text{obj}}$, $P_{\text{oc}} = 1/f'_{\text{oc}}$. Sommier : $P = P_{\text{obj}} + P_{\text{oc}}$.
\end{enumerate}
\end{methode}

\courssub{Problème inverse : Choisir les composants pour un grossissement donné}

C'est un classique du Probatoire : « On dispose d'objectifs $4\times, 10\times, 40\times, 100\times$ et d'oculaires $5\times, 10\times, 15\times$. Choisir la combinaison pour obtenir $G \approx 600$ en fonctionnement normal. »

\begin{exoresolu}[Choix de la configuration pour $G = 600$]
\textbf{Données :} $\Delta = 160\,\text{mm}$. Objectifs disponibles : $4\times, 10\times, 40\times, 100\times$. Oculaires disponibles : $5\times, 10\times, 15\times$. Fonctionnement normal.
\textbf{Recherche :} Couple (Objectif, Oculaire) tel que $G = |\gamma_{\text{obj}}| \times G_{\text{oc}} \approx 600$.

\textbf{Raisonnement :}
On teste les produits $|\gamma_{\text{obj}}| \times G_{\text{oc}}$ :
\begin{itemize}
    \item $4 \times 5 = 20$ ; $4 \times 10 = 40$ ; $4 \times 15 = 60$ (trop faibles)
    \item $10 \times 5 = 50$ ; $10 \times 10 = 100$ ; $10 \times 15 = 150$ (trop faibles)
    \item $40 \times 5 = 200$ ; $40 \times 10 = 400$ ; $40 \times 15 = \mathbf{600}$ ✓
    \item $100 \times 5 = 500$ ; $100 \times 10 = 1000$ ; $100 \times 15 = 1500$ (trop forts)
\end{itemize}

\tcblower

\textbf{Solution :} Choisir l'\textbf{objectif $40\times$} et l'\textbf{oculaire $15\times$}.
\textbf{Vérification des focales et puissance :}
\begin{itemize}
    \item $f'_{\text{obj}} = 160/40 = 4,0\,\text{mm}$ ; $P_{\text{obj}} = 1/0,004 = 250\,\delta$.
    \item $f'_{\text{oc}} = 250/15 \approx 16,7\,\text{mm}$ ; $P_{\text{oc}} = 1/0,0167 \approx 60\,\delta$.
    \item $P = 250 + 60 = 310\,\delta$.
    \item $G_{\text{normal}} = 40 \times 15 = 600$. $G_{\text{confort}} = 40 \times 16 = 640$.
\end{itemize}
\end{exoresolu}

\begin{attention}[Pièges classiques au Probatoire]
\begin{enumerate}
    \item \textbf{Unités mélangées :} Utiliser $\Delta = 16\,\text{cm}$ et $f'_{\text{obj}} = 4\,\text{mm}$ sans conversion. \textbf{Toujours convertir $\Delta$ et $f'_{\text{obj}}$ en mm (ou cm) avant de diviser.}
    \item \textbf{Addition des grossissements :} Écrire $G = 40 + 10 = 50$ au lieu de $40 \times 10 = 400$. Zéro pointé.
    \item \textbf{Confusion $P$ et $G$ :} Calculer $P = 290\,\delta$ et répondre « Le grossissement est $290\times$ ». Faux. $G=400$, $P=290\,\delta$.
    \item \textbf{Oubli du $+1$ en confort :} La formule $G = \frac{\Delta}{f'_{\text{obj}}} \times \frac{250}{f'_{\text{oc}}}$ est \textbf{uniquement} pour le fonctionnement normal (image à l'infini). Si l'énoncé précise « œil accommodé » ou « image au point proche », il faut ajouter $+1$ au facteur oculaire.
    \item \textbf{Marquage objectif :} Le nombre gravé sur l'objectif (ex: 40) est $|\gamma_{\text{obj}}|$, pas $f'_{\text{obj}}$ en mm (sauf coïncidence pour le 40× où $f'=4\,\text{mm}$). Pour un $100\times$, $f'_{\text{obj}} = 1,6\,\text{mm}$, pas $100\,\text{mm}$ !
\end{enumerate}
\end{attention}

\courssub{Synthèse visuelle : Tableau récapitulatif des calculs}

\begin{popfigure}
\begin{tikzpicture}[
    node distance=2mm and 4mm,
    box/.style={draw=popdark, thick, rounded corners=3pt, fill=popcream, minimum height=9mm, text width=2.8cm, align=center, font=\small},
    boxh/.style={box, fill=popblueL, font=\small\bfseries, text width=2.8cm},
    boxg/.style={box, fill=popgreenL, font=\small\bfseries, text width=2.8cm},
    boxo/.style={box, fill=poporangeL, font=\small\bfseries, text width=2.8cm},
    arrow/.style={-Stealth, thick, popblue},
    arrowr/.style={-Stealth, thick, poporange},
]
% En-têtes
\node[boxh] (h1) {Élément};
\node[boxh, right=of h1] (h2) {Objectif};
\node[boxh, right=of h2] (h3) {Oculaire};
\node[boxh, right=of h3] (h4) {Microscope Complet};

% Ligne 1 : Marquage
\node[boxg, below=of h1] (l1) {Marquage normalisé};
\node[box, below=of h2, align=center] (o1){\texttt{40 / 0,65}\\\texttt{160 / 0,17}};
\node[box, below=of h3] (oc1) {\texttt{WF 10$\times$ / 18}};
\node[box, below=of h4] (m1) {--};

% Ligne 2 : Focale
\node[boxg, below=of l1] (l2) {Focale $f'$};
\node[box, below=of o1] (o2) {$f'_{\text{obj}} = \frac{\Delta}{40} = \mathbf{4,0\,mm}$};
\node[box, below=of oc1] (oc2) {$f'_{\text{oc}} = \frac{250}{10} = \mathbf{25\,mm}$};
\node[box, below=of m1] (m2) {--};

% Ligne 3 : Puissance (dioptries)
\node[boxg, below=of l2] (l3) {Puissance $P$ ($\delta$)};
\node[box, below=of o2] (o3) {$P_{\text{obj}} = \frac{1}{0,004} = \mathbf{250\,\delta}$};
\node[box, below=of oc2] (oc3) {$P_{\text{oc}} = \frac{1}{0,025} = \mathbf{40\,\delta}$};
\node[box, below=of m2] (m3) {$P = P_{\text{obj}} + P_{\text{oc}} = \mathbf{290\,\delta}$};

% Ligne 4 : Grossissement partiel
\node[boxg, below=of l3] (l4) {Grossissement partiel};
\node[box, below=of o3] (o4) {$|\gamma_{\text{obj}}| = \mathbf{40\times}$};
\node[box, below=of oc3] (oc4) {$G_{\text{oc}} = \mathbf{10\times}$ (normal)};
\node[box, below=of m3] (m4) {--};

% Ligne 5 : Grossissement total
\node[boxg, below=of l4] (l5) {Grossissement total $G$};
\node[box, below=of o4] (o5) {--};
\node[box, below=of oc4] (oc5) {--};
\node[box, below=of m4, align=center] (m5){
    Normal : $40 \times 10 = \mathbf{400}$\\
    Confort : $40 \times 11 = \mathbf{440}$
};

% Flèches de calcul
\draw[arrowr] (o1.south) -- (o2.north) node[midway, right, font=\tiny, poporange] {$\Delta=160$mm};
\draw[arrowr] (oc1.south) -- (oc2.north) node[midway, right, font=\tiny, poporange] {$D=250$mm};
\draw[arrowr] (o2.south) -- (o3.north) node[midway, right, font=\tiny, poporange] {$f'$ en m};
\draw[arrowr] (oc2.south) -- (oc3.north) node[midway, right, font=\tiny, poporange] {$f'$ en m};
\draw[arrowr] (o3.south) -- (o4.north) node[midway, right, font=\tiny, poporange] {Marquage direct};
\draw[arrowr] (oc3.south) -- (oc4.north) node[midway, right, font=\tiny, poporange] {Marquage direct};
\draw[arrow, bend left=15] (o3.east) -- node[above, font=\tiny, popblue] {$+$} (m3.west);
\draw[arrow, bend left=15] (oc3.east) -- node[above, font=\tiny, popblue] {$+$} (m3.west);
\draw[arrow, bend left=15] (o4.east) -- node[above, font=\tiny, popblue] {$\times$} (m5.west);
\draw[arrow, bend left=15] (oc4.east) -- node[above, font=\tiny, popblue] {$\times$} (m5.west);
\end{tikzpicture}
\legende{Tableau récapitulatif : du marquage normalisé aux grandeurs caractéristiques du microscope (exemple objectif $40\times$, oculaire $10\times$, $\Delta=160\,\text{mm}$). Les flèches orange indiquent les déductions depuis le marquage ; les flèches bleues les combinaisons finales (addition pour $P$, multiplication pour $G$).}
\end{popfigure}

\begin{savaistu}[Le microscope au Cameroun : de la recherche à l'enseignement]
Les laboratoires de l'\textbf{IRAD} (Institut de Recherche Agricole pour le Développement) à Nkolbisson ou de l'\textbf{IMPM} (Institut de Recherche Médicale et d'Études des Plantes Médicinales) à Yaoundé utilisent des microscopes à fluorescence ou confocaux dont les objectifs ont des corrections complexes (plan-apochromatiques) et des ouvertures numériques $NA > 1,4$ (immersion huile). Leur grossissement utile peut dépasser $1000\times$ (limite de diffraction). En revanche, dans nos lycées (Lycée Leclerc, Lycée Technique de Douala, etc.), les microscopes « à miroir » ou à éclairage intégré (type XSZ-107BN) équipés d'objectifs achromatiques $4\times, 10\times, 40\times$ et d'oculaires $10\times, 15\times$ couvrent parfaitement le programme : observation de cellules végétales (oignon, mousse), coupes histologiques, hématies, bactéries (coloration de Gram). La maîtrise du calcul de $G$ et $P$ permet au technicien de choisir la bonne combinaison pour voir un noyau de cellule ($\sim 10\,\mu\text{m}$) ou une bactérie ($\sim 1\,\mu\text{m}$).
\end{savaistu}

\begin{exercice}[Entraînement Probatoire : Microscope à immersion]
Un microscope possède un tube optique de longueur $\Delta = 160\,\text{mm}$. On utilise un objectif à immersion huile marqué \texttt{100 / 1,25 / 160 / 0,17} et un oculaire marqué \texttt{15$\times$}.
\begin{enumerate}
    \item Déterminer les focales $f'_{\text{obj}}$ et $f'_{\text{oc}}$ en mm.
    \item Calculer la puissance commerciale $P$ de l'instrument (en $\delta$).
    \item Calculer le grossissement commercial $G$ en fonctionnement normal et en fonctionnement confort.
    \item L'ouverture numérique de l'objectif est $NA = 1,25$. Sachant que l'indice de l'huile d'immersion est $n = 1,515$, calculer l'angle maximal d'ouverture $i_{\max}$ de l'objectif. (Rappel : $NA = n \sin i_{\max}$).
\end{enumerate}
\end{exercice}

\begin{corrige}[Exercice : Microscope à immersion]
\begin{enumerate}
    \item \textbf{Focales :}
    \begin{itemize}
        \item Objectif : Marquage $100\times \Rightarrow |\gamma_{\text{obj}}| = 100$. $f'_{\text{obj}} = \frac{\Delta}{|\gamma_{\text{obj}}|} = \frac{160}{100} = \mathbf{1,6\,mm}$.
        \item Oculaire : Marquage $15\times \Rightarrow G_{\text{oc}} = 15$. $f'_{\text{oc}} = \frac{250}{15} \approx \mathbf{16,7\,mm}$.
    \end{itemize}
    \item \textbf{Puissance commerciale :}
    Conversion en mètres : $f'_{\text{obj}} = 1,6 \times 10^{-3}\,\text{m}$ ; $f'_{\text{oc}} = 16,7 \times 10^{-3}\,\text{m}$.
    $P_{\text{obj}} = \frac{1}{1,6 \times 10^{-3}} = 625\,\delta$.
    $P_{\text{oc}} = \frac{1}{16,7 \times 10^{-3}} \approx 60\,\delta$.
    $P = 625 + 60 = \mathbf{685\,\delta}$.
    \item \textbf{Grossissement commercial :}
    \begin{itemize}
        \item Normal : $G = |\gamma_{\text{obj}}| \times G_{\text{oc}} = 100 \times 15 = \mathbf{1500}$.
        \item Confort : $G = 100 \times (15 + 1) = \mathbf{1600}$.
    \end{itemize}
    \item \textbf{Angle d'ouverture maximal :}
    $NA = n \sin i_{\max} \Rightarrow \sin i_{\max} = \frac{NA}{n} = \frac{1,25}{1,515} \approx 0,825$.
    $i_{\max} = \arcsin(0,825) \approx \mathbf{55,6^\circ}$.
    Cet angle élevé (proche de $90^\circ$) justifie l'usage de l'huile ($n \approx 1,5$) pour capter un maximum de lumière diffractée et atteindre la limite de résolution théorique ($\delta x \approx \frac{\lambda}{2 NA}$).
\end{enumerate}
\end{corrige}

\coursec{Mise au point : théorie et pratique}

Dans la section précédente, nous avons appris à quantifier les performances d'un microscope : sa puissance commerciale $P_c$ et son grossissement commercial $G_c$. Mais ces grandeurs ne prennent tout leur sens que si l'observateur voit une image \emph{nette}. Or, quiconque a déjà approché son œil d'un microscope au laboratoire du lycée le sait : sans une manipulation préalable des deux vis latérales, on ne voit qu'une tache lumineuse floue. Cette manipulation, c'est la \cle{mise au point}. C'est elle que nous allons maintenant étudier, d'abord sur le plan théorique (où doit se former l'image intermédiaire ?), puis sur le plan pratique (comment y parvenir sans casser la lame ?).

\courssub{Le principe de la mise au point}

\textbf{Une observation pour commencer.} Prends une loupe et regarde un texte à travers elle : tant que la distance loupe–texte n'est pas la bonne, les lettres sont floues ou invisibles. Le microscope fonctionne de la même manière, mais en deux étages : l'objectif fabrique une image intermédiaire $A_1B_1$, et l'oculaire joue le rôle de la loupe pour observer cette image intermédiaire. Pour que l'image finale soit nette et confortable, il faut donc que l'image intermédiaire tombe \emph{au bon endroit} par rapport à l'oculaire.

\begin{definition}[Mise au point d'un microscope]
La \cle{mise au point} d'un microscope consiste à régler la distance entre l'objet observé et l'objectif, de façon que l'\cle{image intermédiaire} $A_1B_1$ donnée par l'objectif se forme dans le \cle{plan focal objet de l'oculaire} (réglage dit \emph{normal}, image finale rejetée à l'infini) ou dans son voisinage immédiat (réglage de confort, image finale à distance finie adaptée à l'œil de l'observateur).
\end{definition}

\begin{aretenir}[Condition fondamentale]
L'image intermédiaire $A_1B_1$ doit coïncider avec le plan focal objet de l'oculaire, c'est-à-dire :
\[ A_1 \equiv F_{2} \quad \text{(foyer objet de l'oculaire).} \]
Alors l'oculaire renvoie l'image définitive $A'B'$ à l'\emph{infini} : l'œil normal observe sans accommoder, sans fatigue. C'est la condition de \cle{vision à l'infini}.
\end{aretenir}

\courssub{Théorie : où doit se trouver l'objet ?}

Notons, pour l'objectif (lentille $L_1$ de centre $O_1$ et de distance focale image $f'_{obj}$) :
\begin{itemize}
\item $p_{obj} = \overline{O_1A}$ la distance algébrique objet–objectif (négative, l'objet étant réel) ;
\item $p'_{obj} = \overline{O_1A_1}$ la distance algébrique objectif–image intermédiaire (positive).
\end{itemize}

\begin{propriete}[Condition de netteté (formule de conjugaison de Descartes)]
La mise au point est réalisée lorsque la position de l'objet vérifie la relation de conjugaison de l'objectif :
\[ \frac{1}{p'_{obj}} - \frac{1}{p_{obj}} = \frac{1}{f'_{obj}}. \]
Pour le réglage normal (image intermédiaire sur le foyer objet de l'oculaire), en notant $\Delta = \overline{F'_{1}F_{2}}$ l'\emph{intervalle optique} (distance entre le foyer image de l'objectif et le foyer objet de l'oculaire), on a $p'_{obj} = f'_{obj} + \Delta$, d'où :
\[ p_{obj} = -\,\frac{f'_{obj}\,(f'_{obj} + \Delta)}{\Delta} \approx -\,f'_{obj}\left(1 + \frac{f'_{obj}}{\Delta}\right). \]
Comme $f'_{obj} \ll \Delta$ (par exemple $f'_{obj} = \SI{4}{mm}$ et $\Delta = \SI{160}{mm}$), l'objet se trouve \emph{très légèrement en avant du foyer objet de l'objectif} : $\lvert p_{obj} \rvert \approx \num{1,025}\,f'_{obj}$ dans cet exemple.
\end{propriete}

\textbf{Conséquence pratique capitale.} La position de l'objet qui donne une image nette est \emph{unique} et extrêmement précise. Un calcul rapide le montre : si l'on déplace l'objet de $\delta = \SI{10}{\micro m}$, l'image intermédiaire se déplace de $\delta' \approx \left(\dfrac{\Delta}{f'_{obj}}\right)^2 \delta = \left(\dfrac{160}{4}\right)^2 \times \SI{10}{\micro m} = \SI{16}{mm}$ ! Le facteur d'amplification des déplacements est le carré du grandissement de l'objectif ($G_{obj} \approx -\Delta/f'_{obj}$). C'est pourquoi la mise au point exige des vis de très grande précision.

\begin{exemplebox}[Position de l'objet pour un objectif 40$\times$]
Un objectif de distance focale $f'_{obj} = \SI{4,0}{mm}$ est associé à un intervalle optique $\Delta = \SI{160}{mm}$. Pour le réglage normal :
\[ p_{obj} = -\frac{\num{4,0}\times(\num{4,0} + 160)}{160} = -\frac{\num{4,0}\times 164}{160} = -\SI{4,1}{mm}. \]
L'objet doit donc être placé à seulement $\SI{0,1}{mm}$ en avant du foyer objet de l'objectif, situé lui-même à $\SI{4,0}{mm}$ de la lentille. La marge de manœuvre est de l'ordre du micromètre : d'où la nécessité de la vis micrométrique.
\end{exemplebox}

\courssub{L'action mécanique : les vis de mise au point}

Sur un microscope réel, on ne déplace pas l'objet à la main : ce sont deux vis coaxiales, situées de part et d'autre du pied du microscope, qui déplacent la \cle{platine} (ou, sur certains modèles, le tube porte-objectif) le long de l'axe optique. Cela modifie la distance $p_{obj}$ entre l'objet et l'objectif.

\begin{definition}[Les deux vis de mise au point]
\begin{itemize}
\item La vis \cle{macrométrique} (grosse vis) : déplacement \emph{rapide} de la platine, de l'ordre de quelques millimètres par tour. Elle sert à amener grossièrement l'image intermédiaire dans le voisinage du plan focal objet de l'oculaire.
\item La vis \cle{micrométrique} (vis fine) : déplacement \emph{lent et précis}, de l'ordre de quelques micromètres par fraction de tour. Elle sert à l'ajustement final : netteté parfaite, c'est-à-dire $A_1$ exactement sur $F_2$ (ou légèrement décalé pour le confort de l'œil).
\end{itemize}
\end{definition}

\textbf{Sens de rotation.} Sur la plupart des microscopes, tourner les vis dans le sens horaire (vu de la droite) fait monter la platine, c'est-à-dire \emph{rapproche} l'objectif de la lame (diminue $\lvert p_{obj} \rvert$). Mais cette convention varie d'un modèle à l'autre : il faut toujours vérifier sur l'instrument, grâce aux flèches gravées près des boutons, avant toute manipulation.

\begin{popfigure}
\begin{tikzpicture}[scale=0.85, >=Stealth]
% --- Styles ---
\tikzset{
  lens/.style={draw=popblue, very thick, fill=popblue!10},
  focal/.style={fill=popgreen, circle, inner sep=1.5pt},
  objpt/.style={fill=popdark, circle, inner sep=1.2pt},
  imgpt/.style={fill=popink, circle, inner sep=1.2pt},
  ray/.style={popink, dashed, thick},
  rayconstr/.style={poporange, thick},
  axis/.style={popline, thick},
  label/.style={font=\small, text=popdark},
  arr/.style={->, thick, popdark}
}

% ============================================================
% CASE 1 : Object TOO FAR (|p| large) -> Image BEFORE F2
% ============================================================
\begin{scope}[xshift=0cm]
% Axis
\draw[axis] (-1.5,0) -- (9,0) node[right, label]{$x$};
% Lenses
\draw[lens] (0,-1) -- (0,1) node[above, popblue, font=\small]{Obj};
\draw[lens] (7,-1.2) -- (7,1.2) node[above, poporange, font=\small]{Oc};
% Focal points
\node[focal, label=below:$F'_1$] at (1,0) {};
\node[focal, label=below:$F_2$] at (5,0) {}; % F2 at x=5 (Delta = 4, f_obj=1)
% Object
\node[objpt, label=below:$A$] (A1) at (-1.2,0) {};
% Image (calculated: p=-1.2, f=1 -> 1/p' = 1 - 1/1.2 = 0.166 -> p'=6. Wait. Need p' < 5 (F2 pos).
% Let f_obj = 1. F2 at 5 => Delta = 4. p'_target = 5.
% Case 1: Object too far -> p < -1.14 (ideal). Say p = -1.5.
% 1/p' = 1 - 1/1.5 = 0.333 -> p' = 3. Image at x=3 (Before F2 at 5).
\node[imgpt, label=below:$A_1$] (A1img) at (3,0) {};
% Rays construction (Object height 0.6)
\coordinate (A1t) at (-1.2, 0.6);
\coordinate (O1) at (0,0);
\coordinate (F1p) at (1,0);
% Ray 1: Parallel -> F'1
\draw[rayconstr] (A1t) -- (0, 0.6) -- (F1p);
% Ray 2: Through Center
\draw[rayconstr] (A1t) -- (O1) -- (A1img);
% Image height: gamma = -p'/p = -3/-1.5 = 2. h' = -1.2.
\coordinate (A1img_t) at (3, -1.2);
\draw[ray] (A1img) -- (A1img_t);
\draw[ray] (0,0.6) -- (A1img_t);
% Object arrow
\draw[arr] (A1) -- (A1t);
% Image arrow
\draw[arr, popink] (A1img) -- (A1img_t);
% Eyepiece rays (diverging from A1 towards F2? No, A1 before F2 -> rays converge after A1 then diverge? 
% Object for oc is A1 (real). A1 is before F2. So for Oc, object is between lens and F2? No, Oc at 7. F2 at 5. A1 at 3.
% Object for Oc is at 3. F2 at 5. Object is BEFORE F2 (left of F2). Real object for Oc.
% Rays from A1_t go through Oc. Since object is before F2, image is virtual, upright, magnified (loupe).
% Draw rays emerging from Oc.
\draw[popgreen, dashed, thick] (A1img_t) -- (7, -0.5) -- (8.5, -1.2);
\draw[popgreen, dashed, thick] (A1img_t) -- (7, 0.2) -- (8.5, 0.8);
\node[font=\small\bfseries, popdark] at (3.5,-2) {1) Objet \emph{trop loin} : $A_1$ avant $F_2$ (flou)};
\end{scope}

% ============================================================
% CASE 2 : Object TOO CLOSE (|p| small) -> Image AFTER F2
% ============================================================
\begin{scope}[xshift=10.5cm]
\draw[axis] (-1.5,0) -- (9,0) node[right, label]{$x$};
\draw[lens] (0,-1) -- (0,1) node[above, popblue, font=\small]{Obj};
\draw[lens] (7,-1.2) -- (7,1.2) node[above, poporange, font=\small]{Oc};
\node[focal, label=below:$F'_1$] at (1,0) {};
\node[focal, label=below:$F_2$] at (5,0) {};
% Object: p = -1.05 (closer to lens than ideal -1.14)
\node[objpt, label=below:$A$] (A2) at (-1.05,0) {};
% Image: 1/p' = 1 - 1/1.05 = 0.0476 -> p' = 21. Way off scale.
% Let's adjust scale or values for visual.
% Let's keep F2 at 5. Target p'=5.
% Case 2: Object too close -> p = -1.05. p' = 21. Image at 21.
% We can't show 21. We show it "way past F2".
% Let's change params for drawing: f_obj=1.5. F2 at 6. Delta=4.5. Target p'=6. Ideal p = -1.5*6/4.5 = -2.
% Case 1 (Far): p=-2.5. 1/p' = 1/1.5 - 1/2.5 = 0.666 - 0.4 = 0.266 -> p'=3.75. (Before F2=6).
% Case 2 (Close): p=-1.8. 1/p' = 0.666 - 0.555 = 0.111 -> p'=9. (After F2=6).
% This fits better in 0..9 range.
% Redraw Case 1 with these params? Too much work. I'll stick to qualitative correct arrows.
% For Case 2 here: Image is to the right of F2.
\node[imgpt, label=below:$A_1$] (A2img) at (6.5,0) {}; % Just past F2 (at 5)
\coordinate (A2t) at (-1.05, 0.5);
\coordinate (O1b) at (0,0);
\coordinate (F1pb) at (1,0);
% Ray 1 Parallel
\draw[rayconstr] (A2t) -- (0, 0.5) -- (F1pb);
% Ray 2 Center
\draw[rayconstr] (A2t) -- (O1b) -- (A2img);
% Gamma = -6.5/-1.05 = 6.19. h' = -3.1.
\coordinate (A2img_t) at (6.5, -3.1);
\draw[ray] (A2img) -- (A2img_t);
\draw[ray] (0,0.5) -- (A2img_t);
\draw[arr] (A2) -- (A2t);
\draw[arr, popink] (A2img) -- (A2img_t);
% Eyepiece rays: Object A1 at 6.5. F2 at 5. Object is AFTER F2 (right of F2).
% For Oc (at 7), object is virtual (on right side)? No. Light goes left to right.
% Obj at 0. Image A1 at 6.5. Light travels to right.
% Oc at 7. Light hits Oc. Object for Oc is A1 at 6.5 (left of Oc). Real object.
% F2 (Oc focal) at 5. Object (6.5) is BETWEEN F2 and Oc.
% Loupe regime: Virtual image at infinity? No, object inside focal length -> virtual image on same side (left), upright.
% Rays emerge divergent.
\draw[popgreen, dashed, thick] (A2img_t) -- (7, -1.5) -- (8.5, -2.5);
\draw[popgreen, dashed, thick] (A2img_t) -- (7, 0.5) -- (8.5, 1.5);
\node[font=\small\bfseries, popdark] at (3.5,-2) {2) Objet \emph{trop près} : $A_1$ après $F_2$ (flou)};
\end{scope}

% ============================================================
% CASE 3 : PERFECT FOCUS
% ============================================================
\begin{scope}[xshift=21cm]
\draw[axis] (-1.5,0) -- (9,0) node[right, label]{$x$};
\draw[lens] (0,-1) -- (0,1) node[above, popblue, font=\small]{Obj};
\draw[lens] (7,-1.2) -- (7,1.2) node[above, poporange, font=\small]{Oc};
\node[focal, label=below:$F'_1$] at (1,0) {};
\node[focal, label=below:$F_2 \equiv A_1$] (F2c) at (5,0) {};
% Ideal object: p = -1.25 (for f=1, Delta=4 -> p'=5). 1/5 - 1/p = 1 -> -1/p = 0.8 -> p=-1.25.
\node[objpt, label=below:$A$] (A3) at (-1.25,0) {};
\node[imgpt] (A3img) at (5,0) {}; % On F2
\coordinate (A3t) at (-1.25, 0.5);
\coordinate (O1c) at (0,0);
\coordinate (F1pc) at (1,0);
% Ray 1 Parallel
\draw[rayconstr] (A3t) -- (0, 0.5) -- (F1pc);
% Ray 2 Center
\draw[rayconstr] (A3t) -- (O1c) -- (A3img);
% Gamma = -5/-1.25 = 4. h' = -2.
\coordinate (A3img_t) at (5, -2);
\draw[ray] (A3img) -- (A3img_t);
\draw[ray] (0,0.5) -- (A3img_t);
\draw[arr] (A3) -- (A3t);
\draw[arr, popink] (A3img) -- (A3img_t);
% Eyepiece: Object at F2. Rays emerge PARALLEL.
\draw[popgreen, thick, ->] (A3img_t) -- (7, -2) -- (9, -2);
\draw[popgreen, thick, ->] (A3img_t) -- (7, 0) -- (9, 0); % Chief ray through center
\draw[popgreen, thick, ->] (A3img_t) -- (7, -1) -- (9, -1);
\node[font=\small\bfseries, popgreen] at (3.5,-2) {3) Mise au point parfaite : $A_1$ sur $F_2$, image à l'infini};
\end{scope}
\end{tikzpicture}
\legende{Principe dynamique de la mise au point (construction géométrique rigoureuse). L'objectif (foyer image $F'_1$) forme l'image intermédiaire $A_1$. L'oculaire a son foyer objet en $F_2$. \textbf{1)} Objet trop loin de l'objectif ($|p|$ grand) $\rightarrow$ $A_1$ se forme \emph{avant} $F_2$. \textbf{2)} Objet trop près ($|p|$ petit) $\rightarrow$ $A_1$ se forme \emph{après} $F_2$. \textbf{3)} Mise au point normale : $A_1 \equiv F_2$, l'oculaire renvoie des faisceaux parallèles (image à l'infini). Les rayons de construction (orange) passent par le centre optique ou sont parallèles à l'axe incident.}
\end{popfigure}

\courssub{Parfocalité et parcentrage}

Les microscopes modernes portent plusieurs objectifs (par exemple 4$\times$, 10$\times$, 40$\times$, 100$\times$) montés sur une tourelle pivotante appelée \cle{revolver}. Deux qualités mécaniques rendent le changement d'objectif confortable :

\begin{definition}[Parfocalité et parcentrage]
\begin{itemize}
\item La \cle{parfocalité} : les objectifs montés sur le revolver sont conçus pour que l'image intermédiaire se forme (quasiment) dans le même plan quel que soit l'objectif utilisé. Après une mise au point soigneuse avec un objectif, le passage à un autre objectif ne demande qu'une retouche à la vis micrométrique — il est inutile (et déconseillé) de refaire la mise au point macrométrique.
\item Le \cle{parcentrage} : les objectifs sont réglés en usine pour que le centre du champ observé reste le même lorsqu'on change d'objectif. Ce réglage n'est pas accessible à l'utilisateur ; il garantit que la zone d'intérêt centrée à faible grossissement reste visible à fort grossissement.
\end{itemize}
\end{definition}

\courssub{Procédure pratique de mise au point au laboratoire}

\begin{methode}[Procédure de mise au point réelle]
\begin{enumerate}
\item Mettre en place l'objectif de \emph{plus faible grossissement} (4$\times$ ou 10$\times$) en faisant tourner le revolver jusqu'au déclic.
\item \textbf{En regardant de côté} (jamais dans l'oculaire !), actionner la vis \emph{macrométrique} pour monter la platine (ou descendre l'objectif) jusqu'à ce que l'objectif frôle la lame, sans la toucher.
\item Regarder maintenant dans l'oculaire et \emph{écarter} lentement la platine de l'objectif (toujours à la vis macrométrique) jusqu'à l'apparition de l'image.
\item Centrer la zone d'intérêt dans le champ en déplaçant doucement la lame à la main ou avec le chariot.
\item Affiner la netteté avec la vis \emph{micrométrique} uniquement.
\item Pour changer d'objectif : tourner le revolver (parfocalité) puis retoucher \emph{uniquement} la vis micrométrique.
\end{enumerate}
\end{methode}

\begin{attention}[Danger : ne jamais monter la platine en regardant dans l'oculaire]
L'erreur classique du débutant consiste à rapprocher l'objectif de la lame \emph{tout en regardant dans l'oculaire}. Comme on ne voit rien tant que la mise au point n'est pas faite, on continue de monter... jusqu'à ce que la frontale de l'objectif écrase la lame : lame cassée, objectif rayé, et parfois éclats de verre. Règle d'or : l'approche macrométrique se fait \textbf{en regardant de côté}, et l'on ne cherche l'image qu'en \emph{écartant} la platine de l'objectif.
\end{attention}

\begin{attention}[Ne pas confondre mise au point et réglage dioptrique]
La \cle{mise au point} (vis macro/micrométrique) positionne l'image intermédiaire par rapport à l'oculaire : elle dépend de l'échantillon. Le \cle{réglage dioptrique} de l'oculaire (bague crantée sur l'oculaire) compense le défaut de vision de l'observateur (myopie, hypermétropie) : il ne dépend que de l'œil. Il se règle \emph{une fois pour toutes}, sur un oculaire, en observant un réticule ou une préparation nette ; il ne faut pas le retoucher à chaque observation.
\end{attention}

\begin{savaistu}[La vis micrométrique, un instrument de mesure]
Sur les microscopes de qualité, la vis micrométrique a un pas de $\SI{0,5}{mm}$ ou $\SI{0,25}{mm}$ par tour, et son tambour gradué permet une lecture au $1/100$ de tour : la précision de déplacement vertical atteint donc $\SI{2,5}{\micro m}$ à $\SI{5}{\micro m}$. On peut ainsi mesurer l'\emph{épaisseur} d'un objet transparent : on fait la mise au point sur la face supérieure, puis sur la face inférieure, et l'on lit la différence au tambour (méthode dite de Berek, utilisée aussi en cristallographie). C'est le même principe que le palmer (vis micrométrique) utilisé en atelier de mécanique !
\end{savaistu}

\courssub{Mise au point sur un schéma de principe}

Au Probatoire, on te demandera souvent de \emph{construire} la marche des rayons et de vérifier graphiquement la mise au point. La démarche est la suivante :

\begin{methode}[Vérifier ou réaliser la mise au point sur un schéma]
\begin{enumerate}
\item Placer l'objectif $L_1$, l'oculaire $L_2$, leurs foyers, et l'objet $AB$.
\item Tracer les deux rayons utiles issus de $B$ à travers l'objectif : le rayon parallèle à l'axe (qui ressort par $F'_1$) et le rayon passant par $O_1$ (non dévié). Leur intersection donne $B_1$, image intermédiaire.
\item Comparer la position de $A_1$ à celle de $F_2$ :
\begin{itemize}
\item si $A_1$ est \emph{avant} $F_2$ (trop près de l'objectif) : l'objet est trop loin de $L_1$, il faut l'éloigner légèrement (diminuer $|p_{obj}|$) ;
\item si $A_1$ est \emph{après} $F_2$ : l'objet est trop près, il faut le rapprocher (augmenter $|p_{obj}|$) ;
\item si $A_1 \equiv F_2$ : la mise au point est réalisée, l'oculaire renvoie l'image finale à l'infini (rayons émergents parallèles entre eux).
\end{itemize}
\item Sur le schéma, déplacer l'objet (ou l'objectif) le long de l'axe et reprendre le tracé jusqu'à ce que $A_1$ coïncide avec $F_2$.
\end{enumerate}
\end{methode}

\begin{exoresolu}[Mise au point d'un microscope de laboratoire]
Un microscope est modélisé par un objectif de distance focale $f'_{obj} = \SI{5,0}{mm}$ et un oculaire de distance focale $f'_{oc} = \SI{25}{mm}$, séparés de telle sorte que l'intervalle optique vaut $\Delta = \SI{160}{mm}$. Un préparateur place une lame de telle sorte que l'objet $A$ soit à $\SI{5,3}{mm}$ devant le centre $O_1$ de l'objectif.
\begin{enumerate}
\item Calculer la position $p'_{obj} = \overline{O_1A_1}$ de l'image intermédiaire. La mise au point normale est-elle réalisée ?
\item De combien et dans quel sens faut-il déplacer la lame pour obtenir une image nette à l'infini ?
\end{enumerate}
\tcblower
\textbf{1.} La relation de conjugaison de Descartes appliquée à l'objectif s'écrit :
\[ \frac{1}{p'_{obj}} = \frac{1}{f'_{obj}} + \frac{1}{p_{obj}} = \frac{1}{\num{5,0}} + \frac{1}{-\num{5,3}} \quad (\text{en mm}). \]
D'où :
\[ \frac{1}{p'_{obj}} = \frac{\num{5,3} - \num{5,0}}{\num{5,0}\times\num{5,3}} = \frac{\num{0,3}}{\num{26,5}} \;\Rightarrow\; p'_{obj} = \frac{\num{26,5}}{\num{0,3}} \approx \SI{88,3}{mm}. \]
Pour la mise au point normale, il faudrait :
\[ p'_{obj} = f'_{obj} + \Delta = \num{5,0} + 160 = \SI{165}{mm}. \]
Or $\SI{88,3}{mm} < \SI{165}{mm}$ : l'image intermédiaire se forme \emph{avant} le foyer objet $F_2$ de l'oculaire. La mise au point n'est pas réalisée : l'image finale serait floue (l'oculaire recevrait des faisceaux convergents et formerait une image virtuelle mal placée).

\textbf{2.} Il faut augmenter $p'_{obj}$, donc rapprocher l'objet de l'objectif (diminuer $\lvert p_{obj} \rvert$). La position correcte de l'objet est :
\[ p_{obj} = -\frac{f'_{obj}\,(f'_{obj}+\Delta)}{\Delta} = -\frac{\num{5,0}\times 165}{160} \approx -\SI{5,16}{mm}. \]
Il faut donc rapprocher la lame de l'objectif de :
\[ \num{5,30} - \num{5,16} \approx \SI{0,14}{mm} = \SI{140}{\micro m}. \]
Un déplacement de seulement $\SI{140}{\micro m}$ suffit : cela justifie l'usage de la vis micrométrique (la vis macrométrique, à quelques millimètres par tour, serait bien trop grossière pour ce réglage).
\end{exoresolu}

\begin{exercice}[Pour s'entraîner]
Un microscope possède un objectif de distance focale $f'_{obj} = \SI{16}{mm}$ et un intervalle optique $\Delta = \SI{160}{mm}$.
\begin{enumerate}
\item Calculer la distance objet–objectif $p_{obj}$ pour une observation à l'infini (réglage normal).
\item L'opérateur éloigne la lame de $\SI{50}{\micro m}$ par rapport à cette position. Calculer la nouvelle position de l'image intermédiaire et conclure sur la netteté.
\item Pourquoi dit-on que l'objet est « pratiquement au foyer de l'objectif » ?
\end{enumerate}
\end{exercice}

\begin{corrige}[Exercice : mise au point]
\textbf{1.} $p_{obj} = -\dfrac{f'_{obj}(f'_{obj}+\Delta)}{\Delta} = -\dfrac{16\times 176}{160} = -\SI{17,6}{mm}$.

\textbf{2.} Nouvelle distance : $p_{obj} = -\SI{17,65}{mm}$. Alors :
\[ \frac{1}{p'_{obj}} = \frac{1}{16} - \frac{1}{\num{17,65}} = \frac{\num{17,65}-16}{16\times\num{17,65}} = \frac{\num{1,65}}{\num{282,4}} \;\Rightarrow\; p'_{obj} \approx \SI{171}{mm}. \]
L'image intermédiaire se déplace de $\SI{171}{mm} - \SI{176}{mm} = -\SI{5}{mm}$ : elle tombe $\SI{5}{mm}$ \emph{avant} $F_2$. Un déplacement de l'objet de $\SI{50}{\micro m}$ a provoqué un déplacement de l'image $100$ fois plus grand (facteur $\left(\Delta/f'_{obj}\right)^2 = 100$) : l'image devient floue, il faut retoucher la vis micrométrique.

\textbf{3.} $\lvert p_{obj} \rvert = \SI{17,6}{mm} \approx f'_{obj} = \SI{16}{mm}$ : l'écart relatif n'est que de $f'_{obj}/\Delta = 10\,\%$ environ, et il est d'autant plus faible que l'objectif est plus convergent. L'objet est donc placé quasiment au foyer objet de l'objectif, légèrement en avant.
\end{corrige}

\begin{aretenir}[L'essentiel de la mise au point]
\begin{itemize}
\item Image nette à l'infini $\Leftrightarrow$ image intermédiaire $A_1$ \emph{sur} le foyer objet $F_2$ de l'oculaire.
\item On règle la distance objet–objectif avec la vis \emph{macrométrique} (approche grossière, en regardant de côté) puis la vis \emph{micrométrique} (netteté parfaite).
\item La position correcte de l'objet est $p_{obj} \approx -f'_{obj}\left(1 + \dfrac{f'_{obj}}{\Delta}\right)$ : l'objet est presque au foyer de l'objectif, et la tolérance est de l'ordre du micromètre.
\item Grâce à la \emph{parfocalité}, le changement d'objectif sur le revolver ne demande qu'une retouche micrométrique ; le \emph{parcentrage} conserve le centre du champ.
\end{itemize}
\end{aretenir}

\coursec{Longueur du tube, encombrement et limites (Complément Probatoire)}

Après avoir maîtrisé la mise au point et le calcul des grossissements (section 4), nous devons comprendre pourquoi tous les microscopes ne se valent pas. La qualité de l'image finale ne dépend pas seulement du grossissement affiché sur les lentilles, mais de paramètres géométriques rigoureux (longueur de tube, ouverture numérique) et de limites physiques fondamentales (diffraction, aberrations). C'est ce qui sépare un jouet d'un instrument de recherche, et c'est exactement ce que le Probatoire attend de vous : \textbf{relier les caractéristiques gravées sur l'objectif à la performance réelle de l'instrument}.

\courssub{Longueur optique du tube et longueur mécanique}

Dans un microscope à tube fini (standard historique encore très présent dans nos lycées, comme les modèles Olympus CH-2 ou les microscopes chinois type XSZ), la distance mécanique entre l'épaulement de l'objectif (plan de référence où il se visse) et le plan de l'anneau de l'oculaire (où s'insère l'oculaire) est normalisée.

\begin{definition}[Longueur mécanique du tube (Standard RMS)]
C'est la distance physique, en millimètres, entre l'épaulement de l'objectif et le plan de l'anneau de l'oculaire. La norme \textbf{RMS (Royal Microscopical Society)} impose :
\[
L_{\text{mécanique}} = \SI{160}{mm}.
\]
C'est une constante constructive pour les microscopes dits « à tube fini ».
\end{definition}

Mais l'optique géométrique ne raisonne pas sur les épaulements métalliques, elle raisonne sur les \textbf{plans principaux} des systèmes optiques.

\begin{definition}[Longueur optique du tube ($\Delta_{\text{optique}}$ ou $L$)]
C'est la distance algébrique entre le foyer image de l'objectif ($F'_{\text{obj}}$) et le foyer objet de l'oculaire ($F_{\text{oc}}$). On la note souvent $\Delta$ (Delta) ou $L_{\text{optique}}$.
Pour des lentilles minces (approximation souvent utilisée en Première), les plans principaux coïncident avec les lentilles. Si $f'_{\text{obj}}$ et $f'_{\text{oc}}$ sont les focales image (positives) de l'objectif et de l'oculaire :
\[
L_{\text{mécanique}} \approx L_{\text{optique}} + f'_{\text{obj}} + f'_{\text{oc}} \quad \Longrightarrow \quad \Delta = L_{\text{optique}} = \SI{160}{mm} - f'_{\text{obj}} - f'_{\text{oc}}.
\]
\end{definition}

\begin{attention}[Piège : Confusion $\Delta$ et $L_{\text{mécanique}}$]
Au Probatoire, on vous donne souvent $L = \SI{160}{mm}$ comme « longueur du tube ». \textbf{Vérifiez toujours} s'il s'agit de la longueur mécanique (standard RMS) ou de l'intervalle optique $\Delta$.
\begin{itemize}
    \item Si l'énoncé dit « longueur du tube $L = \SI{160}{mm}$ » sans précision et que les focales sont données : c'est la longueur mécanique. Vous devez calculer $\Delta = \SI{160}{mm} - f'_{\text{obj}} - f'_{\text{oc}}$.
    \item Si l'énoncé dit « intervalle optique $\Delta = \dots$ » : utilisez directement cette valeur pour le grossissement $G_{\text{obj}} = -\Delta / f'_{\text{obj}}$.
\end{itemize}
Ne confondez jamais les deux sous peine de perdre tous les points de calcul.
\end{attention}

\begin{exemplebox}[Calcul de l'intervalle optique $\Delta$]
Un microscope scolaire (tube fini RMS) possède un objectif de focale $f'_{\text{obj}} = \SI{4,0}{mm}$ (gravure 40×) et un oculaire de focale $f'_{\text{oc}} = \SI{25,0}{mm}$ (gravure 10×).
Calculer l'intervalle optique $\Delta$ et le grossissement de l'objectif.
\begin{align*}
\Delta &= L_{\text{mécanique}} - f'_{\text{obj}} - f'_{\text{oc}} \\
       &= \SI{160}{mm} - \SI{4,0}{mm} - \SI{25,0}{mm} = \SI{131}{mm}. \\
G_{\text{obj}} &= -\frac{\Delta}{f'_{\text{obj}}} = -\frac{131}{4,0} \approx -32,8 \quad (\text{au lieu de } -40 \text{ si on utilisait } 160 \text{ mm à tort}).
\end{align*}
Le grossissement commercial 40× est une valeur normalisée (basée sur $\Delta_{\text{std}} = \SI{160}{mm}$), le grossissement réel dépend de $\Delta$ effectif.
\end{exemplebox}

\textbf{Microscopes à optique à l'infini (Modernes / Recherche).}\par
Sur les microscopes modernes (Leica, Zeiss, Nikon, Olympus recherche), l'objectif est corrigé pour donner une image \textbf{à l'infini} (faisceau sortant collimaté). Il n'y a plus de $\Delta$ fixe.
Un composant supplémentaire, la \textbf{lentille tube} (ou tube-lens, focale $f'_{\text{tube}} \approx \SI{200}{mm}$ pour Nikon/Leica, $\SI{180}{mm}$ pour Olympus), forme l'image intermédiaire dans son plan focal. L'oculaire observe ensuite cette image.
\begin{itemize}
    \item Avantage : On peut insérer des accessoires (filtres fluorescence, lames de Nomarski pour contraste de phase DIC, polariseurs) dans l'espace « infini » entre objectif et lentille tube sans perturber la correction optique.
    \item Le grossissement de l'objectif devient $G_{\text{obj}} = -f'_{\text{tube}} / f'_{\text{obj}}$.
\end{itemize}

\begin{popfigure}
\begin{tikzpicture}[scale=0.9, every node/.style={font=\small}]
    % Microscope à tube fini (haut)
    \draw[thick, popdark] (0,0) -- (8,0) node[midway, below=2mm] {Microscope à tube fini (RMS 160 mm)};
    % Objectif
    \draw[fill=popblueL] (0.5,-0.3) rectangle (1.5,0.3) node[midway] {Obj};
    % Oculaire
    \draw[fill=poporangeL] (7,-0.3) rectangle (8,0.3) node[midway] {Ocu};
    % Foyers
    \draw[dashed, popblue] (1.5,0) -- (2.5,0) node[midway, above] {$f'_{\text{obj}}$};
    \draw[dashed, poporange] (6,0) -- (7,0) node[midway, above] {$f'_{\text{oc}}$};
    % Delta
    \draw[<->, popgreen, thick] (2.5,0.6) -- (6,0.6) node[midway, above] {$\Delta = L_{\text{optique}}$};
    % Longueur mécanique
    \draw[<->, poppurple, thick] (0.5,-0.8) -- (8,-0.8) node[midway, below] {$L_{\text{méc}} = \SI{160}{mm}$};

    % Microscope à l'infini (bas)
    \begin{scope}[yshift=-3.5cm]
    \draw[thick, popdark] (0,0) -- (9,0) node[midway, below=2mm] {Microscope à optique à l'infini (Modernes)};
    \draw[fill=popblueL] (0.5,-0.3) rectangle (1.5,0.3) node[midway] {Obj};
    \draw[fill=popgreenL] (4.5,-0.3) rectangle (5.5,0.3) node[midway] {Lentille Tube};
    \draw[fill=poporangeL] (8,-0.3) rectangle (9,0.3) node[midway] {Ocu};
    % Faisceaux
    \draw[->, popblue, thick] (1.5,0.2) -- (4.5,0.2) node[midway, above] {Faisceau collimaté (infini)};
    \draw[->, popblue, thick] (1.5,-0.2) -- (4.5,-0.2);
    \draw[->, popgreen, thick] (5.5,0.2) -- (6.5,0) -- (8,0.2);
    \draw[->, popgreen, thick] (5.5,-0.2) -- (6.5,0) -- (8,-0.2);
    % Focales
    \draw[dashed, popblue] (1.5,0) -- (0.5,0); % objet au foyer obj
    \draw[dashed, popgreen] (5.5,0) -- (6.5,0) node[midway, above] {$f'_{\text{tube}}$};
    \draw[dashed, poporange] (7,0) -- (8,0) node[midway, above] {$f'_{\text{oc}}$};
    \end{scope}
\end{tikzpicture}
\legende{Comparaison des géométries : Tube fini (standard scolaire) vs Optique à l'infini (standard recherche).}
\end{popfigure}

\courssub{Pouvoir séparateur (Résolution) et Ouverture Numérique}

Grossir une image floue ne sert à rien. La limite fondamentale de tout instrument d'optique est la \textbf{diffraction}. Deux points sources très proches produisent des taches d'Airy qui se chevauchent. Le critère de Rayleigh définit la distance minimale $d$ pour les distinguer.

\begin{definition}[Ouverture Numérique (NA — Numerical Aperture)]
C'est la grandeur clé qui caractérise la capacité d'un objectif à collecter la lumière et à résoudre les détails. Elle est gravée sur le corps de l'objectif (ex: 40/0,65 ou 100/1,30 Oil).
\[
\boxed{\text{NA} = n \sin(\alpha_{\text{max}})}
\]
où :
\begin{itemize}
    \item $n$ : indice de réfraction du milieu entre l'objet (lame) et la lentille frontale de l'objectif.
    \item $\alpha_{\text{max}}$ : demi-angle d'ouverture maximal du cône de lumière entrant dans l'objectif.
\end{itemize}
\textbf{Valeurs typiques :} Air ($n=1$) $\Rightarrow$ NA max $\approx 0,95$. Huile d'immersion ($n=1,515$) $\Rightarrow$ NA max $\approx 1,4 - 1,45$.
\end{definition}

\begin{propriete}[Critère de Rayleigh — Pouvoir séparateur (Admis au Probatoire)]
Deux points sont \textbf{juste résolus} si le maximum central de la tache de diffraction de l'un coïncide avec le premier minimum de l'autre.
La distance minimale résoluble (pouvoir séparateur) dans le plan de l'objet est :
\[
\boxed{d_{\text{min}} = 0,61 \frac{\lambda}{\text{NA}}}
\]
$\lambda$ : longueur d'onde dans le vide (on prend souvent $\lambda = \SI{550}{nm}$, vert, sensibilité max de l'œil).
$d_{\text{min}}$ est en mètres (ou nm, $\mu$m) si $\lambda$ est dans la même unité.
\emph{Analyse dimensionnelle :} $[\lambda] = L$, $[NA] = 1$ $\Rightarrow$ $[d] = L$. Correct.
\end{propriete}

\begin{popfigure}
\begin{tikzpicture}
\begin{axis}[
    width=\linewidth, height=8cm,
    xlabel={Ouverture Numérique NA}, ylabel={Résolution $d_{\text{min}}$ (\si{\micro m})},
    domain=0.1:1.45, samples=100,
    xmin=0, xmax=1.5, ymin=0, ymax=3.5,
    grid=major, grid style={popline},
    axis line style={popdark},
    tick style={popdark},
    label style={font=\small, popdark},
    tick label style={font=\footnotesize, popdark},
    legend style={at={(0.98,0.98)}, anchor=north east, font=\small, fill=popcream, draw=popline}
]
% Courbe théorique lambda = 550 nm = 0.55 um
\addplot[popcurve, thick, domain=0.1:1.45] {0.61*0.55/x};
\addlegendentry{$d = 0,61 \lambda / \text{NA}$ ($\lambda=\SI{550}{nm}$)}

% Points repères objectifs standards
\addplot[only marks, mark=*, mark size=2.5, popblue] coordinates {(0.10, 3.355)}; % 4x NA 0.1
\addplot[only marks, mark=square*, mark size=2.5, poporange] coordinates {(0.65, 0.516)}; % 40x NA 0.65
\addplot[only marks, mark=triangle*, mark size=2.5, popgreen] coordinates {(1.30, 0.258)}; % 100x Oil NA 1.30
\addplot[only marks, mark=diamond*, mark size=2.5, poppurple] coordinates {(1.40, 0.240)}; % 100x Oil NA 1.40

% Annotations
\node[popblue, anchor=south west, font=\tiny] at (axis cs:0.10,3.355) {4$\times$ (NA 0,1) : 3,3 $\mu$m};
\node[poporange, anchor=north west, font=\tiny] at (axis cs:0.65,0.516) {40$\times$ (NA 0,65) : 0,5 $\mu$m};
\node[popgreen, anchor=north west, font=\tiny] at (axis cs:1.30,0.258) {100$\times$ Huile (NA 1,30) : 0,26 $\mu$m};

% Zone grossissement vide (indicatif)
\addplot[poppinkL, fill opacity=0.1, domain=1.0:1.45] {0} \closedcycle;
\node[poppink, rotate=90, anchor=south, font=\small\bfseries] at (axis cs:1.42, 1.5) {Risque grossissement vide si $G > 1000\times\text{NA}$};
\end{axis}
\end{tikzpicture}
\legende{Évolution du pouvoir séparateur $d_{\text{min}}$ en fonction de l'Ouverture Numérique (NA) pour $\lambda = \SI{550}{nm}$. Les points correspondent aux objectifs standards.}
\end{popfigure}

\begin{exemplebox}[Résolution réelle d'un objectif 100× Huile]
Objectif : \textbf{100× / 1,30 Oil} (NA = 1,30, $n_{\text{huile}}=1,515$). $\lambda = \SI{550}{nm}$.
\begin{align*}
d_{\text{min}} &= 0,61 \times \frac{\SI{550}{nm}}{1,30} = 0,61 \times 423 \text{ nm} \approx \SI{258}{nm} = \SI{0,258}{\micro m}.
\end{align*}
Cet objectif peut distinguer deux bactéries espacées de $\SI{0,26}{\micro m}$.
Si on utilise le même objectif \textbf{sans huile} (air, $n=1$, NA max $\approx 0,95$ mais l'objectif n'est pas corrigé pour l'air, NA effectif chute drastiquement $\sim 0,6$), la résolution tombe à $\sim \SI{0,56}{\micro m}$ : \textbf{perte de facteur 2}. L'huile n'est pas optionnelle pour les objectifs « Oil » !
\end{exemplebox}

\courssub{Grossissement utile et grossissement vide}

Le grossissement total $G = G_{\text{obj}} \times G_{\text{oc}}$ agrandit la tache de diffraction. L'œil humain a une acuité angulaire d'environ $1'$ (un minute d'angle) $\approx \SI{3e-4}{rad}$. À la distance de vision distincte $D = \SI{250}{mm}$, l'œil sépare deux points s'ils sont écartés de $d_{\text{œil}} \approx D \times \SI{3e-4}{rad} \approx \SI{0,075}{mm} = \SI{75}{\micro m}$.

Pour que le détail résolu par l'objectif ($d_{\text{min}}$) devienne visible à l'œil, il faut :
\[
G \times d_{\text{min}} \gtrsim d_{\text{œil}} \quad \Rightarrow \quad G \gtrsim \frac{\SI{75}{\micro m}}{d_{\text{min}}}.
\]
En utilisant $d_{\text{min}} = 0,61 \lambda / \text{NA}$ avec $\lambda = \SI{550}{nm}$ :
\[
G_{\text{min utile}} \approx \frac{75 \times 10^{-6}}{0,61 \times 550 \times 10^{-9} / \text{NA}} \approx 224 \times \text{NA}.
\]
On retient généralement la règle empirique (validée par l'expérience) :

\begin{aretenir}[Grossissement utile maximal]
\[
\boxed{G_{\text{max utile}} \approx 500 \text{ à } 1000 \times \text{NA}}
\]
\begin{itemize}
    \item \textbf{Minimum utile} ($\approx 500 \times \text{NA}$) : Le détail résolu devient juste visible par l'œil.
    \item \textbf{Maximum utile} ($\approx 1000 \times \text{NA}$) : Au-delà, l'image est plus grande mais \textbf{pas plus nette}. C'est le \textbf{grossissement vide} (empty magnification) : on agrandit le flou de diffraction.
\end{itemize}
\end{aretenir}

\begin{attention}[Erreur classique au Probatoire : « Mon microscope fait 2500× ! »]
Un élève achète un oculaire 25× pour son objectif 100× (NA=1,30). $G = 2500\times$.
$G_{\text{max utile}} = 1000 \times 1,30 = 1300\times$.
\textbf{Résultat :} Grossissement vide. L'image est grande, floue, sombre, et ne révèle \textbf{aucune information supplémentaire} par rapport à un grossissement 1000× (oculaire 10×). Perte de luminosité ($L \propto 1/G^2$) et de champ de vision en prime.
\textbf{Conseil :} Choisissez l'oculaire pour que $G_{\text{total}} \in [500\text{NA}, 1000\text{NA}]$. Pour NA=1,30 : $G \in [650, 1300]$. Oculaire 10× ($G=1000\times$) est idéal. Oculaire 15× ($G=1500\times$) est déjà limite.
\end{attention}

\begin{exemplebox}[Choix de l'oculaire pour un objectif 40× (NA=0,65)]
$G_{\text{obj}} \approx 40\times$.
$G_{\text{max utile}} \in [500 \times 0,65 ; 1000 \times 0,65] = [325 ; 650]$.
Oculaires disponibles : 5×, 10×, 15×, 20×.
\begin{itemize}
    \item Oculaire 5× : $G=200\times$ \textbf{Trop faible} (détail non visible par l'œil).
    \item Oculaire 10× : $G=400\times$ \textbf{Bon} (dans la plage utile).
    \item Oculaire 15× : $G=600\times$ \textbf{Acceptable} (proche limite haute).
    \item Oculaire 20× : $G=800\times$ \textbf{Grossissement vide}.
\end{itemize}
\end{exemplebox}

\courssub{Aberrations, corrections et épaisseur de lame}

Un objectif n'est pas une lentille unique. C'est un système complexe (jusqu'à 15 lentilles) corrigé pour plusieurs défauts.

\begin{definition}[Principales aberrations corrigées]
\begin{itemize}
    \item \textbf{Aberration chromatique :} $f'$ dépend de $\lambda$ (dispersion). Image colorée sur les bords.
        \begin{itemize}
            \item \textbf{Achromate :} Corrigé pour 2 longueurs d'onde (rouge, bleu). Foyer commun pour le rouge et le bleu. Standard scolaire. Gravure : \textbf{Achro} ou rien.
            \item \textbf{Apochromate (Apo) :} Corrigé pour 3 longueurs d'onde (rouge, vert, bleu) + correction sphérique améliorée. Haute performance. Gravure : \textbf{Plan Apo} ou \textbf{Apo}.
            \item \textbf{Plan-achromate (Plan) :} Champ d'image \textbf{plan} (correction de la courbure de champ). Net jusqu'aux bords. Essentiel pour photo/microscopie numérique. Gravure : \textbf{Plan}.
        \end{itemize}
    \item \textbf{Aberration sphérique :} Rayons marginaux focalisent différemment des paraxiaux. Corrigée par combinaisons de verres (couronnes/flints) ou lentilles asphériques.
\end{itemize}
\end{definition}

\begin{attention}[L'épaisseur du couvre-objet : 0,17 mm standard !]
Les objectifs à fort grossissement (40×, 60×, 100×) sont conçus pour travailler \textbf{à travers un couvre-objet d'épaisseur standard $e = \SI{0,17}{mm}$} (indice $n \approx 1,52$).
\begin{itemize}
    \item Gravure sur l'objectif : \textbf{0,17} ou \textbf{/0} (signifie corrigé pour 0,17 mm).
    \item Si vous mettez une lame plus épaisse (ex: $\SI{0,5}{mm}$) ou pas de lame : \textbf{aberration sphérique introduite}, perte nette de résolution et de contraste, surtout en 100× huile.
    \item Objectifs « sans couvre-objet » (gravure \textbf{0} ou \textbf{NCG - No Cover Glass}) existent pour métallurgie ou cultures en milieu liquide (fond de boîte de Pétri $\sim \SI{1}{mm}$).
\end{itemize}
\end{attention}

\begin{savaistu}[Le Microscope Électronique : Dépasser la limite de la lumière]
La résolution optique est limitée par $\lambda$ (visible $\approx \SI{500}{nm}$ $\Rightarrow d_{\text{min}} \approx \SI{200}{nm}$).
Le \textbf{Microscope Électronique à Transmission (MET)} utilise un faisceau d'électrons accélérés à $\SI{100}{kV}$ - $\SI{300}{kV}$.
Longueur d'onde de Broglie : $\lambda = h/p \approx \SI{0,0037}{nm}$ (à 100 kV).
Résolution théorique $\ll \SI{1}{nm}$. Pratique : $\sim \SI{0,1}{nm}$ (résolution atomique !).
Le \textbf{Microscope Électronique à Balayage (MEB)} : résolution $\sim \SI{1}{nm} - \SI{10}{nm}$, image 3D de surface.
\textbf{Contraintes :} Vide poussé ($10^{-4}$ à $10^{-6}$ Pa), échantillon conducteur (métallisation or/palladium pour MEB), coupes ultrafines ($< \SI{100}{nm}$) pour MET. Coût : millions de FCFA. Indispensable en virologie (IRAD, Centre Pasteur Cameroun), science des matériaux, nanotechnologies.
\end{savaistu}

\courssub{Exercice résolu type Probatoire : Analyse complète d'un microscope}

\begin{exoresolu}[Caractérisation d'un microscope de laboratoire]
\textbf{Énoncé :} On utilise un microscope à tube fini (standard RMS $L_{\text{méc}} = \SI{160}{mm}$) équipé de :
\begin{itemize}
    \item Objectif : \textbf{Plan 40× / 0,65 / 0,17} (focale $f'_{\text{obj}} = \SI{4,0}{mm}$).
    \item Oculaire : \textbf{WF 10× / 20} (focale $f'_{\text{oc}} = \SI{25,0}{mm}$, champ apparent $2\omega' = 20$ mm).
    \item Éclairage : Lampe LED blanche ($\lambda_{\text{moy}} = \SI{550}{nm}$), diaphragme d'ouverture réglé au max.
    \item Préparation : Lame standard ($e = \SI{1,0}{mm}$) + Couvre-objet standard ($e = \SI{0,17}{mm}$).
\end{itemize}
\textbf{Questions :}
\begin{enumerate}
    \item Calculer l'intervalle optique $\Delta$ et le grossissement réel de l'objectif $G_{\text{obj}}$. Comparer à la gravure 40×.
    \item Calculer le grossissement total $G$ et le diamètre du champ objet visible.
    \item Calculer le pouvoir séparateur $d_{\text{min}}$ (critère de Rayleigh).
    \item Déterminer si le grossissement total est utile, vide ou insuffisant.
    \item L'utilisateur retire le couvre-objet pour observer une goutte d'eau sans lame. Quel défaut apparaît ? Pourquoi ?
\end{enumerate}

\tcblower
\textbf{Solution détaillée :}

\textbf{1. Intervalle optique et grossissement objectif.}
\begin{align*}
\Delta &= L_{\text{méc}} - f'_{\text{obj}} - f'_{\text{oc}} \\
       &= \SI{160}{mm} - \SI{4,0}{mm} - \SI{25,0}{mm} = \SI{131}{mm}. \\
G_{\text{obj}} &= -\frac{\Delta}{f'_{\text{obj}}} = -\frac{131}{4,0} = -32,75 \approx -32,8.
\end{align*}
\textit{Commentaire :} La gravure « 40× » correspond à une définition normalisée basée sur $\Delta_{\text{std}} = \SI{160}{mm}$ ($G_{\text{norm}} = -160/4 = -40$). Le grossissement réel est inférieur de $\approx 18\%$ car l'oculaire prend de la place dans le tube.

\textbf{2. Grossissement total et champ objet.}
$G_{\text{oc}} = \frac{\SI{250}{mm}}{f'_{\text{oc}}} = \frac{250}{25} = 10$ (vérifie la gravure 10×).
$G = G_{\text{obj}} \times G_{\text{oc}} = (-32,75) \times 10 = -327,5 \approx \textbf{-328}$.
Diamètre champ objet $D_{\text{obj}} = \frac{\text{Champ apparent oculaire}}{|G_{\text{obj}}|} = \frac{\SI{20}{mm}}{32,75} \approx \SI{0,61}{mm} = \SI{610}{\micro m}$.
(L'oculaire WF 10×/20 a un diaphragme de champ de 20 mm de diamètre dans son plan focal).

\textbf{3. Pouvoir séparateur (Résolution).}
NA = 0,65 (air). $\lambda = \SI{550}{nm}$.
$d_{\text{min}} = 0,61 \frac{\lambda}{\text{NA}} = 0,61 \times \frac{550}{0,65} \approx \SI{516}{nm} = \SI{0,516}{\micro m}$.
\textit{Interprétation :} Deux points distants de $\SI{0,52}{\micro m}$ sur la préparation sont juste distinguables dans l'image finale.

\textbf{4. Utilité du grossissement.}
$G_{\text{max utile}} \approx 1000 \times \text{NA} = 1000 \times 0,65 = 650$.
$G_{\text{min utile}} \approx 500 \times \text{NA} = 325$.
Notre $G = 328$.
\textbf{Conclusion :} Le grossissement est \textbf{juste au seuil minimal utile} ($328 \approx 325$). Le détail résolu ($0,52 \mu m$) est agrandi à $328 \times 0,52 \approx \SI{170}{\mu m}$ sur la rétine, juste au-dessus du seuil de l'œil ($\sim \SI{75}{\mu m}$). C'est acceptable mais limite. Un oculaire 15× ($G \approx 490$) serait plus confortable pour l'observation. Un oculaire 5× ($G \approx 164$) serait \textbf{insuffisant} (grossissement sous-optimal).

\textbf{5. Retrait du couvre-objet (Épaisseur 0 au lieu de 0,17 mm).}
L'objectif est gravé \textbf{/0,17} (corrigé pour couvre-objet 0,17 mm).
Sans couvre-objet, le milieu entre l'objet et la lentille frontale est l'air ($n=1$) sur une distance où l'objectif « attend » du verre ($n=1,52$).
\textbf{Conséquence :} Introduction d'une \textbf{aberration sphérique positive} (sous-correction). Les rayons marginaux sont moins déviés que prévu, ils focalisent plus loin que les rayons paraxiaux.
\textbf{Résultat visuel :} Image \textbf{floue}, perte de contraste, « voile » lumineux, résolution effective très dégradée (bien pire que $d_{\text{min}}$ théorique). L'objectif 40×/0,65 devient inutilisable pour du travail fin sans son couvre-objet standard.
\end{exoresolu}

\begin{exercice}[Entraînement : Objectif 100× Huile]
Un microscope à tube fini ($L_{\text{méc}} = \SI{160}{mm}$) possède un objectif \textbf{100× / 1,30 Oil / 0,17} ($f'_{\text{obj}} = \SI{1,6}{mm}$) et un oculaire \textbf{10×} ($f'_{\text{oc}} = \SI{25}{mm}$).
\begin{enumerate}
    \item Calculer $\Delta$ et $G_{\text{obj}}$ réel. Pourquoi la gravure 100× diffère-t-elle ?
    \item Calculer $d_{\text{min}}$ pour $\lambda = \SI{550}{nm}$.
    \item Déterminer la plage de grossissement utile. L'oculaire 10× est-il adapté ? Et un oculaire 15× ?
    \item L'utilisateur met de l'huile d'immersion ($n=1,515$) mais utilise un couvre-objet d'épaisseur $\SI{0,13}{mm}$ (trop fin). Expliquer l'effet sur l'image.
\end{enumerate}
\end{exercice}

\begin{corrige}[Exercice : Objectif 100× Huile]
\begin{enumerate}
    \item $\Delta = 160 - 1,6 - 25 = \SI{133,4}{mm}$. $G_{\text{obj}} = -133,4 / 1,6 = -83,4$. Gravure 100× basée sur $\Delta_{\text{std}}=160$ mm ($160/1,6=100$).
    \item $d_{\text{min}} = 0,61 \times 550 / 1,30 \approx \SI{258}{nm}$.
    \item $G_{\text{utile}} \in [650 ; 1300]$. $G(10\times) = 834$ \textbf{Parfait (au milieu de la plage)}. $G(15\times) = 1251$ \textbf{Acceptable (proche max)}. $G(20\times) = 1668$ \textbf{Vide}.
    \item Couvre-objet trop fin ($\SI{0,13}{mm}$ au lieu de $\SI{0,17}{mm}$) $\Rightarrow$ Aberration sphérique \textbf{positive} (sous-correction). Les rayons marginaux focalisent \textbf{après} les paraxiaux (plus loin). Image floue, anneaux de diffraction visibles. Qualité dégradée. Il faut impérativement $e = \SI{0,17}{mm} \pm \SI{0,01}{mm}$.
\end{enumerate}
\end{corrige}

\coursec{Résolution de problèmes types (Probatoire)}

\courssub{Transition et stratégie générale}

Après avoir étudié la géométrie du microscope, ses grandeurs caractéristiques (grandissement, puissance, longueur de tube) et la mise au point, nous abordons maintenant la résolution systématique des problèmes types que vous rencontrerez au Probatoire. La clé de la réussite réside dans une \textbf{méthode rigoureuse} : identifier les données, choisir le mode de fonctionnement (normal ou confort), appliquer les bonnes formules sans confusion d'unités, et vérifier la cohérence des résultats.

\begin{methode}[Stratégie universelle « Microscope au Probatoire »]
\begin{enumerate}
    \item \textbf{Lister les données} en les convertissant \emph{systématiquement} en unités SI (mètres pour les puissances, millimètres pour les grandissements) ou en les gardant cohérentes selon la formule utilisée.
    \item \textbf{Identifier le mode de fonctionnement} : \emph{normal} (image finale à l'infini, œil relâché) ou \emph{confort} (image finale à 25 cm, distance de vision distincte). L'énoncé le précise souvent (« microscope réglé en vision normale » ou « pour un œil normal »).
    \item \textbf{Choisir les formules} adaptées (voir boîte récapitulative ci-dessous).
    \item \textbf{Calculer étape par étape} : puissances, grandissements partiels, grandissement total, longueur du tube.
    \item \textbf{Vérifier les ordres de grandeur} : $f'_{\text{obj}}$ de l'ordre du mm, $f'_{\text{oc}}$ de l'ordre de la cm, $\Delta \approx 160$ mm, $G$ entre 50 et 1500, $P$ entre 50 et 200 $\delta$.
    \item \textbf{Phraser la réponse finale} avec unité, encadrée, et commentaire physique si demandé (signe du grandissement, pouvoir séparateur, etc.).
\end{enumerate}
\end{methode}

\begin{aretenir}[Formules de référence — À connaître par cœur]
Pour un microscope de longueur d'intervalle optique $\Delta$ (souvent 160 mm), objectif de focale image $f'_{\text{obj}}$ et oculaire de focale image $f'_{\text{oc}}$ :
\begin{align*}
    \text{Puissances :} \quad & P_{\text{obj}} = \frac{1}{f'_{\text{obj}}},\quad P_{\text{oc}} = \frac{1}{f'_{\text{oc}}},\quad P = P_{\text{obj}} + P_{\text{oc}} \quad (\text{en dioptries } \delta,\ f' \text{ en m}) \\
    \text{Grandissement objectif :} \quad & \gamma_{\text{obj}} = -\frac{\Delta}{f'_{\text{obj}}} \quad (\Delta,\ f'_{\text{obj}} \text{ en mm}) \\
    \text{Grandissement oculaire (normal) :} \quad & G_{\text{oc}} = \frac{250}{f'_{\text{oc}}} \quad (f'_{\text{oc}} \text{ en mm}) \\
    \text{Grandissement oculaire (confort) :} \quad & G_{\text{oc}} = \frac{250}{f'_{\text{oc}}} + 1 \quad (f'_{\text{oc}} \text{ en mm}) \\
    \text{Grandissement commercial :} \quad & G = |\gamma_{\text{obj}}| \times G_{\text{oc}} \\
    \text{Longueur mécanique (approx.) :} \quad & L \approx \Delta + f'_{\text{obj}} + f'_{\text{oc}} \quad (\text{tout en mm})
\end{align*}
\end{aretenir}

\begin{attention}[Piège mortel : les unités !]
C'est la \textbf{première cause de perte de points} au Probatoire.
\begin{itemize}
    \item Pour les \textbf{puissances} ($P$, $P_{\text{obj}}$, $P_{\text{oc}}$) : les focales \textbf{doivent être en mètres}. Exemple : $f'_{\text{obj}} = 4\ \text{mm} = 0,004\ \text{m} \Rightarrow P_{\text{obj}} = 250\ \delta$.
    \item Pour les \textbf{grandissements} ($\gamma_{\text{obj}}$, $G_{\text{oc}}$, $G$) : on garde \textbf{généralement les millimètres} car $\Delta = 160\ \text{mm}$ et $250\ \text{mm}$ (distance de vision distincte) sont en mm. Exemple : $\gamma_{\text{obj}} = -160/4 = -40$ (sans unité).
    \item \textbf{Ne jamais mélanger m et mm dans une même formule.} Convertissez tout au début.
\end{itemize}
\end{attention}

\begin{popfigure}
\centering
\begin{tikzpicture}[
    node distance=1.2cm and 1.5cm,
    startstop/.style={rectangle, rounded corners, minimum width=3cm, minimum height=1cm, text centered, draw=popblue, fill=popblueL, thick},
    process/.style={rectangle, minimum width=3.5cm, minimum height=1cm, text centered, draw=popdark, fill=popcream, thick},
    decision/.style={diamond, aspect=1.5, minimum width=2.5cm, minimum height=1cm, text centered, draw=poporange, fill=poporangeL, thick},
    arrow/.style={-Stealth, thick, popdark},
    check/.style={rectangle, rounded corners, minimum width=3cm, minimum height=1cm, text centered, draw=popgreen, fill=popgreenL, thick}
]
\node (start) [startstop] {DÉPART : Lire l'énoncé};
\node (data) [process, below of=start] {Lister données + convertir unités};
\node (mode) [decision, below of=data, align=center]{Mode précisé ?\\Normal / Confort};
\node (formG) [process, below left of=mode, xshift=-1.5cm, align=center]{Formules G :\\$\gamma_{\text{obj}}=-\Delta/f'_{\text{obj}}$\\$G_{\text{oc}}=250/f'_{\text{oc}}$ (+1 si confort)\\$G=|\gamma_{\text{obj}}|\times G_{\text{oc}}$};
\node (formP) [process, below right of=mode, xshift=1.5cm, align=center]{Formules P :\\$P_{\text{obj}}=1/f'_{\text{obj}}$\\$P_{\text{oc}}=1/f'_{\text{oc}}$\\$P=P_{\text{obj}}+P_{\text{oc}}$};
\node (calc) [process, below of=mode, yshift=-2.5cm] {Calculer étape par étape};
\node (verify) [check, below of=calc, align=center]{Vérifier ordres de grandeur\\+ cohérence signes};
\node (end) [startstop, below of=verify] {RÉPONSE FINALE encadrée};

\draw [arrow] (start) -- (data);
\draw [arrow] (data) -- (mode);
\draw [arrow] (mode) -- node[left, pos=0.5, text width=2cm, align=center] {\small Normal} (formG);
\draw [arrow] (mode) -- node[right, pos=0.5, text width=2cm, align=center] {\small Confort} (formG);
\draw [arrow] (mode) -- (formP);
\draw [arrow] (formG) -- (calc);
\draw [arrow] (formP) -- (calc);
\draw [arrow] (calc) -- (verify);
\draw [arrow] (verify) -- (end);
\end{tikzpicture}
\legende{Arbre de décision pour la résolution de problèmes de microscope au Probatoire.}
\end{popfigure}

\courssub{Type 1 : Données $f'_{\text{obj}}$, $f'_{\text{oc}}$, $\Delta$ — Calculer $G$, $P$, $L$}

C'est le type le plus classique. On vous donne les focales image de l'objectif et de l'oculaire (souvent en mm ou cm) et l'intervalle optique $\Delta$ (souvent 160 mm). On demande le grandissement commercial (normal et/ou confort), la puissance commerciale et la longueur du tube.

\begin{exemplebox}[Type 1 — Exemple complet]
\textbf{Énoncé.} Un microscope est constitué d'un objectif de focale image $f'_{\text{obj}} = 4,0\ \text{mm}$ et d'un oculaire de focale image $f'_{\text{oc}} = 25\ \text{mm}$. L'intervalle optique est $\Delta = 160\ \text{mm}$.
\begin{enumerate}
    \item Calculer la puissance commerciale $P$ (en $\delta$).
    \item Calculer le grandissement commercial $G$ en vision normale et en vision de confort.
    \item Déterminer la longueur mécanique approximative $L$ du microscope.
\end{enumerate}
\end{exemplebox}

\begin{exoresolu}[Type 1 — Solution détaillée]
\textbf{Données (conversion en SI pour les puissances) :}
$f'_{\text{obj}} = 4,0\ \text{mm} = 4,0 \times 10^{-3}\ \text{m}$,
$f'_{\text{oc}} = 25\ \text{mm} = 2,5 \times 10^{-2}\ \text{m}$,
$\Delta = 160\ \text{mm}$ (on garde mm pour les grandissements).

\textbf{1. Puissance commerciale $P$ :}
\[
P_{\text{obj}} = \frac{1}{f'_{\text{obj}}} = \frac{1}{4,0 \times 10^{-3}} = 250\ \delta
\]
\[
P_{\text{oc}} = \frac{1}{f'_{\text{oc}}} = \frac{1}{2,5 \times 10^{-2}} = 40\ \delta
\]
\[
\boxed{P = P_{\text{obj}} + P_{\text{oc}} = 250 + 40 = 290\ \delta}
\]

\textbf{2. Grandissement commercial $G$ :}
Grandissement de l'objectif (toujours négatif, image inversée) :
\[
\gamma_{\text{obj}} = -\frac{\Delta}{f'_{\text{obj}}} = -\frac{160}{4,0} = -40
\]

\begin{itemize}
    \item \textbf{Vision normale} (image finale à l'infini, œil relâché) :
    \[
    G_{\text{oc}} = \frac{250}{f'_{\text{oc}}} = \frac{250}{25} = 10
    \]
    \[
    \boxed{G_{\text{normal}} = |\gamma_{\text{obj}}| \times G_{\text{oc}} = 40 \times 10 = 400}
    \]

    \item \textbf{Vision de confort} (image finale à 25 cm, œil accommodé au maximum) :
    \[
    G_{\text{oc}} = \frac{250}{f'_{\text{oc}}} + 1 = 10 + 1 = 11
    \]
    \[
    \boxed{G_{\text{confort}} = |\gamma_{\text{obj}}| \times G_{\text{oc}} = 40 \times 11 = 440}
    \]
\end{itemize}

\textbf{3. Longueur mécanique $L$ :}
\[
L \approx \Delta + f'_{\text{obj}} + f'_{\text{oc}} = 160 + 4,0 + 25 = 189\ \text{mm}
\]
\[
\boxed{L \approx 18,9\ \text{cm}}
\]

\textbf{Vérification :} $P \approx 290\ \delta$ (cohérent : objectif fort, oculaire modéré), $G \approx 400$ (microscope standard), $L \approx 19\ \text{cm}$ (encombrement réaliste).
\end{exoresolu}

\courssub{Type 2 : Données $G$ visé, $\Delta$, $f'_{\text{oc}}$ — Déterminer $f'_{\text{obj}}$ et choisir dans un catalogue}

Ici, on inverse le problème : on veut obtenir un grandissement donné, on connaît l'oculaire et $\Delta$, il faut trouver la focale de l'objectif nécessaire, puis sa puissance, et enfin choisir l'objectif standard le plus proche (catalogue : $4\times$, $10\times$, $40\times$, $100\times$).

\begin{exemplebox}[Type 2 — Choix dans un catalogue]
\textbf{Énoncé.} On souhaite réaliser un microscope de grandissement commercial $G = 600$ en vision normale. L'intervalle optique est $\Delta = 160\ \text{mm}$ et l'oculaire disponible a une focale $f'_{\text{oc}} = 20\ \text{mm}$. Le catalogue propose des objectifs de grandissements nominaux $4\times$, $10\times$, $40\times$, $100\times$ (pour $\Delta = 160\ \text{mm}$).
\begin{enumerate}
    \item Calculer la focale image $f'_{\text{obj}}$ nécessaire.
    \item En déduire la puissance $P_{\text{obj}}$.
    \item Quel objectif du catalogue choisir ? Justifier.
\end{enumerate}
\end{exemplebox}

\begin{exoresolu}[Type 2 — Solution détaillée]
\textbf{Données :} $G = 600$ (normal), $\Delta = 160\ \text{mm}$, $f'_{\text{oc}} = 20\ \text{mm}$.

\textbf{1. Grandissement oculaire (normal) :}
\[
G_{\text{oc}} = \frac{250}{f'_{\text{oc}}} = \frac{250}{20} = 12,5
\]

\textbf{2. Grandissement objectif requis :}
\[
G = |\gamma_{\text{obj}}| \times G_{\text{oc}} \quad \Rightarrow \quad |\gamma_{\text{obj}}| = \frac{G}{G_{\text{oc}}} = \frac{600}{12,5} = 48
\]

\textbf{3. Focale image de l'objectif :}
\[
|\gamma_{\text{obj}}| = \frac{\Delta}{f'_{\text{obj}}} \quad \Rightarrow \quad f'_{\text{obj}} = \frac{\Delta}{|\gamma_{\text{obj}}|} = \frac{160}{48} \approx 3,33\ \text{mm}
\]
\[
\boxed{f'_{\text{obj}} \approx 3,3\ \text{mm}}
\]

\textbf{4. Puissance de l'objectif :}
\[
P_{\text{obj}} = \frac{1}{f'_{\text{obj}}} = \frac{1}{3,33 \times 10^{-3}} \approx 300\ \delta
\]
\[
\boxed{P_{\text{obj}} \approx 300\ \delta}
\]

\textbf{5. Choix dans le catalogue :}
Les objectifs du catalogue sont caractérisés par leur grandissement nominal $\gamma_{\text{cat}} = -\Delta/f'_{\text{cat}}$ pour $\Delta = 160\ \text{mm}$.
\begin{itemize}
    \item $4\times \Rightarrow f' = 40\ \text{mm}$
    \item $10\times \Rightarrow f' = 16\ \text{mm}$
    \item $40\times \Rightarrow f' = 4,0\ \text{mm}$
    \item $100\times \Rightarrow f' = 1,6\ \text{mm}$
\end{itemize}
Notre $f'_{\text{obj}} \approx 3,3\ \text{mm}$ est le plus proche de $4,0\ \text{mm}$ (objectif $40\times$).
Avec l'objectif $40\times$ ($f'_{\text{obj}} = 4,0\ \text{mm}$), on obtient $|\gamma_{\text{obj}}| = 40$ et $G = 40 \times 12,5 = 500$.
Avec l'objectif $100\times$ ($f'_{\text{obj}} = 1,6\ \text{mm}$), on obtient $|\gamma_{\text{obj}}| = 100$ et $G = 100 \times 12,5 = 1250$.

\textbf{Choix :} L'objectif \textbf{$40\times$} est le plus adapté pour approcher $G = 600$ (on obtient 500). L'objectif $100\times$ donnerait un grossissement trop fort (1250) et une profondeur de champ trop faible.
\[
\boxed{\text{Choisir l'objectif } 40\times \ (f'_{\text{obj}} = 4,0\ \text{mm})}
\]
\end{exoresolu}

\courssub{Type 3 : Marquage objectif (ex: 40/0,65/160/0,17) et oculaire (10×) — Interprétation complète}

C'est un \textbf{classique du Probatoire}. Il faut savoir lire le marquage normalisé d'un objectif : \textbf{Grossissement / Ouverture numérique / Longueur de tube / Épaisseur couvre-objet}.

\begin{definition}[Marquage normalisé d'un objectif (ex: 40/0,65/160/0,17)]
\begin{itemize}
    \item \textbf{40} : Grandissement nominal $\gamma_{\text{nom}} = -40$ (pour $\Delta = 160\ \text{mm}$).
    \item \textbf{0,65} : Ouverture numérique $NA = n \sin \alpha_{\text{max}} = 0,65$ (air, $n=1$).
    \item \textbf{160} : Longueur de tube mécanique normalisée $L = 160\ \text{mm}$ (ancienne norme) ou intervalle optique $\Delta = 160\ \text{mm}$.
    \item \textbf{0,17} : Épaisseur du couvre-objet (en mm) pour laquelle l'objectif est corrigé. \emph{Important :} si on observe sans couvre-objet ou avec une épaisseur différente, des aberrations sphériques apparaissent.
\end{itemize}
\end{definition}

\begin{exemplebox}[Type 3 — Lecture de marquage]
\textbf{Énoncé.} Un objectif porte le marquage \textbf{40/0,65/160/0,17}. L'oculaire est marqué \textbf{10×}. On suppose $\Delta = 160\ \text{mm}$.
\begin{enumerate}
    \item Déterminer $f'_{\text{obj}}$, $f'_{\text{oc}}$, $P_{\text{obj}}$, $P_{\text{oc}}$, $P$.
    \item Calculer $G$ en vision normale.
    \item Calculer la distance minimale résolue $d_{\text{min}}$ pour $\lambda = 550\ \text{nm}$ (lumière verte).
    \item Interpréter le chiffre 0,17.
\end{enumerate}
\end{exemplebox}

\begin{exoresolu}[Type 3 — Solution détaillée]
\textbf{Données :} Objectif $40/0,65/160/0,17$, Oculaire $10\times$, $\Delta = 160\ \text{mm}$.

\textbf{1. Focales et puissances :}
\begin{itemize}
    \item Objectif : $\gamma_{\text{nom}} = -40 = -\Delta/f'_{\text{obj}} \Rightarrow f'_{\text{obj}} = 160/40 = 4,0\ \text{mm} = 4,0 \times 10^{-3}\ \text{m}$.
    \[
    P_{\text{obj}} = \frac{1}{4,0 \times 10^{-3}} = 250\ \delta
    \]
    \item Oculaire : $G_{\text{oc}} = 10 = 250/f'_{\text{oc}} \Rightarrow f'_{\text{oc}} = 250/10 = 25\ \text{mm} = 2,5 \times 10^{-2}\ \text{m}$.
    \[
    P_{\text{oc}} = \frac{1}{2,5 \times 10^{-2}} = 40\ \delta
    \]
    \item Puissance commerciale :
    \[
    \boxed{P = P_{\text{obj}} + P_{\text{oc}} = 250 + 40 = 290\ \delta}
    \]
\end{itemize}

\textbf{2. Grandissement commercial (normal) :}
\[
\gamma_{\text{obj}} = -40,\quad G_{\text{oc}} = 10 \quad \Rightarrow \quad \boxed{G = 40 \times 10 = 400}
\]

\textbf{3. Pouvoir séparateur (critère de Rayleigh) :}
$NA = 0,65$ (air, $n=1$), $\lambda = 550\ \text{nm} = 5,5 \times 10^{-7}\ \text{m}$.
\[
d_{\text{min}} = 0,61 \frac{\lambda}{NA} = 0,61 \times \frac{5,5 \times 10^{-7}}{0,65} \approx 5,16 \times 10^{-7}\ \text{m} = 516\ \text{nm}
\]
\[
\boxed{d_{\text{min}} \approx 520\ \text{nm} \ (0,52\ \mu\text{m})}
\]

\textbf{4. Interprétation du 0,17 :}
Ce chiffre indique que l'objectif est \textbf{corrigé pour un couvre-objet d'épaisseur $0,17\ \text{mm}$} (couvre-objet standard n°1,5). Si on observe une lame sans couvre-objet ou avec un couvre-objet d'épaisseur différente (ex: 0,5 mm), l'objectif introduit une \textbf{aberration sphérique} qui dégrade l'image. Pour les objectifs à forte ouverture numérique ($NA > 0,4$), cette correction est cruciale. Les objectifs « sans couvre-objet » portent la mention « 0 » ou « – » à la place de 0,17.
\end{exoresolu}

\begin{savaistu}[Pourquoi 0,17 mm exactement ?]
Historiquement, les couvre-objets standards (norme ISO 8255) ont une épaisseur nominale de $0,17\ \text{mm}$ (tolérance $\pm 0,01\ \text{mm}$). Les objectifs à fort $NA$ (plan-apochromatiques, immersion huile) sont conçus pour que les rayons traversant le couvre-objet convergent parfaitement en son absence d'aberration sphérique. En microscopie moderne (microscopes à tube infini), la correction est souvent déplacée dans le tube optique (objectifs « ∞ »).
\end{savaistu}

\courssub{Type 4 : Mise au point — Déplacement de l'objectif pour passer du normal au confort}

La mise au point se fait en déplaçant l'objectif (ou la platine) pour que l'image intermédiaire se forme au foyer image de l'oculaire (vision normale) ou légèrement en avant (vision confort). On utilise la relation de conjugaison de Newton ou de Descartes pour l'objectif.

\begin{propriete}[Relations de mise au point — Objectif]
Pour l'objectif de focale $f'_{\text{obj}}$, soit $p_{\text{obj}} = \overline{F_{\text{obj}}A}$ (distance objet-foyer objet, algébrique, négative si objet réel) et $p'_{\text{obj}} = \overline{F'_{\text{obj}}A'}$ (distance image-foyer image, algébrique, positive si image réelle).
\begin{itemize}
    \item \textbf{Relation de Newton :} $p_{\text{obj}} \cdot p'_{\text{obj}} = -f'_{\text{obj}}{}^2$
    \item \textbf{Relation de Descartes :} $\frac{1}{p_{\text{obj}}} + \frac{1}{p'_{\text{obj}}} = \frac{1}{f'_{\text{obj}}}$ (avec signes algébriques : $p_{\text{obj}} < 0$, $p'_{\text{obj}} > 0$, $f'_{\text{obj}} > 0$)
\end{itemize}
L'image intermédiaire $A'B'$ doit se former au foyer objet de l'oculaire $F_{\text{oc}}$ pour la vision normale, ou à une distance $p_{\text{oc}}$ de l'oculaire telle que $p'_{\text{oc}} = -25\ \text{cm}$ pour la vision de confort.
\end{propriete}

\begin{exemplebox}[Type 4 — Mise au point normal → confort]
\textbf{Énoncé.} Un microscope a un objectif de focale $f'_{\text{obj}} = 4,0\ \text{mm}$. En vision normale, l'objet est placé à une distance $p_{\text{obj}} = -4,2\ \text{mm}$ du foyer objet de l'objectif (c'est-à-dire $\overline{F_{\text{obj}}A} = -4,2\ \text{mm}$). L'oculaire a une focale $f'_{\text{oc}} = 25\ \text{mm}$.
\begin{enumerate}
    \item Calculer $p'_{\text{obj}}$ en vision normale. Vérifier que l'image intermédiaire est bien au foyer objet de l'oculaire.
    \item Pour passer en vision de confort (image finale à $25\ \text{cm}$), on déplace l'objectif. Calculer le nouveau $p'_{\text{obj}}$.
    \item En déduire le déplacement $\Delta p_{\text{obj}}$ de l'objectif (ou de la platine). Sens du déplacement ?
\end{enumerate}
\end{exemplebox}

\begin{exoresolu}[Type 4 — Solution détaillée]
\textbf{Données :} $f'_{\text{obj}} = 4,0\ \text{mm}$, $p_{\text{obj}} = -4,2\ \text{mm}$ (normal), $f'_{\text{oc}} = 25\ \text{mm}$, $\Delta = 160\ \text{mm}$.

\textbf{1. Vision normale :}
Relation de Newton : $p_{\text{obj}} \cdot p'_{\text{obj}} = -f'_{\text{obj}}{}^2$
\[
(-4,2) \cdot p'_{\text{obj}} = -(4,0)^2 = -16 \quad \Rightarrow \quad p'_{\text{obj}} = \frac{16}{4,2} \approx 3,81\ \text{mm}
\]
L'image intermédiaire est à $3,81\ \text{mm}$ du foyer image de l'objectif. Comme $\Delta = 160\ \text{mm}$, l'image est à $160 - 3,81 \approx 156,2\ \text{mm}$ de l'objectif, soit très proche du foyer objet de l'oculaire (à $160\ \text{mm}$ de $F'_{\text{obj}}$). C'est cohérent.

\textbf{2. Vision de confort :}
Pour l'oculaire, image finale à $p'_{\text{oc}} = -250\ \text{mm}$ (virtuelle, même côté que l'objet pour l'oculaire).
Relation de Descartes pour l'oculaire : $\frac{1}{p_{\text{oc}}} + \frac{1}{p'_{\text{oc}}} = \frac{1}{f'_{\text{oc}}}$
\[
\frac{1}{p_{\text{oc}}} + \frac{1}{-250} = \frac{1}{25} \quad \Rightarrow \quad \frac{1}{p_{\text{oc}}} = \frac{1}{25} + \frac{1}{250} = \frac{11}{250}
\]
\[
p_{\text{oc}} = \frac{250}{11} \approx 22,7\ \text{mm}
\]
L'image intermédiaire doit donc se former à $22,7\ \text{mm}$ de l'oculaire (côté objet), c'est-à-dire à une distance $p'_{\text{obj}} = \Delta - p_{\text{oc}} = 160 - 22,7 = 137,3\ \text{mm}$ du foyer image de l'objectif.
\[
\boxed{p'_{\text{obj}}(\text{confort}) \approx 137,3\ \text{mm}}
\]

\textbf{3. Déplacement de l'objectif :}
En vision normale : $p'_{\text{obj}} \approx 3,8\ \text{mm}$.
En vision confort : $p'_{\text{obj}} \approx 137,3\ \text{mm}$.
L'image intermédiaire s'éloigne de l'objectif de $\Delta p'_{\text{obj}} = 137,3 - 3,8 = 133,5\ \text{mm}$.
Pour maintenir la conjugaison, il faut \textbf{rapprocher l'objectif de l'objet} (diminuer $|p_{\text{obj}}|$).
Calculons le nouveau $p_{\text{obj}}$ (Newton) : $p_{\text{obj}} \cdot 137,3 = -16 \Rightarrow p_{\text{obj}} \approx -0,116\ \text{mm}$.
L'objet passe de $-4,2\ \text{mm}$ à $-0,116\ \text{mm}$ du foyer : il se rapproche de $4,1\ \text{mm}$ du foyer.
\[
\boxed{\text{Déplacement de l'objectif (ou platine) : } \approx 4,1\ \text{mm vers l'objet (rapprochement)}}
\]

\begin{attention}[Piège : sens du déplacement]
Pour passer du \textbf{normal au confort}, l'image intermédiaire doit se former \textbf{plus loin} de l'objectif (plus près de l'oculaire). Comme l'objectif est convergent, il faut \textbf{rapprocher l'objet de l'objectif} (diminuer la distance objet-objectif). En pratique, on \textbf{monte la platine} (ou on descend l'objectif) de quelques mm. Au Probatoire, précisez toujours : « On rapproche l'objectif de l'objet » ou « On monte la platine ».
\end{attention}
\end{exoresolu}

\courssub{Type 5 : Résolution — Pouvoir séparateur, observation de bactéries, virus, atomes}

Le pouvoir séparateur (résolution) est la plus petite distance $d_{\text{min}}$ entre deux points distincts que l'instrument peut discriminer. Critère de Rayleigh pour un microscope : $d_{\text{min}} = 0,61 \frac{\lambda}{NA}$.

\begin{exemplebox}[Type 5 — Résolution et observabilité]
\textbf{Énoncé.} Un microscope a un objectif d'ouverture numérique $NA = 1,30$ (immersion huile, $n=1,52$). On utilise une lumière verte $\lambda = 550\ \text{nm}$.
\begin{enumerate}
    \item Calculer la distance minimale résolue $d_{\text{min}}$.
    \item Peut-on distinguer deux bactéries de taille $1\ \mu\text{m}$ espacées de $0,5\ \mu\text{m}$ ?
    \item Peut-on distinguer des virus de $100\ \text{nm}$ ?
    \item Peut-on voir des atomes ($0,1\ \text{nm}$) ?
    \item Que proposez-vous pour améliorer la résolution ?
\end{enumerate}
\end{exemplebox}

\begin{exoresolu}[Type 5 — Solution détaillée]
\textbf{Données :} $NA = 1,30$, $\lambda = 550\ \text{nm} = 5,5 \times 10^{-7}\ \text{m}$.

\textbf{1. Distance minimale résolue :}
\[
d_{\text{min}} = 0,61 \frac{\lambda}{NA} = 0,61 \times \frac{5,5 \times 10^{-7}}{1,30} \approx 2,58 \times 10^{-7}\ \text{m} = 258\ \text{nm}
\]
\[
\boxed{d_{\text{min}} \approx 260\ \text{nm} \ (0,26\ \mu\text{m})}
\]

\textbf{2. Bactéries ($1\ \mu\text{m}$, espacement $0,5\ \mu\text{m} = 500\ \text{nm}$) :}
$500\ \text{nm} > 260\ \text{nm}$ \Rightarrow \textbf{OUI}, on peut les distinguer (résolution suffisante).

\textbf{3. Virus ($100\ \text{nm}$) :}
$100\ \text{nm} < 260\ \text{nm}$ \Rightarrow \textbf{NON}, on ne peut pas les distinguer (taille inférieure à la résolution). Ils apparaîtront comme des taches floues de $\sim 260\ \text{nm}$.

\textbf{4. Atomes ($0,1\ \text{nm}$) :}
$0,1\ \text{nm} \ll 260\ \text{nm}$ \Rightarrow \textbf{ABSOLUMENT PAS}. La microscopie optique ne voit pas les atomes. Il faut de la microscopie électronique (MET, MEB) ou à effet tunnel (STM).

\textbf{5. Améliorer la résolution :}
D'après $d_{\text{min}} = 0,61 \frac{\lambda}{NA}$, on peut :
\begin{itemize}
    \item Diminuer $\lambda$ : utiliser de la lumière bleue ($\lambda \approx 450\ \text{nm}$) ou UV (mais attention aux dommages biologiques et à l'absorption du verre).
    \item Augmenter $NA = n \sin \alpha_{\text{max}}$ :
    \begin{itemize}
        \item Augmenter $n$ : \textbf{immersion huile} ($n=1,52$) ou eau ($n=1,33$) au lieu d'air ($n=1$). Déjà fait ici ($NA=1,30$).
        \item Augmenter $\alpha_{\text{max}}$ : objectif à plus grand angle (haute ouverture), mais limité par la géométrie ($\sin \alpha_{\text{max}} \le 1 \Rightarrow NA_{\text{max}} = n$). En huile, $NA_{\text{max}} \approx 1,52$.
    \end{itemize}
    \item \textbf{Microscopie à fluorescence / super-résolution (STED, PALM/STORM) :} contourne la limite de diffraction (prix Nobel 2014).
\end{itemize}
\end{exoresolu}

\begin{savaistu}[Microscopie au Cameroun — Contexte réel]
Au Centre Pasteur du Cameroun (Yaoundé) et dans les laboratoires de l'Université de Yaoundé I, les microscopes à immersion huile ($100\times$, $NA=1,30$) sont standards pour le diagnostic du paludisme (recherche de \emph{Plasmodium} dans le sang, taille $\sim 1-2\ \mu\text{m}$) et de la tuberculose (bacilles de Koch, $\sim 0,5 \times 2\ \mu\text{m}$). La résolution de $\sim 250\ \text{nm}$ est amplement suffisante. Pour les virus (VIH, Ebola, COVID-19), on utilise la PCR (biologie moléculaire) ou la microscopie électronique au Centre de Recherche en Santé de la Défense (CRESAR) ou à l'étranger. Le microscope optique reste l'outil de base, robuste et peu coûteux, pour l'hématologie, la parasitologie et la bactériologie de routine.
\end{savaistu}

\courssub{Type 6 : Aberration chromatique — Objectif simple vs doublet achromate}

Un objectif simple (une lentille) a une focale qui dépend de la longueur d'onde (dispersion) : $f'_{\text{bleu}} < f'_{\text{rouge}}$. L'image d'un objet blanc présente des franges colorées. La solution est le \textbf{doublet achromate} (couronne + flint).

\begin{exemplebox}[Type 6 — Aberration chromatique]
\textbf{Énoncé.} Un objectif simple en verre couronne a une focale $f'_{\text{rouge}} = 4,2\ \text{mm}$ pour $\lambda_R = 656\ \text{nm}$ (raie C) et $f'_{\text{bleu}} = 3,8\ \text{mm}$ pour $\lambda_B = 486\ \text{nm}$ (raie F).
\begin{enumerate}
    \item Calculer la dispersion de focale $\Delta f' = f'_{\text{rouge}} - f'_{\text{bleu}}$.
    \item Expliquer pourquoi l'image d'un objet blanc est colorée.
    \item On réalise un doublet achromate en associant cette lentille couronne ($V_1 = 60$, $f'_1 = 4,0\ \text{mm}$ pour $\lambda_d = 589\ \text{nm}$) à une lentille flint ($V_2 = 36$). Calculer la focale $f'_2$ de la lentille flint pour obtenir un doublet achromate ($f'_{\text{rouge}} = f'_{\text{bleu}}$).
\end{enumerate}
\end{exemplebox}

\begin{exoresolu}[Type 6 — Solution détaillée]
\textbf{Données :} $f'_{\text{rouge}} = 4,2\ \text{mm}$, $f'_{\text{bleu}} = 3,8\ \text{mm}$, $V_1 = 60$ (couronne), $V_2 = 36$ (flint), $f'_1 = 4,0\ \text{mm}$ (ligne jaune sodium).

\textbf{1. Dispersion de focale :}
\[
\Delta f' = f'_{\text{rouge}} - f'_{\text{bleu}} = 4,2 - 3,8 = 0,4\ \text{mm}
\]
\[
\boxed{\Delta f' = 0,4\ \text{mm}}
\]

\textbf{2. Image colorée :}
Le verre a un indice de réfraction plus grand pour le bleu que pour le rouge ($n_B > n_R$). Donc la lentille est plus convergente pour le bleu : $f'_{\text{bleu}} < f'_{\text{rouge}}$. L'image bleue se forme \textbf{plus près} de la lentille que l'image rouge. Sur un écran placé pour le vert, on voit une tache bleue nette au centre et un halo rouge flou (ou inversement selon la position de l'écran). C'est l'\textbf{aberration chromatique longitudinale}.

\textbf{3. Doublet achromate :}
Condition d'achromatisme pour deux lentilles minces jointes : $\frac{P_1}{V_1} + \frac{P_2}{V_2} = 0$ où $P = 1/f'$.
$P_1 = 1/f'_1 = 1/0,004 = 250\ \delta$.
\[
\frac{250}{60} + \frac{P_2}{36} = 0 \quad \Rightarrow \quad P_2 = -36 \times \frac{250}{60} = -150\ \delta
\]
\[
f'_2 = \frac{1}{P_2} = \frac{1}{-150} \approx -6,67 \times 10^{-3}\ \text{m} = -6,67\ \text{mm}
\]
\[
\boxed{f'_2 \approx -6,7\ \text{mm} \text{ (lentille divergente en flint)}}
\]
Focale du doublet : $P = P_1 + P_2 = 250 - 150 = 100\ \delta \Rightarrow f' = 10\ \text{mm}$.
Le doublet corrige l'aberration chromatique pour le rouge et le bleu (même focale), mais une \textbf{aberration résiduelle} (secondaire) subsiste pour le vert (spectre apochromatique nécessite 3 lentilles).
\end{exoresolu}

\courssub{Exercices types Probatoire — Entraînement}

\begin{exercice}[Exercice 1 — Calcul complet]
Un microscope a un objectif de puissance $100\ \delta$ et un oculaire de focale $2,0\ \text{cm}$. L'intervalle optique est $\Delta = 16\ \text{cm}$.
\begin{enumerate}
    \item Calculer les focales $f'_{\text{obj}}$ et $f'_{\text{oc}}$ en mm.
    \item Calculer le grandissement commercial en vision normale et en vision de confort.
    \item Calculer la puissance commerciale $P$.
    \item Calculer la longueur mécanique $L$.
\end{enumerate}
\end{exercice}

\begin{corrige}[Exercice 1]
\textbf{Données :} $P_{\text{obj}} = 100\ \delta \Rightarrow f'_{\text{obj}} = 1/100 = 0,01\ \text{m} = 10\ \text{mm}$.
$f'_{\text{oc}} = 2,0\ \text{cm} = 20\ \text{mm}$.
$\Delta = 16\ \text{cm} = 160\ \text{mm}$.
\begin{enumerate}
    \item $f'_{\text{obj}} = 10\ \text{mm}$, $f'_{\text{oc}} = 20\ \text{mm}$.
    \item $\gamma_{\text{obj}} = -160/10 = -16$.
    Normal : $G_{\text{oc}} = 250/20 = 12,5 \Rightarrow G = 16 \times 12,5 = \boxed{200}$.
    Confort : $G_{\text{oc}} = 12,5 + 1 = 13,5 \Rightarrow G = 16 \times 13,5 = \boxed{216}$.
    \item $P_{\text{oc}} = 1/0,02 = 50\ \delta$. $P = 100 + 50 = \boxed{150\ \delta}$.
    \item $L \approx 160 + 10 + 20 = 190\ \text{mm} = \boxed{19,0\ \text{cm}}$.
\end{enumerate}
\end{corrige}

\begin{exercice}[Exercice 2 — Choix d'objectif]
On veut obtenir $G = 800$ en vision normale avec $\Delta = 160\ \text{mm}$ et un oculaire $15\times$. Catalogue objectifs : $4\times$, $10\times$, $40\times$, $100\times$ (pour $\Delta=160\ \text{mm}$). Lequel choisir ? Calculer le $G$ réel obtenu.
\end{exercice}

\begin{corrige}[Exercice 2]
$G_{\text{oc}} = 15$. $|\gamma_{\text{obj}}| = 800/15 \approx 53,3$.
Objectif $40\times \Rightarrow |\gamma|=40 \Rightarrow G=600$.
Objectif $100\times \Rightarrow |\gamma|=100 \Rightarrow G=1500$.
Choix : \textbf{objectif $40\times$} (plus proche de 800, bien que 600 soit en dessous). L'objectif $100\times$ donne un grossissement trop fort (vide).
$G_{\text{réel}} = \boxed{600}$.
\end{corrige}

\begin{exercice}[Exercice 3 — Marquage et résolution]
Objectif marqué \textbf{100/1,30/160/0,17}. Oculaire \textbf{10×}. $\lambda = 550\ \text{nm}$.
\begin{enumerate}
    \item $f'_{\text{obj}}$, $f'_{\text{oc}}$, $P$, $G$ (normal).
    \item $d_{\text{min}}$.
    \item Peut-on voir un virus de $120\ \text{nm}$ ? Une bactérie de $1\ \mu\text{m}$ ?
    \item Signification du 0,17 ?
\end{enumerate}
\end{exercice}

\begin{corrige}[Exercice 3]
\begin{enumerate}
    \item $\gamma_{\text{obj}}=-100 \Rightarrow f'_{\text{obj}}=1,6\ \text{mm}$. $G_{\text{oc}}=10 \Rightarrow f'_{\text{oc}}=25\ \text{mm}$.
    $P_{\text{obj}}=625\ \delta$, $P_{\text{oc}}=40\ \delta$, $P=665\ \delta$.
    $G = 100 \times 10 = \boxed{1000}$.
    \item $d_{\text{min}} = 0,61 \times 550 / 1,30 \approx 258\ \text{nm}$.
    \item Virus $120\ \text{nm} < 258\ \text{nm}$ : \textbf{Non}. Bactérie $1000\ \text{nm} > 258\ \text{nm}$ : \textbf{Oui}.
    \item Épaisseur du couvre-objet pour laquelle l'objectif est corrigé ($0,17\ \text{mm}$ = standard n°1,5).
\end{enumerate}
\end{corrige}

\begin{exercice}[Exercice 4 — Mise au point]
Objectif $f'_{\text{obj}} = 5,0\ \text{mm}$. En vision normale, $p_{\text{obj}} = -5,5\ \text{mm}$. Oculaire $f'_{\text{oc}} = 20\ \text{mm}$. $\Delta = 160\ \text{mm}$.
Calculer le déplacement de la platine pour passer en vision de confort ($p'_{\text{oc}} = -25\ \text{cm}$). Sens ?
\end{exercice}

\begin{corrige}[Exercice 4]
Normal : $p'_{\text{obj}} = -f'^2/p_{\text{obj}} = -25/(-5,5) = 4,55\ \text{mm}$.
Confort : $1/p_{\text{oc}} = 1/20 + 1/250 = 0,054 \Rightarrow p_{\text{oc}} = 18,5\ \text{mm}$.
$p'_{\text{obj}} = \Delta - p_{\text{oc}} = 160 - 18,5 = 141,5\ \text{mm}$.
Nouveau $p_{\text{obj}} = -f'^2/p'_{\text{obj}} = -25/141,5 = -0,177\ \text{mm}$.
Déplacement objet : $|-5,5 - (-0,177)| = 5,32\ \text{mm}$.
\boxed{\text{Rapprocher la platine de l'objectif de } \approx 5,3\ \text{mm}}.
\end{corrige}

\courssub{Check-list finale avant de rendre la copie (Probatoire)}

\begin{aretenir}[Check-list « Zéro faute » — Microscope]
\begin{itemize}
    \item[$\square$] \textbf{Unités cohérentes} : Puissances en $\delta$ ($f'$ en m) ; Grandissements avec $\Delta$ et $f'$ en mm.
    \item[$\square$] \textbf{Signes respectés} : $\gamma_{\text{obj}} < 0$ (image inversée), $p_{\text{obj}} < 0$, $p'_{\text{obj}} > 0$, $p'_{\text{oc}} < 0$ (image virtuelle).
    \item[$\square$] \textbf{Mode précisé} : « Vision normale » ou « Vision de confort » écrit explicitement avant le calcul de $G_{\text{oc}}$.
    \item[$\square$] \textbf{Formule $G$} : $G = |\gamma_{\text{obj}}| \times G_{\text{oc}}$ (produit, \textbf{jamais} somme).
    \item[$\square$] \textbf{Formule $P$} : $P = P_{\text{obj}} + P_{\text{oc}}$ (somme des puissances).
    \item[$\square$] \textbf{Réponse encadrée} : $\boxed{G = ...}$, $\boxed{P = ...\ \delta}$, $\boxed{L = ...\ \text{cm}}$.
    \item[$\square$] \textbf{Justifications} : Signe de $\gamma$ (image inversée), rôle du couvre-objet (0,17 mm), immersion huile (augmente $NA$), aberration chromatique (doublet achromate).
\end{itemize}
\end{aretenir}

\begin{attention}[Derniers pièges à éviter]
\begin{itemize}
    \item Confondre \textbf{grandissement nominal de l'objectif} (ex: $40\times$) et \textbf{grandissement commercial du microscope} (ex: $400\times$). Le premier est $\gamma_{\text{obj}}$, le second est $G$.
    \item Oublier le \textbf{$+1$} pour la vision de confort : $G_{\text{oc}} = 250/f'_{\text{oc}} + 1$.
    \item Utiliser $\Delta = 160\ \text{mm}$ en \textbf{mètres} dans la formule de $\gamma_{\text{obj}}$ (donnerait $\gamma_{\text{obj}} = -0,16$ au lieu de $-160$).
    \item Dire « l'image est droite » au lieu de « l'image est inversée » (grandissement négatif).
    \item Confondre \textbf{pouvoir séparateur} ($d_{\text{min}}$) et \textbf{grandissement}. Un grand $G$ ne garantit pas une bonne résolution si $NA$ est faible (grossissement vide).
\end{itemize}
\end{attention}

\begin{savaistu}[L'huile d'immersion au Probatoire — Le point bonus]
Si l'énoncé mentionne « objectif à immersion huile » ou donne $NA > 1$ (ex: 1,30), sachez que :
\begin{itemize}
    \item L'huile ($n \approx 1,52$) remplit l'espace entre l'objectif et le couvre-objet.
    \item Cela supprime la réflexion/réfraction à l'interface air-verre, augmentant l'angle $\alpha_{\text{max}}$ capté.
    \item $NA = n \sin \alpha_{\text{max}}$ peut dépasser 1 (max $\approx 1,4-1,5$).
    \item \textbf{Résultat :} $d_{\text{min}}$ divisé par $\sim 1,5$ par rapport à l'air. C'est \textbf{indispensable} pour le $100\times$.
    \item Au Probatoire, on demande souvent : « Pourquoi utilise-t-on de l'huile ? » $\rightarrow$ « Pour augmenter l'indice $n$, donc le $NA$, donc la résolution (diminuer $d_{\text{min}}$). »
\end{itemize}
\end{savaistu}

\vspace{0,5cm}
\noindent\textbf{Fin de la leçon 9 — Le microscope.} Vous maîtrisez maintenant la description, le schéma, les calculs de $G$ et $P$, la mise au point, la lecture des marquages, la résolution et les aberrations. Entraînez-vous sur les exercices types ci-dessus : ce sont \emph{exactement} les questions qui tombent au Probatoire. Bon courage et bon travail !

\newpage

\coursec{Activité d'intégration}

\coursec{Le microscope}

\courssub{Description et rôle du microscope}
Le microscope optique composé est un instrument d'optique destiné à l'observation d'objets de très petite taille (cellules, bactéries, tissus) invisibles à l'œil nu. Il réalise un \textbf{grossissement angulaire} en deux étapes successives grâce à deux systèmes convergents :
\begin{itemize}
    \item L'\textbf{objectif} (face à l'objet), de courte distance focale ($f_{obj}' \approx \text{quelques mm}$), forme une \textbf{image réelle, inversée et grossie} $A'B'$ (image intermédiaire) dans le plan focal objet de l'oculaire.
    \item L'\textbf{oculaire} (face à l'œil), de distance focale plus grande ($f_{oc}' \approx \text{1 à 3 cm}$), fonctionne comme une \textbf{loupe} : il donne une \textbf{image finale virtuelle, inversée par rapport à l'objet initial} (donc redressée par rapport à l'image intermédiaire), située à l'infini (réglage normal) ou au point proche (réglage confort).
\end{itemize}
\cle{Microscope} \cle{Objectif} \cle{Oculaire} \cle{Image intermédiaire} \cle{Image finale}

\begin{savaistu}[Contexte camerounais -- Microscopie et santé publique]
Au Cameroun, le microscope optique est l'outil de diagnostic de première ligne dans les laboratoires d'analyses médicales (hôpitaux, centres de santé, laboratoires privés). Il est indispensable pour :
\begin{itemize}
    \item Le diagnostic du \textbf{paludisme} (recherche de \emph{Plasmodium} sur frottis sanguin épais et mince) -- première cause de consultation.
    \item La \textbf{bactériologie} : identification des bacilles (tuberculose, méningite, choléra), recherche de germes dans les urines, selles, LCR.
    \item L'hématologie (formule leucocytaire), la parasitologie intestinale, la cytologie.
\end{itemize}
La maîtrise de son fonctionnement, de sa mise au point et de ses performances (grossissement, résolution) est une compétence vitale pour les futurs techniciens de laboratoire, biologistes et médecins formés en série C.
\end{savaistu}

\courssub{Schéma de principe et construction géométrique}
\begin{experience}[Construction géométrique de l'image par un microscope -- Réglage normal (image à l'infini)]
\textbf{Matériel (simulation ou banc d'optique) :} Objectif ($f_{obj}' \approx 4\,\text{mm}$), Oculaire ($f_{oc}' \approx 25\,\text{mm}$), Objet (flèche) à $x_{obj} \approx -4,5\,\text{mm}$ de l'objectif, Écran pour image intermédiaire, Règle, Crayon.

\textbf{Protocole :}
\begin{enumerate}
    \item Placer l'objet $AB$ légèrement au-delà du foyer objet de l'objectif ($OF_{obj} < AB < 2f_{obj}'$).
    \item Tracer les rayons émergents de l'objectif : rayon parallèle à l'axe $\to$ passe par $F_{obj}'$ ; rayon passant par le centre optique $O_{obj}$ $\to$ non dévié. Ils se coupent en $A'B'$ : \textbf{image réelle, inversée, grossie} ($G_{obj} > 1$).
    \item Placer l'oculaire de sorte que $A'B'$ se trouve \textbf{dans son plan focal objet} $F_{oc}$ (distance $F_{oc}A'B' = f_{oc}'$).
    \item Tracer les rayons émergents de l'oculaire : depuis $A'B'$, rayon passant par centre $O_{oc}$ ; rayon parallèle à l'axe $\to$ réfracté vers $F_{oc}'$. Les rayons émergents sont \textbf{parallèles} : image finale à l'infini.
    \item L'œil placé derrière l'oculaire reçoit un faisceau parallèle $\to$ vision nette sans accommodation (œil relâché). L'image finale virtuelle $A''B''$ est \textbf{inversée par rapport à $AB$}.
\end{enumerate}

\textbf{Interprétation :} Le grossissement angulaire $G_c$ est le rapport de l'angle sous-tendu par l'image finale à l'infini ($\alpha'$) à l'angle sous-tendu par l'objet au point proche ($\alpha = AB/250$). On démontre $G_c = \frac{\Delta}{f_{obj}'} \times \frac{250}{f_{oc}'}$.
\end{experience}

\begin{popfigure}
\begin{tikzpicture}[>=Stealth, scale=0.8]
% Axes
\draw[thin, ->] (-6,0) -- (10,0) node[right] {Axe optique};
% Objectif
\draw[thick, popblue] (-1,-1.5) -- (-1,1.5) node[above] {$L_{obj}$};
\node[popblue] at (-1,-2) {$F_{obj}$};
\node[popblue] at (-1,2) {$F'_{obj}$};
% Oculaire
\draw[thick, poporange] (7,-1.5) -- (7,1.5) node[above] {$L_{oc}$};
\node[poporange] at (7,-2) {$F_{oc}$};
\node[poporange] at (7,2) {$F'_{oc}$};
% Object AB
\draw[thick, popdark] (-2.5,-0.5) -- (-2.5,0.5);
\node[popdark, left] at (-2.5,0.5) {$A$};
\node[popdark, left] at (-2.5,-0.5) {$B$};
% Rays from Object through Objective
\draw[popgreen, dashed] (-2.5,0.5) -- (-1,0.5) -- (3,-1); % parallel -> F'obj
\draw[popgreen, dashed] (-2.5,0.5) -- (-1,0) -- (3,-1); % center
\draw[popgreen, dashed] (-2.5,-0.5) -- (-1,-0.5) -- (3,1); % parallel -> F'obj
\draw[popgreen, dashed] (-2.5,-0.5) -- (-1,0) -- (3,1); % center
% Intermediate Image A'B'
\draw[thick, poppurple] (3,-1) -- (3,1);
\node[poppurple, right] at (3,1) {$A'$};
\node[poppurple, right] at (3,-1) {$B'$};
% Delta
\draw[<->, popmuted] (0, -2.5) -- (6, -2.5);
\node[popmuted, below] at (3, -2.5) {$\Delta$ (Longueur de tube)};
% Rays from Intermediate Image through Ocular (Normal adjustment: A'B' at F_oc)
% F_oc is at x = 7 - f_oc'. Let's say f_oc' = 2.5cm -> F_oc at x=4.5. But A'B' is at x=3.
% Let's adjust positions for correct geometry: Objective at 0. F'obj at 4mm. Delta=160mm. F_oc at 160mm. Ocular at 185mm.
% Simplified schematic:
% Objective at x=0. F'obj at x=4. Image at x=164 (Delta=160 from F'obj). F_oc at x=164. Ocular at x=189 (f_oc=25).
% Too wide for tikz. Schematic representation:
\draw[popgreen, dashed] (3,1) -- (7,1) -- (10,1); % parallel -> parallel
\draw[popgreen, dashed] (3,1) -- (7,0) -- (10, -0.8); % center
\draw[popgreen, dashed] (3,-1) -- (7,-1) -- (10,-1); % parallel -> parallel
\draw[popgreen, dashed] (3,-1) -- (7,0) -- (10, 0.8); % center
% Eye
\draw[poppink, thick] (10.5,0) circle (0.3);
\draw[poppink, thick] (10.3,-0.1) -- (10.1,-0.3);
\node[poppink, above] at (10.5,0.5) {Œil relâché};
% Final virtual image at infinity (arrows)
\draw[<->, poppink] (9, 1.2) -- (10, 1.2);
\node[poppink, above] at (9.5, 1.3) {$\alpha'$};
\end{tikzpicture}
\legende{Schéma de principe du microscope (réglage normal, image finale à l'infini). L'objectif forme l'image intermédiaire réelle $A'B'$ dans le plan focal objet de l'oculaire. L'oculaire renvoie des faisceaux parallèles.}
\end{popfigure}

\courssub{Grossissement commercial et Puissance commerciale}
\begin{propriete}[Grandeurs caractéristiques du microscope -- Réglage Normal (image à l'infini)]
Soit un microscope de \textbf{longueur de tube mécanique} $\Delta$ (distance entre le foyer image de l'objectif $F'_{obj}$ et le foyer objet de l'oculaire $F_{oc}$), normalisée à $\Delta = \SI{160}{\mm}$ (standard RMS).
\begin{align*}
    \text{Grossissement de l'objectif :}\quad & G_{obj} = \frac{\Delta}{f_{obj}'} \qquad (\Delta \text{ et } f_{obj}' \text{ en mm}) \\
    \text{Grossissement de l'oculaire :}\quad & G_{oc} = \frac{\SI{250}{\mm}}{f_{oc}'} \qquad (f_{oc}' \text{ en mm, image à l'infini}) \\
    \text{Grossissement commercial :}\quad & G_c = G_{obj} \times G_{oc} = \frac{\Delta \times \SI{250}{\mm}}{f_{obj}' f_{oc}'} \\
    \text{Puissance commerciale :}\quad & P_c = G_c \quad (\text{nombre pur, sans unité}) \\
    \text{Vergence de l'objectif :}\quad & V_{obj} = \frac{1}{f_{obj}'} = \frac{\num{1000}}{f_{obj}'(\text{mm})} \quad (\text{en } \delta = \si{\per\meter}) \\
    \text{Vergence de l'oculaire :}\quad & V_{oc} = \frac{1}{f_{oc}'} = \frac{\num{1000}}{f_{oc}'(\text{mm})} \quad (\text{en } \delta)
\end{align*}
\emph{Convention de marquage :} Un objectif marqué « $40\times / 0,65$ » a $G_{obj}=40$ (pour $\Delta=\SI{160}{\mm}$) et $NA=0,65$. Un oculaire marqué « $10\times$ » a $G_{oc}=10$ ($f_{oc}'=\SI{25,0}{\mm}$). Le marquage « $160/0,17$ » indique la correction pour $\Delta=\SI{160}{\mm}$ et épaisseur de couvre-objet $\SI{0,17}{\mm}$.
\end{propriete}

\begin{propriete}[Grossissement effectif -- Réglage Confort (image au point proche $D=\SI{250}{\mm}$)]
L'oculaire fonctionne comme une loupe utilisée au point proche : $G_{oc,eff} = 1 + \frac{D}{f_{oc}'} = 1 + \frac{\SI{250}{\mm}}{f_{oc}'(\text{mm})}$.
Le grossissement effectif total est :
\[ G_{eff} = G_{obj} \times G_{oc,eff} = \frac{\Delta}{f_{obj}'} \left(1 + \frac{\SI{250}{\mm}}{f_{oc}'}\right) = G_c \left(1 + \frac{1}{G_{oc}}\right) \]
\emph{Attention :} $G_{eff} > G_c$. Au Probatoire, \textbf{préciser toujours quel réglage est utilisé} (commercial = image à l'infini ; effectif = image à $\SI{250}{\mm}$).
\end{propriete}

\courssub{Mise au point : théorie et pratique}
\begin{methode}[Mise au point standard (Objectifs secs : $4\times, 10\times, 40\times$)]
\begin{enumerate}
    \item \textbf{Préparation :} Nettoyer les lentilles (objectifs, oculaire, condenseur) avec du papier optique. Vérifier l'épaisseur du couvre-objet ($\SI{0,17}{\mm}$). Poser la lame sur la platine, couvre-objet vers le haut. Centrer la zone d'observation sous l'objectif de plus faible grossissement ($4\times$ ou $10\times$).
    \item \textbf{Mise au point macrométrique (grosse vis) :} En regardant \textbf{de côté} (pas dans l'oculaire), descendre l'objectif (ou monter la platine) au plus près de la lame sans la toucher (distance $\approx \SI{1}{\mm}$).
    \item \textbf{Mise au point micrométrique (petite vis) :} Regarder dans l'oculaire. Tourner la micrométrique pour \textbf{monter} l'objectif (ou descendre la platine) jusqu'à obtenir une image nette. Ne jamais forcer vers le bas.
    \item \textbf{Changement d'objectif :} Basculer sur l'objectif supérieur ($10\times \to 40\times$). \textbf{Retoucher uniquement au micrométrique} (l'image est quasi-nette par conception parfocale). Re-centrer la zone d'intérêt.
    \item \textbf{Éclairage (Köhler approximatif) :} Fermer le diaphragme d'ouverture du condenseur jusqu'à ce que son image apparaisse nette dans le champ, puis l'ouvrir juste au bord du champ pour optimiser contraste et résolution.
\end{enumerate}
\end{methode}

\begin{methode}[Mise au point à l'immersion (Objectif $100\times$ Huile -- Indice $n=\num{1,515}$)]
\begin{enumerate}
    \item Effectuer la mise au point complète sur l'objectif $40\times$ (étapes 1 à 4 ci-dessus). L'image est nette et centrée.
    \item \textbf{Ne pas basculer encore sur $100\times$.} Déposer \textbf{une seule goutte} d'huile d'immersion ($n=\num{1,515}$) sur le centre du couvre-objet, exactement sur la zone éclairée.
    \item Basculer \textbf{lentement} l'objectif $100\times$ jusqu'au clic. La lentille frontale doit tremper dans la goutte d'huile.
    \item \textbf{Interdiction formelle d'utiliser la macrométrique.} Mise au point \textbf{exclusivement au micrométrique} (débattement très court $\approx \SI{0,1}{\mm}$). Tourner très lentement. Si pas d'image après un tour complet : remonter légèrement, vérifier l'absence de bulle d'air, redescendre.
    \item \textbf{Fin de séance (CRITIQUE) :} Remonter l'objectif $100\times$ (ou baisser la platine) \textbf{AVANT} de basculer sur un autre objectif. Essuyer \textbf{immédiatement} l'huile sur la lentille frontale de l'objectif $100\times$ et sur le couvre-objet avec du papier optique imbibé de xylène ou d'alcool isopropylique $\ge 99\%$. Ne jamais laisser l'huile sécher (durcissement, décollement des lentilles, perte de $NA$).
\end{enumerate}
\end{methode}

\courssub{Longueur du tube, encombrement et limites (Complément Probatoire)}
\begin{propriete}[Longueur de tube $\Delta$ et encombrement]
La distance optique $\Delta = \overline{F'_{obj}F_{oc}}$ est fixée par le constructeur ($\SI{160}{\mm}$ standard RMS, $\SI{170}{\mm}$ ou $\infty$ pour les microscopes à l'infini modernes). L'encombrement mécanique $L$ (distance objectif-oculaire) vaut $L \approx \Delta + f_{obj}' + f_{oc}'$.
\begin{itemize}
    \item \textbf{Objectif $100\times$ huile :} $f_{obj}' = \SI{1,6}{\mm}$. $NA = n \sin \alpha \le n = \num{1,515}$ (théorique max $\approx \num{1,4}$). Sans huile ($n=1$), $NA_{eff} \approx 1,0$ $\to$ résolution divisée par $1,4$ et aberrations sphériques majeures (image floue, halo).
    \item \textbf{Grossissement vide (inutile) :} Augmenter $G_c$ en changeant l'oculaire (ex: $25\times$) sans augmenter $NA$ ne révèle aucun détail nouveau. Limite utile : $G_c \approx 500 \times NA$ à $1000 \times NA$.
    \item \textbf{Champ de vision :} Diamètre champ objet $= \frac{\text{Champ oculaire (mm)}}{G_{obj}}$. Un oculaire $15\times$ a un champ apparent plus petit qu'un $10\times$ $\to$ champ objet réduit.
    \item \textbf{Luminosité :} Éclairement image $\propto \left(\frac{NA}{G_c}\right)^2$. Doubler $G_c$ divise la luminosité par 4.
\end{itemize}
\end{propriete}

\courssub{Résolution de problèmes types (Probatoire) -- Activité d'intégration}

\courssub{Objectifs de l'activité}
Cette activité d'intégration vise à mobiliser l'ensemble des compétences acquises dans la leçon sur le microscope (description, schéma, calculs de grossissement et de puissance, mise au point, pouvoir séparateur) pour résoudre une situation professionnelle authentique : la vérification de conformité du matériel d'un lycée camerounais avant une session d'examen pratique du Probatoire. À l'issue de cette activité, vous devez être capable de :
\begin{itemize}
    \item Calculer les distances focales, puissances vergences et grossissements (normal et confort) pour toutes les combinaisons objectifs/oculaires données.
    \item Exploiter l'ouverture numérique pour déterminer le pouvoir séparateur et juger de l'aptitude à observer des bactéries.
    \item Détecter une non-conformité de marquage (anomalie) et en quantifier l'impact sur le grossissement réel.
    \item Rédiger une procédure rigoureuse de mise au point à l'immersion (objectif $100\times$ huile).
    \item Structurer un rapport technique clair, complet et argumenté.
\end{itemize}

\courssub{Situation}
\emph{Contexte :} Nous sommes au Lycée Bilingue d'Etoug-Ebe (Yaoundé), une semaine avant le début des épreuves pratiques du Probatoire, session de juin. M. \textsc{Mvondo}, le chef des travaux de SVT, doit valider les deux microscopes optiques du laboratoire qui serviront aux candidats de la série C pour l'observation de lames de sang (recherche de plasmodium du paludisme) et de cultures bactériennes (identification de bacilles). Les fiches techniques constructeurs sont égarées ; seuls les marquages gravés sur les objectifs et les oculaires sont lisibles. Il constitue des groupes de 3-4 élèves « experts en optique » pour auditer le matériel.

\medskip
\noindent \textbf{Matériel à auditer (données relevées sur les instruments) :}
\begin{center}
\begin{tabularx}{\linewidth}{|l|X|X|}\hline
\textbf{Composant} & \textbf{Microscope A (Marque \emph{Optika} B-150)} & \textbf{Microscope B (Marque \emph{Euromex} Oxion)} \\ \hline
Objectifs disponibles & \begin{itemize}\item Obj 1 : $10\times / 0,25$ (marqué $160/0,17$)\item Obj 2 : $40\times / 0,65$ (marqué $160/0,17$)\end{itemize} & \begin{itemize}\item Obj 1 : $4\times / 0,10$ (marqué $160/0,17$)\item Obj 2 : $40\times / 0,65$ (marqué $160/0,17$)\item Obj 3 : $100\times / 1,25$ \textbf{Huile} (marqué $160/0,17$)\end{itemize} \\ \hline
Oculaires disponibles & Oculaire unique : \textbf{$10\times$} (marqué $f' = \SI{25,0}{\mm}$) & Deux oculaires : \textbf{$10\times$} ($f' = \SI{25,0}{\mm}$) et \textbf{$15\times$} ($f' = \SI{16,7}{\mm}$) \\ \hline
Longueur de tube mécanique $\Delta$ & \multicolumn{2}{c|}{$\Delta = \SI{160}{\mm}$ (standard RMS, distance plan image objectif -- plan objet oculaire)} \\ \hline
\end{tabularx}
\end{center}

\medskip
\noindent \textbf{Anomalie suspectée sur Microscope A :} Sur l'objectif $40\times$, le marquage « $160/0,17$ » indique une correction pour une longueur de tube de $\SI{160}{\mm}$ et une lame de couvre-objet d'épaisseur $\SI{0,17}{\mm}$. Or, l'oculaire présent est un $15\times$ (focale $\SI{16,7}{\mm}$) au lieu du $10\times$ standard prévu pour ce grossissement commercial affiché.

\medskip
\noindent \textbf{Contraintes de l'examen :}
\begin{itemize}
    \item Observation de bactéries (bacilles) de taille typique $\SI{1}{\um}$ à $\SI{2}{\um}$.
    \item Longueur d'onde de référence pour la résolution : $\lambda = \SI{550}{\nm}$ (lumière verte, pic de sensibilité de l'œil).
    \item Distance de vision distincte (confort) : $D = \SI{250}{\mm}$.
    \item Réglage « normal » : image finale à l'infini (œil relâché). Réglage « confort » : image finale à $\SI{250}{\mm}$ (point proche).
\end{itemize}

\courssub{Consignes}
Travailler en groupe et produire un rapport unique structuré répondant aux cinq points suivants :

\begin{enumerate}
    \item \textbf{Tableau récapitulatif des performances optiques.} Pour chaque microscope et chaque combinaison (Objectif + Oculaire) utilisable, calculer :
    \begin{itemize}
        \item La distance focale image de l'objectif $f_{obj}'$ (en mm) et sa vergence $V_{obj}$ (en dioptries).
        \item La distance focale image de l'oculaire $f_{oc}'$ (en mm) et sa vergence $V_{oc}$ (en dioptries).
        \item Le grossissement commercial $G_c$ (réglage normal, image à l'infini).
        \item Le grossissement effectif $G_{eff}$ (réglage confort, image à $\SI{250}{\mm}$).
        \item La puissance commerciale $P_c$ (rappel : $P_c = G_c$ pour un microscope, nombre pur).
    \end{itemize}
    Présenter les résultats dans un grand tableau comparatif.

    \item \textbf{Pouvoir séparateur et choix pour la bactériologie.} 
    \begin{itemize}
        \item Rappeler la formule du pouvoir séparateur (critère de Rayleigh) $\delta$ en fonction de l'ouverture numérique $NA$ et de $\lambda$.
        \item Calculer $\delta$ pour chaque objectif des deux microscopes.
        \item Déterminer quel(s) microscope(s) et quel(s) objectif(s) permettent de \textbf{distinguer} deux bactéries adjacentes de $\SI{1,5}{\um}$ (critère : $\delta < \text{taille}$).
        \item Justifier le choix final de l'instrument pour l'épreuve de bactériologie.
    \end{itemize}

    \item \textbf{Analyse de l'anomalie du Microscope A.} 
    \begin{itemize}
        \item L'objectif $40\times$ est conçu pour donner un grossissement objectif $G_{obj} = 40$ avec $\Delta = \SI{160}{\mm}$. Vérifier sa focale théorique.
        \item L'oculaire monté est un $15\times$ ($f_{oc}' = \SI{16,7}{\mm}$) au lieu du $10\times$ standard. Calculer le \textbf{vrai grossissement commercial} $G_{c,\text{réel}}$ obtenu avec cette combinaison.
        \item Commenter l'impact sur le champ de vision, la luminosité et la validité de l'échelle de mesure (micromètre oculaire) pour les candidats.
    \end{itemize}

    \item \textbf{Procédure de mise au point : Objectif $100\times$ Huile (Microscope B).} Rédiger la procédure pas à pas, en langage impératif, en précisant :
    \begin{itemize}
        \item La préparation de la lame (couvre-objet $\SI{0,17}{\mm}$, immersion).
        \item La séquence de grossissements (bas $\to$ haut) et les boutons à actionner (macrométrique / micrométrique).
        \item La gestion de l'huile d'immersion (indice $n = \num{1,515}$) : quantité, mise en place, nettoyage final.
        \item Les consignes de sécurité (objectif fragile, huile sur les autres objectifs interdite).
    \end{itemize}

    \item \textbf{Synthèse et Recommandation.} Rédiger une conclusion formelle (style rapport technique) indiquant : 
    \begin{itemize}
        \item Le microscope retenu pour l'examen (A ou B ?).
        \item La configuration exacte (Objectif + Oculaire) pour l'épreuve de parasitologie (paludisme) et pour la bactériologie.
        \item La décision concernant l'anomalie du Microscope A (réparation, étiquetage, mise au rebut ?).
    \end{itemize}
\end{enumerate}

\courssub{Ressources à mobiliser}
\begin{propriete}[Formules du microscope -- Réglage Normal (image à l'infini)]
Pour un microscope de longueur de tube $\Delta$ (distance foyer image objectif -- foyer objet oculaire) :
\begin{align*}
    \text{Grossissement objectif :}\quad & G_{obj} = \frac{\Delta}{f_{obj}'} \quad (\Delta \text{ et } f_{obj}' \text{ en mm}) \\
    \text{Grossissement oculaire :}\quad & G_{oc} = \frac{250}{f_{oc}'} \quad (f_{oc}' \text{ en mm}, \text{ image à l'infini}) \\
    \text{Grossissement commercial :}\quad & G_c = G_{obj} \times G_{oc} = \frac{\Delta \times 250}{f_{obj}' f_{oc}'} \\
    \text{Puissance commerciale :}\quad & P_c = G_c \quad (\text{adimensionnel}) \\
    \text{Vergence objectif :}\quad & V_{obj} = \frac{1000}{f_{obj}'(\text{mm})} \quad (\delta) \\
    \text{Vergence oculaire :}\quad & V_{oc} = \frac{1000}{f_{oc}'(\text{mm})} \quad (\delta)
\end{align*}
\emph{Note :} Le marquage « $10\times$ » sur l'objectif signifie $G_{obj} = 10$ pour $\Delta = \SI{160}{\mm}$. Le marquage « $10\times$ » sur l'oculaire signifie $G_{oc} = 10$ (soit $f_{oc}' = \SI{25}{\mm}$).
\end{propriete}

\begin{propriete}[Grossissement effectif -- Réglage Confort (image au point proche $D=\SI{250}{\mm}$)]
L'oculaire fonctionne comme une loupe : $G_{oc,eff} = 1 + \frac{D}{f_{oc}'} = 1 + \frac{250}{f_{oc}'(\text{mm})}$.
Le grossissement effectif total est :
\[ G_{eff} = G_{obj} \times G_{oc,eff} = \frac{\Delta}{f_{obj}'} \left(1 + \frac{250}{f_{oc}'}\right) \]
\emph{Attention :} $G_{eff} > G_c$. Au Probatoire, préciser toujours quel réglage est utilisé.
\end{propriete}

\begin{propriete}[Pouvoir séparateur (Résolution) -- Critère de Rayleigh]
Pour un objectif d'ouverture numérique $NA = n \sin \alpha$ (air : $n=1$ ; huile : $n=\num{1,515}$), la distance minimale résolue entre deux points est :
\[ \delta = 0,61 \frac{\lambda}{NA} \]
avec $\lambda$ en $\text{nm}$ (ou $\mu\text{m}$) et $\delta$ dans la même unité.
Deux objets sont distingués si leur distance centre-à-centre $> \delta$.
\end{propriete}

\begin{exoresolu}[Rapport d'expertise complet -- Microscopes A et B]
\textbf{Énoncé de l'activité :} Voir consignes ci-dessus (Situation, Matériel, Anomalie, Contraintes, 5 points à traiter).
\tcblower
\textbf{Solution détaillée}

\medskip
\noindent \textbf{Étape 1 : Calcul des caractéristiques intrinsèques (Focales, Vergences).}
\medskip
\noindent \underline{Rappel :} Marquage Objectif $G_{obj}\times / NA$ $\Rightarrow$ $f_{obj}' = \frac{\Delta}{G_{obj}}$ avec $\Delta = \SI{160}{\mm}$.
Marquage Oculaire $G_{oc}\times$ $\Rightarrow$ $f_{oc}' = \frac{250}{G_{oc}}$ (mm).

\begin{center}
\begin{tabular}{|l|ccc|ccc|}\hline
 & \multicolumn{3}{c|}{\textbf{Microscope A}} & \multicolumn{3}{c|}{\textbf{Microscope B}} \\ \hline
\textbf{Paramètre} & \textbf{Obj $10\times$} & \textbf{Obj $40\times$} & \textbf{Ocul $10\times$} & \textbf{Obj $4\times$} & \textbf{Obj $40\times$} & \textbf{Obj $100\times$ Huile} \\ \hline
$G_{obj}$ marqué & 10 & 40 & -- & 4 & 40 & 100 \\
$f_{obj}'$ (mm) & $\frac{160}{10}=16,0$ & $\frac{160}{40}=4,00$ & -- & $\frac{160}{4}=40,0$ & $4,00$ & $\frac{160}{100}=1,60$ \\
$V_{obj}$ ($\delta$) & $\frac{1000}{16}=62,5$ & $\frac{1000}{4}=250$ & -- & $\frac{1000}{40}=25,0$ & $250$ & $\frac{1000}{1,6}=625$ \\
$NA$ & 0,25 & 0,65 & -- & 0,10 & 0,65 & 1,25 \\ \hline
$G_{oc}$ marqué & \multicolumn{3}{c|}{10} & \multicolumn{3}{c|}{10 \quad / \quad 15} \\ \hline
$f_{oc}'$ (mm) & \multicolumn{3}{c|}{$\frac{250}{10}=25,0$} & \multicolumn{3}{c|}{25,0 \quad / \quad $\frac{250}{15}=16,7$} \\
$V_{oc}$ ($\delta$) & \multicolumn{3}{c|}{$\frac{1000}{25}=40,0$} & \multicolumn{3}{c|}{40,0 \quad / \quad $\frac{1000}{16,7}=60,0$} \\ \hline
\end{tabular}
\end{center}

\medskip
\noindent \textbf{Étape 2 : Grossissements Commercial ($G_c$) et Effectif ($G_{eff}$).}
\medskip
\noindent Formules : $G_c = G_{obj} \times G_{oc}$ ; $G_{eff} = G_{obj} \times \left(1 + \frac{250}{f_{oc}'}\right) = G_{obj} \times (1 + G_{oc})$.

\begin{center}
\begin{tabular}{|l|cc|cccc|}\hline
\textbf{Combinaison} & \multicolumn{2}{c|}{\textbf{Microscope A}} & \multicolumn{4}{c|}{\textbf{Microscope B}} \\ \hline
 & \textbf{Obj $10\times$ + Ocul $10\times$} & \textbf{Obj $40\times$ + Ocul $10\times$} & \textbf{Obj $4\times$ + Ocul $10\times$} & \textbf{Obj $40\times$ + Ocul $10\times$} & \textbf{Obj $40\times$ + Ocul $15\times$} & \textbf{Obj $100\times$ + Ocul $10\times$} \\ \hline
$G_{obj}$ & 10 & 40 & 4 & 40 & 40 & 100 \\
$G_{oc}$ (normal) & 10 & 10 & 10 & 10 & 15 & 10 \\
\textbf{$G_c = G_{obj}G_{oc}$} & \textbf{100} & \textbf{400} & \textbf{40} & \textbf{400} & \textbf{600} & \textbf{1000} \\
$G_{oc,eff}$ (confort) & $1+10=11$ & $11$ & $11$ & $11$ & $1+15=16$ & $11$ \\
\textbf{$G_{eff} = G_{obj}G_{oc,eff}$} & \textbf{110} & \textbf{440} & \textbf{44} & \textbf{440} & \textbf{640} & \textbf{1100} \\
\textbf{Puissance $P_c = G_c$} & 100 & 400 & 40 & 400 & 600 & 1000 \\ \hline
\end{tabular}
\end{center}

\medskip
\noindent \textbf{Étape 3 : Pouvoir séparateur $\delta$ et choix bactériologie.}
\medskip
\noindent $\lambda = \SI{550}{\nm} = \SI{0,55}{\um}$. Formule : $\delta = 0,61 \frac{\lambda}{NA}$.
\begin{align*}
\text{Micro A -- Obj }10\times (NA=0,25) &: \delta = 0,61 \times \frac{0,55}{0,25} = \SI{1,34}{\um} \\
\text{Micro A -- Obj }40\times (NA=0,65) &: \delta = 0,61 \times \frac{0,55}{0,65} = \SI{0,52}{\um} \\
\text{Micro B -- Obj }4\times (NA=0,10) &: \delta = 0,61 \times \frac{0,55}{0,10} = \SI{3,36}{\um} \\
\text{Micro B -- Obj }40\times (NA=0,65) &: \delta = \SI{0,52}{\um} \quad (\text{identique Micro A}) \\
\text{Micro B -- Obj }100\times \text{ Huile } (NA=1,25) &: \delta = 0,61 \times \frac{0,55}{1,25} = \SI{0,27}{\um}
\end{align*}

\medskip
\noindent \textbf{Analyse :} Taille bactéries $\approx \SI{1,5}{\um}$.
\begin{itemize}
    \item Micro A Obj $10\times$ : $\delta = \SI{1,34}{\um} < \SI{1,5}{\um}$ \textbf{? Limite.} Théoriquement résolu, mais contraste faible, limite de l'œil.
    \item Micro A Obj $40\times$ : $\delta = \SI{0,52}{\um} \ll \SI{1,5}{\um}$ \textbf{Excellent.}
    \item Micro B Obj $4\times$ : $\delta = \SI{3,36}{\um} > \SI{1,5}{\um}$ \textbf{Impossible.}
    \item Micro B Obj $40\times$ : $\delta = \SI{0,52}{\um}$ \textbf{Excellent.}
    \item Micro B Obj $100\times$ Huile : $\delta = \SI{0,27}{\um}$ \textbf{Optimal (détails internes).}
\end{itemize}
\emph{Conclusion :} Les deux microscopes permettent l'observation bactériologique avec leur objectif $40\times$ (air). Le Microscope B offre en plus l'objectif $100\times$ huile pour une résolution supérieure (morphologie fine, flagelles). Le Microscope A n'a pas d'objectif $100\times$.

\medskip
\noindent \textbf{Étape 4 : Anomalie Microscope A (Objectif $40\times$ + Oculaire $15\times$).}
\medskip
\noindent L'objectif $40\times$ est calibré pour $\Delta = \SI{160}{\mm}$ $\Rightarrow$ $f_{obj}' = \SI{4,00}{\mm}$ (vérifié).
L'oculaire présent est $15\times$ ($f_{oc}' = \SI{16,7}{\mm}$) au lieu de $10\times$ ($\SI{25,0}{\mm}$).
\medskip
\noindent \textbf{Vrai grossissement commercial :} $G_{c,\text{réel}} = G_{obj} \times G_{oc,\text{réel}} = 40 \times 15 = \mathbf{600\times}$.
\medskip
\noindent \textbf{Commentaires critiques :}
\begin{itemize}
    \item \textbf{Sur-grossissement de 50\%} par rapport au marquage « $400\times$ » implicite (objectif $40\times$ + oculaire standard $10\times$). Les candidats mesurant une cellule avec un micromètre oculaire étalonné pour $400\times$ trouveront une taille 1,5 fois trop grande $\rightarrow$ \textbf{Erreur systématique majeure}.
    \item \textbf{Champ de vision réduit :} Le champ réel objet $= \frac{\text{Champ oculaire}}{G_{obj}}$. L'oculaire $15\times$ a un champ apparent plus petit qu'un $10\times$ (souvent $\SI{16}{\mm}$ vs $\SI{18-20}{\mm}$). Le champ objet devient très petit, difficile pour localiser les parasites.
    \item \textbf{Luminosité :} Éclairement $\propto \left(\frac{NA}{G_c}\right)^2$. $G_c$ passe de 400 à 600 $\Rightarrow$ luminosité divisée par $(600/400)^2 = 2,25$. Image plus sombre, inconfortable.
    \item \textbf{Correction optique :} L'objectif $40\times$ (marqué $160/0,17$) est corrigé pour un tube de $\SI{160}{\mm}$. L'oculaire $15\times$ a une distance focale plus courte ($\SI{16,7}{\mm}$ vs $\SI{25}{\mm}$). Son plan objet principal est décalé. L'image intermédiaire formée par l'objectif ne tombe pas exactement dans le plan objet de l'oculaire $\rightarrow$ \textbf{aberration sphérique et chromatique résiduelles non corrigées}, perte de piqué aux bords.
\end{itemize}
\emph{Verdict :} Microscope A \textbf{non conforme} pour l'examen en configuration $40\times$. Il faut impérativement remettre l'oculaire $10\times$ d'origine ou étiqueter « Configuration modifiée : $G=600\times$ ».

\medskip
\noindent \textbf{Étape 5 : Procédure mise au point Objectif $100\times$ Huile (Micro B) -- Version Opératoire.}
\begin{enumerate}
    \item \textbf{Préparation lame :} Vérifier épaisseur couvre-objet $\SI{0,17}{\mm}$ (marquage sur l'objectif). Nettoyer lame et couvre-objet (alcool 70\%, essuyer papier optique). Déposer l'échantillon (frottis sang ou culture), couvrir.
    \item \textbf{Mise en place :} Poser la lame sur la platine, couvre-objet vers le haut. Bloquer avec les pinces. Centrer grossièrement la zone au-dessus du condenseur (diaphragme iris ouvert au max, lentille condenseur rentrée).
    \item \textbf{Objectif $4\times$ :} Le sélectionner (tourelle). Baisser la platine (ou monter le tube) au macrométrique jusqu'à approche maximale sans contact (surveiller de côté, distance $\approx \SI{1}{\cm}$). Regarder dans l'oculaire $10\times$ : monter la platine au macrométrique jusqu'à apparition de l'image. Net au micrométrique. Centrer la zone d'intérêt au centre du champ.
    \item \textbf{Objectif $10\times$ :} Basculer. Retoucher \textbf{uniquement au micrométrique}. Re-centrer. Fermer diaphragme iris condenseur jusqu'au bord du champ (Köhler approximatif) pour le contraste.
    \item \textbf{Objectif $40\times$ :} Basculer. \textbf{Uniquement micrométrique}. Image nette. Re-centrer.
    \item \textbf{Immersion :} \textbf{Ne pas basculer encore sur $100\times$.} Déposer \textbf{une seule goutte} d'huile d'immersion ($n=\num{1,515}$, flacon compte-gouttes) sur le centre du couvre-objet, exactement sur la zone éclairée.
    \item \textbf{Objectif $100\times$ :} Basculer \textbf{lentement} jusqu'au clic. L'avant de l'objectif doit tremper dans la goutte d'huile. \textbf{Interdiction formelle d'utiliser le macrométrique.}
    \item \textbf{Mise au point fine :} Tourner le micrométrique \textbf{très lentement} (débattement $\approx \SI{0,1-0,2}{\mm}$) jusqu'à l'image nette. Si pas d'image après 1 tour complet : \textbf{remonter légèrement}, vérifier qu'il y a de l'huile (pas de bulle), redescendre.
    \item \textbf{Observation :} Ajuster diaphragme iris et intensité lampe. Noter grossissement $G_c = 1000\times$ ($G_{eff}=1100\times$).
    \item \textbf{Fin de séance :} Remonter l'objectif $100\times$ (ou baisser platine) \textbf{avant} de basculer sur un autre objectif. Essuyer \textbf{immédiatement} l'huile sur la lentille frontale de l'objectif $100\times$ avec papier optique + xylène (ou alcool isopropylique 99\%). Essuyer l'huile sur le couvre-objet. Ne jamais laisser l'huile sécher (durcissement, perte de collage lentilles). Nettoyer aussi l'objectif $40\times$ s'il a été touché par capillarité.
\end{enumerate}

\medskip
\noindent \textbf{Étape 6 : Synthèse et Recommandation finale.}
\medskip
\noindent \textbf{Rapport technique -- Validation Matériel Probatoire Session Juin -- Lycée Bilingue Etoug-Ebe}

\begin{aretenir}[Tableau de décision final]
\begin{center}
\begin{tabularx}{\linewidth}{|l|X|X|}\hline
\textbf{Épreuve} & \textbf{Configuration Recommandée} & \textbf{Justification} \\ \hline
Parasitologie (Paludisme) & \textbf{Microscope A -- Obj $10\times$ + Ocul $10\times$ ($G_c=100\times$)} & Champ large ($\SI{1,8-2}{\mm}$), luminosité max, balayage rapide frottis sanguin. Grossissement suffisant pour plasmodium ($\SI{2-5}{\um}$). \textbf{Condition :} Oculaire $10\times$ d'origine remis en place. \\ \hline
Bactériologie (Bacilles) & \textbf{Microscope B -- Obj $40\times$ + Ocul $10\times$ ($G_c=400\times$)} & $\delta = \SI{0,52}{\um} \ll \SI{1,5}{\um}$. Résolution confortable. Oculaire $10\times$ standard. \\ \hline
Bactériologie (Morphologie fine) & \textbf{Microscope B -- Obj $100\times$ Huile + Ocul $10\times$ ($G_c=1000\times$)} & $\delta = \SI{0,27}{\um}$. Résolution optimale pour flagelles, parois, inclusion. Procédure immersion stricte. \\ \hline
\end{tabularx}
\end{center}
\end{aretenir}

\medskip
\noindent \textbf{Décision sur l'anomalie Microscope A :}
\begin{itemize}
    \item \textbf{Action immédiate :} Retirer l'oculaire $15\times$ (le ranger en réserve / rechange). Remettre l'oculaire $10\times$ standard (commander si manquant).
    \item \textbf{Étiquetage :} Poser une étiquette « \emph{Config. Examen : Ocul 10× obligatoire} » sur le corps du microscope.
    \item \textbf{Ne pas utiliser} la combinaison $40\times/15\times$ pour l'examen (erreur de mesure, aberrations, luminosité).
    \item \textbf{Microscope B} est l'instrument principal (polyvalent $4\times/40\times/100\times$ huile, deux oculaires conformes). Le réserver aux épreuves exigeant la haute résolution (Bactériologie $100\times$).
\end{itemize}

\medskip
\noindent \textbf{Compétences validées :} Calculs $f', V, G_c, G_{eff}, \delta$ ; Analyse dimensionnelle (mm, $\mu$m, $\delta$) ; Lecture marquages normalisés ($160/0,17$) ; Rédaction procédure technique ; Esprit critique (détection anomalie, impact métrologique).
\end{exoresolu}

\courssub{Bilan de l'activité}
Cette activité d'intégration a permis de mettre en évidence les points clés suivants, essentiels pour la réussite au Probatoire :

\begin{aretenir}[Bilan des acquis mobilisés]
\begin{itemize}
\item \textbf{Modélisation complète :} On passe du marquage constructeur ($G\times/NA$, $160/0,17$) aux grandeurs physiques réelles ($f', V, \delta$) par l'application rigoureuse des formelles $G_{obj}=\Delta/f_{obj}'$, $G_{oc}=250/f_{oc}'$, $\delta=0,61\lambda/NA$.
\item \textbf{Distinction cruciale :} \textbf{Grossissement Commercial $G_c$ (image à l'infini, référence constructeur)} vs \textbf{Grossissement Effectif $G_{eff}$ (image à 250 mm, confort visuel réel)}. $G_{eff} = G_c \times (1 + 1/G_{oc}) \approx 1,1 G_c$ pour un oculaire $10\times$. Ne jamais les confondre dans un énoncé.
\item \textbf{Pouvoir séparateur :} C'est l'ouverture numérique $NA$ (et non le grossissement) qui détermine la résolution. Un grossissement $1000\times$ avec $NA=0,65$ (objectif $40\times$ + oculaire $25\times$ « vide ») ne résout pas mieux que $400\times$ ($NA=0,65$). L'objectif $100\times$ huile ($NA=1,25$) est le seul à franchir la barre des $\SI{0,3}{\um}$.
\item \textbf{Rigueur métrologique :} L'anomalie du Microscope A illustre qu'un instrument d'optique est un \textbf{système cohérent}. Changer un composant (oculaire) sans recalculer/réétiqueter invalide les mesures (micrométrie) et dégrade l'image (aberrations, luminosité). Compétence « Savoir-être » : rigueur, traçabilité, sécurité.
\item \textbf{Compétence expérimentale :} La procédure de mise au point à l'immersion est un savoir-faire technique évalué en pratique. La séquence « Bas grossissement $\to$ Centrage $\to$ Montée progressive $\to$ Huile $\to$ Micrométrique seul $\to$ Nettoyage immédiat » est un standard international (ISO 8037) à maîtriser parfaitement.
\end{itemize}
\end{aretenir}

\begin{attention}[Pièges Probatoire identifiés dans l'activité]
\begin{itemize}
\item \textbf{Unités :} $\Delta = \SI{160}{\mm}$, $f'$ en $\text{mm}$, $\lambda$ en $\text{nm}$ ou $\mu\text{m}$, $\delta$ en $\mu\text{m}$. Convertir \textbf{avant} de calculer. $1\,\text{mm} = 10^3\,\mu\text{m} = 10^6\,\text{nm}$.
\item \textbf{Formule $G_{oc}$ :} $G_{oc} = 250/f_{oc}'(\text{mm})$ pour le \textbf{commercial} (image $\infty$). $G_{oc,eff} = 1 + 250/f_{oc}'(\text{mm})$ pour le \textbf{confort} (image $\SI{250}{\mm}$). Lire l'énoncé : « Grossissement commercial » $\Rightarrow$ formule sans le $1+$.
\item \textbf{Puissance commerciale :} $P_c = G_c$ (nombre pur). Ne pas confondre avec la vergence $V = 1/f'$ (en $\delta = \si{\per\meter}$).
\item \textbf{Marquage $160/0,17$ :} Indique la distance tube mécanique ($\Delta$) et l'épaisseur du couvre-objet pour laquelle l'objectif est corrigé (aberration sphérique). Utiliser un couvre-objet $\SI{1,2}{\mm}$ (lame épaisse) dégrade l'image aux forts grossissements ($40\times, 100\times$).
\item \textbf{Immersion :} L'huile ($n=\num{1,515}$) supprime l'interface air/verre ($n=1$), augmentant $NA = n \sin \alpha$ au max $\approx 1,4$. Sans huile, l'objectif $100\times$ a un $NA$ effectif $\approx 1,0$ (air) $\rightarrow$ résolution divisée par 1,4, et aberrations sphériques massives (image floue, halo).
\end{itemize}
\end{attention}

\newpage

\coursec{Méthodes et exercices résolus}

\courssub{Méthodes et exercices résolus}

\begin{methode}[Construire le schéma de principe et expliquer le fonctionnement]
\textbf{Objectif :} Représenter le trajet de la lumière dans un microscope composé et justifier la formation de l'image finale.
\begin{enumerate}
    \item \textbf{Placer les lentilles :} Tracer l'axe optique. Placer l'objectif (convergent, focale $f_O$ courte) et l'oculaire (convergent, focale $f_{O'}$ plus grande) séparés par la distance $\overline{OO'} = \Delta + f_O + f_{O'}$ ($\Delta$ = intervalle optique).
    \item \textbf{Placer l'objet :} L'objet $AB$ (réel, inversé par convention) est placé entre $F_O$ et $2F_O$ (légèrement au-delà du foyer objet de l'objectif) : $f_O < \overline{OA} < 2f_O$.
    \item \textbf{Image intermédiaire par l'objectif :} Construire l'image $A'B'$ de $AB$ par l'objectif (réelle, inversée, agrandie). Elle se forme au plan focal image de l'oculaire $F_{O'}$ (pour une observation \emph{au repos}, image finale à l'infini) ou légèrement avant (pour une observation \emph{au point proche}).
    \item \textbf{Image finale par l'oculaire :} L'oculaire fonctionne comme une loupe. L'image intermédiaire $A'B'$ sert d'objet à l'oculaire. Construire l'image finale $A''B''$ (virtuelle, droite par rapport à $A'B'$, donc inversée par rapport à $AB$).
    \item \textbf{Conclusion :} Le microscope donne une image finale \textbf{virtuelle, inversée} par rapport à l'objet initial et \textbf{très agrandie}.
\end{enumerate}
\cle{Schéma de principe} \cle{Image intermédiaire} \cle{Image finale virtuelle} \cle{Inversée}
\end{methode}

\begin{exoresolu}[Schéma de principe et fonctionnement (Type Probatoire 2021)]
\textbf{Énoncé :} On réalise le schéma de principe d'un microscope composé d'un objectif $L_O$ de focale $f_O = \SI{4,0}{mm}$ et d'un oculaire $L_{O'}$ de focale $f_{O'} = \SI{25}{mm}$. La distance entre les lentilles est $\overline{OO'} = \SI{160}{mm}$. L'objet $AB$ de hauteur $\SI{0,10}{mm}$ est placé à $\SI{4,5}{mm}$ de l'objectif. L'observation se fait \emph{au repos} (œil accommodé sur l'infini).
\begin{enumerate}
    \item Construire le schéma de principe à l'échelle (échelle suggérée : \SI{1}{cm} pour \SI{10}{mm} sur l'axe optique, \SI{1}{cm} pour \SI{0,1}{mm} en hauteur). Tracer les rayons particuliers.
    \item Qualifier l'image intermédiaire $A'B'$ (réelle/virtuelle, droite/inversée, grossissement $\gamma_O$).
    \item Qualifier l'image finale $A''B''$ (réelle/virtuelle, droite/inversée par rapport à $AB$).
    \item Calculer le grossissement commercial $G_c$ de cet instrument.
\end{enumerate}

\tcblower
\textbf{Solution détaillée :}
\begin{enumerate}
    \item \textbf{Construction géométrique (description) :}
    \begin{itemize}
        \item Axe optique horizontal. Origine $O$ (objectif). $O'$ à $\SI{16}{cm}$ (échelle \SI{1}{cm}/\SI{10}{mm}).
        \item Foyers : $F_O$ à $\SI{4,0}{mm}$ ($\SI{0,4}{cm}$) à droite de $O$. $F_{O'}$ à $\SI{25}{mm}$ ($\SI{2,5}{cm}$) à gauche de $O'$.
        \item Objet $AB$ : $A$ sur l'axe, $B$ à $\SI{0,10}{mm}$ ($\SI{1,0}{cm}$ à l'échelle hauteur) au-dessus. Distance $OA = \SI{4,5}{mm}$ ($\SI{0,45}{cm}$). $\rightarrow$ $AB$ est entre $F_O$ ($\SI{4,0}{mm}$) et $2F_O$ ($\SI{8,0}{mm}$). Correct.
        \item \textbf{Objectif $L_O$ :} Rayon $B \to$ parallèle à l'axe $\to$ réfracté par $F'_O$. Rayon $B \to O$ (centre optique) $\to$ non dévié. Intersection $\to B'$. $A'$ sur l'axe.
        \item Position $A'$ : Relation de conjugaison $\frac{1}{f_O} = \frac{1}{OA} + \frac{1}{OA'}$. $\frac{1}{4} = \frac{1}{4,5} + \frac{1}{OA'} \Rightarrow \frac{1}{OA'} = \frac{1}{4} - \frac{1}{4,5} = \frac{4,5-4}{18} = \frac{0,5}{18} = \frac{1}{36}$. Donc $OA' = \SI{36}{mm}$ ($\SI{3,6}{cm}$).
        \item Grossissement objectif $\gamma_O = -\frac{OA'}{OA} = -\frac{36}{4,5} = -8$. $A'B'$ est réelle (de l'autre côté de la lentille), inversée ($\gamma_O < 0$), agrandie ($|\gamma_O|=8$). Hauteur $A'B' = 8 \times \SI{0,10}{mm} = \SI{0,80}{mm}$.
        \item \textbf{Oculaire $L_{O'}$ :} Observation au repos $\Rightarrow$ objet de l'oculaire ($A'B'$) devrait être au foyer objet $F_{O'}$. Vérification : $O'A' = OO' - OA' = 160 - 36 = \SI{124}{mm}$. $F_{O'}O' = f_{O'} = \SI{25}{mm}$. $A'$ n'est pas en $F_{O'}$ ! \textbf{Attention :} L'énoncé impose $\overline{OO'}=160$ mm et $OA=4,5$ mm. Cela ne correspond pas à une observation standard au repos avec ces valeurs exactes. \emph{Au Probatoire, on construit le schéma avec les données données, on ne force pas la mise au point standard.} On trace les rayons pour l'oculaire : $B' \to$ parallèle à l'axe $\to$ réfracté par $F'_{O'}$ (foyer image oculaire). $B' \to O'$ $\to$ non dévié. Ces deux rayons émergents sont divergents. Leur prolongement en arrière donne l'image finale virtuelle $A''B''$.
    \end{itemize}
    \item \textbf{Image intermédiaire $A'B'$ :} \textbf{Réelle} (formée par convergence de rayons réels), \textbf{inversée} par rapport à $AB$ ($\gamma_O = -8$), \textbf{agrandie} ($|\gamma_O| = 8$).
    \item \textbf{Image finale $A''B''$ :} \textbf{Virtuelle} (formée par prolongement de rayons émergents divergents), \textbf{inversée} par rapport à $AB$ (l'oculaire ne renverse pas l'image intermédiaire déjà inversée), \textbf{agrandie}.
    \item \textbf{Grossissement commercial $G_c$ :} Par définition, pour un microscope, $G_c = \frac{\Delta}{f_O} \times \frac{D}{f_{O'}}$ avec $D = \SI{250}{mm}$ (distance de vision distincte) et $\Delta = \overline{F'_O F_{O'}}$ (intervalle optique).
    \[ \Delta = \overline{OO'} - f_O - f_{O'} = 160 - 4 - 25 = \SI{131}{mm} \]
    \[ G_c = \frac{131}{4} \times \frac{250}{25} = 32,75 \times 10 = \mathbf{327,5} \approx \mathbf{328} \]
    (Sans unité). L'image finale est 328 fois plus grande que l'objet vu à l'œil nu à \SI{250}{mm}.
\end{enumerate}
\cle{Grossissement commercial} \cle{Intervalle optique $\Delta$} \cle{Observation au repos}
\end{exoresolu}

\begin{methode}[Calculer le grossissement commercial $G_c$]
\textbf{Formule fondamentale :} $G_c = \gamma_O \times \gamma_{O'}$ (valeur absolue).
\begin{enumerate}
    \item \textbf{Identifier le mode d'observation :}
    \begin{itemize}
        \item \textbf{Au repos (image à l'infini) :} $\gamma_{O'} = \frac{D}{f_{O'}}$ avec $D = \SI{250}{mm}$. $\gamma_O = -\frac{\Delta}{f_O}$ (si objet au foyer objet de l'objectif, approximation standard) ou calcul exact via conjugaison.
        \item \textbf{Au point proche (image à $D$) :} $\gamma_{O'} = 1 + \frac{D}{f_{O'}}$. $\gamma_O$ ajusté pour que l'image intermédiaire soit à $p'_{O'} = \frac{D f_{O'}}{D+f_{O'}}$ de l'oculaire.
    \end{itemize}
    \item \textbf{Calculer $\Delta$ (intervalle optique) :} $\Delta = \overline{OO'} - f_O - f_{O'}$. Si $\overline{OO'}$ non donné, on utilise la définition standard : $\Delta = \SI{160}{mm}$ (norme mécanique tube) sauf indication contraire.
    \item \textbf{Appliquer la formule standard (au repos) :} $G_c = \frac{\Delta}{f_O} \times \frac{D}{f_{O'}}$.
    \item \textbf{Vérifier l'homogénéité des unités :} $f_O, f_{O'}, \Delta, D$ en \textbf{même unité} (mm ou m). $G_c$ est sans unité.
    \item \textbf{Arrondir :} Nombre entier généralement (ex: $\times 400$).
\end{enumerate}
\cle{Grossissement commercial $G_c$} \cle{Distance de vision distincte $D=250$ mm} \cle{Intervalle optique $\Delta$}
\end{methode}

\begin{exoresolu}[Calcul du grossissement commercial et choix d'instrument]
\textbf{Énoncé :} Un biologiste au Centre Pasteur de Yaoundé doit observer des bactéries de taille $\SI{2}{\mu m}$. Il dispose de deux microscopes :
\begin{itemize}
    \item \textbf{Micro A :} Objectif $f_O = \SI{4,0}{mm}$, Oculaire $f_{O'} = \SI{25}{mm}$, $\Delta = \SI{160}{mm}$.
    \item \textbf{Micro B :} Objectif $f_O = \SI{16}{mm}$, Oculaire $f_{O'} = \SI{25}{mm}$, $\Delta = \SI{160}{mm}$.
\end{itemize}
L'observation se fait au repos ($D = \SI{250}{mm}$). L'œil nu a une puissance séparatrice d'environ $\SI{1}{min}$ d'angle (soit $\approx \SI{3e-4}{rad}$), ce qui correspond à une taille minimale résolue d'environ $\SI{0,1}{mm}$ à \SI{250}{mm}.
\begin{enumerate}
    \item Calculer $G_{c,A}$ et $G_{c,B}$.
    \item Quelle taille apparente auront les bactéries avec chaque microscope ?
    \item Lequel choisir pour voir les bactéries distinctement ? Justifier.
\end{enumerate}

\tcblower
\textbf{Solution détaillée :}
\begin{enumerate}
    \item \textbf{Calcul des grossissements commerciaux (au repos) :}
    Formule : $G_c = \frac{\Delta}{f_O} \times \frac{D}{f_{O'}}$.
    \textbf{Micro A :} $f_O = \SI{4,0}{mm}, f_{O'} = \SI{25}{mm}, \Delta = \SI{160}{mm}, D = \SI{250}{mm}$.
    \[ G_{c,A} = \frac{160}{4,0} \times \frac{250}{25} = 40 \times 10 = \mathbf{400} \]
    \textbf{Micro B :} $f_O = \SI{16}{mm}$.
    \[ G_{c,B} = \frac{160}{16} \times \frac{250}{25} = 10 \times 10 = \mathbf{100} \]
    \item \textbf{Taille apparente des bactéries (taille de l'image finale vue à l'infini) :}
    La taille apparente correspond à la taille de l'image virtuelle finale si elle était placée à \SI{250}{mm}. Elle vaut $G_c \times \text{taille objet}$.
    Taille objet $h = \SI{2}{\mu m} = \SI{2e-3}{mm}$.
    \textbf{Micro A :} $h'_{A} = 400 \times \SI{2e-3}{mm} = \SI{0,80}{mm}$.
    \textbf{Micro B :} $h'_{B} = 100 \times \SI{2e-3}{mm} = \SI{0,20}{mm}$.
    \item \textbf{Choix :} L'œil nu résout $\SI{0,10}{mm}$ à \SI{250}{mm}.
    \begin{itemize}
        \item Micro A donne $\SI{0,80}{mm} > \SI{0,10}{mm}$ : \textbf{Les bactéries sont visibles distinctement}.
        \item Micro B donne $\SI{0,20}{mm} > \SI{0,10}{mm}$ : \textbf{Les bactéries sont aussi visibles} (mais limite, détail faible).
    \end{itemize}
    \textbf{Conclusion :} On choisira le \textbf{Micro A} ($G_c = 400$) car il offre une taille apparente plus grande ($\SI{0,80}{mm}$), permettant une meilleure observation des détails (meilleure résolution angulaire effective). De plus, son objectif à focale plus courte ($\SI{4}{mm}$) a généralement une plus grande ouverture numérique, donc une meilleure résolution physique (limite de diffraction).
\end{enumerate}
\cle{Puissance séparatrice de l'œil} \cle{Taille apparente} \cle{Choix d'instrument}
\end{exoresolu}

\begin{methode}[Calculer la puissance commerciale $P_c$ (en dioptries)]
\textbf{Définition :} La puissance commerciale $P_c$ (en dioptries, $\delta$ ou D) est la puissance équivalente du système objectif+oculaire pour la vision au repos. Elle caractérise la "force" de l'instrument.
\begin{enumerate}
    \item \textbf{Formule de base :} $P_c = \frac{G_c}{D}$ avec $D = \SI{0,250}{m}$ (distance de vision distincte en \textbf{mètres}).
    \item \textbf{Formule pratique (au repos) :} En remplaçant $G_c = \frac{\Delta}{f_O} \frac{D}{f_{O'}}$, on obtient :
    \[ P_c = \frac{\Delta}{f_O f_{O'}} \]
    \textbf{Attention :} $f_O, f_{O'}, \Delta$ doivent être en \textbf{mètres} pour obtenir $P_c$ en $\text{m}^{-1}$ (dioptries).
    \item \textbf{Unité :} Dioptrie ($\delta$ ou D) = $\text{m}^{-1}$.
    \item \textbf{Vérification dimensionnelle :} $[\Delta] = L$, $[f_O] = L$, $[f_{O'}] = L \Rightarrow [\frac{\Delta}{f_O f_{O'}}] = L^{-1} = \text{Dioptrie}$. Correct.
\end{enumerate}
\cle{Puissance commerciale $P_c$} \cle{Dioptrie ($\delta$)} \cle{Formule $P_c = \frac{\Delta}{f_O f_{O'}}$}
\end{methode}

\begin{exoresolu}[Calcul de la puissance commerciale et vérification dimensionnelle]
\textbf{Énoncé :} Un microscope a un intervalle optique $\Delta = \SI{160}{mm}$. L'objectif a une puissance $P_O = \SI{250}{\delta}$ et l'oculaire une puissance $P_{O'} = \SI{40}{\delta}$.
\begin{enumerate}
    \item Déduire les focales $f_O$ et $f_{O'}$ en mètres et en mm.
    \item Calculer la puissance commerciale $P_c$ de ce microscope (observation au repos) en utilisant la formule $P_c = \frac{\Delta}{f_O f_{O'}}$.
    \item Vérifier le résultat avec la formule de la puissance équivalente de deux lentilles minces séparées par $d = \Delta + f_O + f_{O'}$ : $P_{eq} = P_O + P_{O'} - d P_O P_{O'}$. Relier $P_c$ et $P_{eq}$.
    \item Calculer le grossissement commercial $G_c$.
\end{enumerate}

\tcblower
\textbf{Solution détaillée :}
\begin{enumerate}
    \item \textbf{Focales :} $f = \frac{1}{P}$ (avec $P$ en $\text{m}^{-1}$, $f$ en m).
    $P_O = \SI{250}{\delta} \Rightarrow f_O = \frac{1}{250} = \SI{0,0040}{m} = \mathbf{\SI{4,0}{mm}}$.
    $P_{O'} = \SI{40}{\delta} \Rightarrow f_{O'} = \frac{1}{40} = \SI{0,025}{m} = \mathbf{\SI{25}{mm}}$.
    \item \textbf{Calcul de $P_c$ (Formule standard) :} $P_c = \frac{\Delta}{f_O f_{O'}}$.
    Conversion en mètres : $\Delta = \SI{160}{mm} = \SI{0,160}{m}$, $f_O = \SI{0,0040}{m}$, $f_{O'} = \SI{0,025}{m}$.
    \[ P_c = \frac{0,160}{0,0040 \times 0,025} = \frac{0,160}{0,00010} = \mathbf{1600\,\delta} \]
    \item \textbf{Vérification (Puissance équivalente du système) :}
    La distance entre les centres optiques est $d = \Delta + f_O + f_{O'} = 0,160 + 0,004 + 0,025 = \SI{0,189}{m}$.
    La puissance équivalente de l'association (objectif + oculaire) est donnée par la formule de Gullstrand :
    \[ P_{eq} = P_O + P_{O'} - d P_O P_{O'} = 250 + 40 - 0,189 \times 250 \times 40 \]
    \[ P_{eq} = 290 - 1890 = \mathbf{-1600\,\delta} \]
    Le signe négatif indique que le système global se comporte comme une lentille divergente pour un objet à l'infini (ce qui n'est pas le cas d'usage du microscope).
    La \textbf{puissance commerciale} est définie comme l'inverse de la distance de vision distincte multipliée par le grossissement : $P_c = G_c/D$. Pour un microscope utilisé au repos, on a la relation $P_c = -P_{eq} = \mathbf{1600\,\delta}$.
    \item \textbf{Grossissement commercial :} $G_c = P_c \times D = 1600 \times 0,250 = \mathbf{400}$.
    (Vérification : $G_c = \frac{160}{4} \times \frac{250}{25} = 40 \times 10 = 400$. Cohérent.)
\end{enumerate}
\cle{Analyse dimensionnelle} \cle{Puissance équivalente vs Puissance commerciale} \cle{Unité SI : mètre}
\end{exoresolu}

\begin{methode}[Mise au point : Théorie (Relations de conjugaison) et Pratique]
\textbf{But :} Obtenir une image nette de l'objet observé.
\begin{enumerate}
    \item \textbf{Théorie (Deux degrés de liberté) :}
    \begin{itemize}
        \item \textbf{Mise au point "macro" (grosse mise au point) :} On déplace l'ensemble \textbf{objectif + oculaire} (le tube) par rapport à la platine (objet) pour placer l'objet au foyer objet de l'objectif (ou légèrement au-delà). Cela forme l'image intermédiaire au plan focal de l'oculaire (vision au repos) ou à la distance voulue (vision au point proche).
        \item \textbf{Mise au point "micro" (fine mise au point) :} On déplace légèrement l'\textbf{oculaire} par rapport à l'objectif (variation de $\Delta$) pour placer exactement l'image intermédiaire au foyer de l'oculaire (vision au repos) ou compenser l'accommodation de l'observateur.
    \end{itemize}
    \item \textbf{Relations de conjugaison (Calcul exact) :}
    \begin{enumerate}
        \item \textbf{Objectif :} $\frac{1}{f_O} = \frac{1}{p_O} + \frac{1}{p'_O}$. $p_O = \overline{OA}$ (distance objet-objectif, $>0$). $p'_O = \overline{OA'}$ (distance objectif-image intermédiaire).
        \item \textbf{Oculaire :} $\frac{1}{f_{O'}} = \frac{1}{p_{O'}} + \frac{1}{p'_{O'}}$.
        \item \textbf{Condition de montage :} $p_{O'} = \overline{O'A'} = \overline{OO'} - p'_O$.
        \item \textbf{Condition d'observation :}
        \begin{itemize}
            \item Au repos : $p'_{O'} = -\infty \Rightarrow p_{O'} = f_{O'}$.
            \item Au point proche : $p'_{O'} = -D \Rightarrow p_{O'} = \frac{D f_{O'}}{D+f_{O'}}$.
        \end{itemize}
    \end{enumerate}
    \item \textbf{Pratique (Manipulation) :}
    \begin{enumerate}
        \item Éclairer l'échantillon (diaphragme iris / condenseur).
        \item Objectif de plus faible grossissement (ex $\times 4$ ou $\times 10$).
        \item \textbf{Grosse mise au point :} Monter le tube (ou descendre la platine) au maximum (rapprocher objectif/objet). Regarder dans l'oculaire. Descendre le tube (ou monter la platine) \textbf{lentement} jusqu'à apparition de l'image nette. \textbf{Jamais l'inverse} (risque de casser la lame/objetif).
        \item \textbf{Fine mise au point :} Ajuster la molette fine pour netteté parfaite.
        \item Centrer la zone d'intérêt. Passer à l'objectif supérieur (ex $\times 40$). \textbf{Refaire seulement la fine mise au point} (l'objectif est parfocal).
    \end{enumerate}
\end{enumerate}
\cle{Mise au point macro/micro} \cle{Parfocalité} \cle{Sécurité : ne pas casser la lame}
\end{methode}

\begin{exoresolu}[Mise au point : Calcul des positions pour observation au point proche]
\textbf{Énoncé :} Un microscope a $f_O = \SI{5,0}{mm}$, $f_{O'} = \SI{20}{mm}$. L'observateur veut observer au point proche ($D = \SI{250}{mm}$). L'objet est placé à $p_O = \SI{5,5}{mm}$ de l'objectif.
\begin{enumerate}
    \item Calculer la position de l'image intermédiaire $A'B'$ par rapport à l'objectif ($p'_O$).
    \item Déduire la distance $\overline{O'A'}$ (objet pour l'oculaire) nécessaire pour former l'image finale à $D$.
    \item En déduire la distance entre lentilles $\overline{OO'}$ (encombrement optique) pour cette mise au point.
    \item Calculer le grossissement total $\gamma$ dans cette configuration.
\end{enumerate}

\tcblower
\textbf{Solution détaillée :}
\begin{enumerate}
    \item \textbf{Objectif :} Relation de Descartes : $\frac{1}{f_O} = \frac{1}{p_O} + \frac{1}{p'_O}$.
    $f_O = \SI{5,0}{mm}$, $p_O = \SI{5,5}{mm}$.
    $\frac{1}{p'_O} = \frac{1}{5,0} - \frac{1}{5,5} = \frac{5,5 - 5,0}{27,5} = \frac{0,5}{27,5} = \frac{1}{55}$.
    $p'_O = \mathbf{\SI{55}{mm}}$. (Image réelle, de l'autre côté de l'objectif).
    Grossissement objectif $\gamma_O = -\frac{p'_O}{p_O} = -\frac{55}{5,5} = \mathbf{-10}$.
    \item \textbf{Oculaire (Vision au point proche) :} Image finale virtuelle en $A''$ : $p'_{O'} = -D = -\SI{250}{mm}$.
    Relation de Descartes : $\frac{1}{f_{O'}} = \frac{1}{p_{O'}} + \frac{1}{p'_{O'}}$.
    $\frac{1}{20} = \frac{1}{p_{O'}} + \frac{1}{-250} \Rightarrow \frac{1}{p_{O'}} = \frac{1}{20} + \frac{1}{250} = \frac{12,5 + 1}{250} = \frac{13,5}{250} = \frac{27}{500}$.
    $p_{O'} = \frac{500}{27} \approx \mathbf{\SI{18,52}{mm}}$.
    (C'est la distance $\overline{O'A'}$, objet réel pour l'oculaire, donc $A'$ est entre $O'$ et $F_{O'}$).
    Grossissement oculaire $\gamma_{O'} = -\frac{p'_{O'}}{p_{O'}} = -\frac{-250}{18,52} \approx \mathbf{+13,5}$.
    \item \textbf{Distance entre lentilles $\overline{OO'}$ :}
    $\overline{OO'} = p'_O + p_{O'} = 55 + 18,52 = \mathbf{\SI{73,52}{mm}}$.
    Intervalle optique $\Delta = \overline{OO'} - f_O - f_{O'} = 73,52 - 5 - 20 = \SI{48,52}{mm}$.
    \item \textbf{Grossissement total $\gamma$ :} $\gamma = \gamma_O \times \gamma_{O'} = (-10) \times (+13,5) = \mathbf{-135}$.
    L'image finale est inversée (signe -), agrandie 135 fois.
    \textbf{Note :} Grossissement commercial standard (au repos, $\Delta=160$) : $G_c = \frac{160}{5} \times \frac{250}{20} = 32 \times 12,5 = 400$. Ici $\gamma=135$ car $\Delta$ est petit (rapprochement lentilles pour mise au point au point proche avec cet objet loin du foyer).
\end{enumerate}
\cle{Relation de Descartes} \cle{Vision au point proche} \cle{Distance inter-lentilles}
\end{exoresolu}

\begin{methode}[Calculer la longueur du tube $\Delta$ et l'encombrement $L$]
\textbf{Définitions :}
\begin{itemize}
    \item \textbf{Intervalle optique (Longueur de tube optique) $\Delta$ :} Distance entre le foyer image de l'objectif $F'_O$ et le foyer objet de l'oculaire $F_{O'}$. $\Delta = \overline{F'_O F_{O'}}$.
    \item \textbf{Encombrement optique $L$ (ou longueur mécanique du tube) :} Distance entre les centres optiques $O$ et $O'$. $L = \overline{OO'} = \Delta + f_O + f_{O'}$.
    \item \textbf{Norme standard (DIN/ISO) :} $\Delta = \SI{160}{mm}$ (microscopes "finis" anciens) ou $\Delta = \infty$ (microscopes "à l'infini" modernes, tube lentille). Au Probatoire, on suppose $\Delta = \SI{160}{mm}$ sauf précision.
\end{itemize}
\textbf{Méthode de calcul :}
\begin{enumerate}
    \item Identifier les données : $f_O, f_{O'}, G_c$ (ou $P_c$), $\Delta$ (donné ou standard).
    \item Si $G_c$ connu et $\Delta$ inconnu (au repos) : $G_c = \frac{\Delta}{f_O} \frac{D}{f_{O'}} \Rightarrow \Delta = G_c \frac{f_O f_{O'}}{D}$.
    \item Calculer $L = \Delta + f_O + f_{O'}$.
    \item Vérifier la compatibilité mécanique (ex: $L$ doit rentrer dans le tube du microscope).
\end{enumerate}
\cle{Intervalle optique $\Delta$} \cle{Encombrement $L$} \cle{Norme $\Delta = 160$ mm}
\end{methode}

\begin{exoresolu}[Longueur du tube et encombrement (Probatoire 2019)]
\textbf{Énoncé :} On souhaite concevoir un microscope de grossissement commercial $G_c = 600$ (observation au repos). L'objectif disponible a une focale $f_O = \SI{4,0}{mm}$. L'oculaire a une focale $f_{O'} = \SI{25}{mm}$. La distance de vision distincte est $D = \SI{250}{mm}$. L'intervalle optique standard est $\Delta = \SI{160}{mm}$.
\begin{enumerate}
    \item Calculer l'intervalle optique $\Delta$ nécessaire pour obtenir ce grossissement avec ces lentilles.
    \item Calculer l'encombrement optique $L = \overline{OO'}$ correspondant.
    \item Comme $\Delta$ nécessaire diffère de la norme ($\SI{160}{mm}$), proposer une solution technique réaliste (changement d'objectif ou d'oculaire) pour atteindre $G_c \approx 600$ avec $\Delta = \SI{160}{mm}$.
\end{enumerate}

\tcblower
\textbf{Solution détaillée :}
\begin{enumerate}
    \item \textbf{Calcul de $\Delta$ nécessaire :} Formule au repos : $G_c = \frac{\Delta}{f_O} \times \frac{D}{f_{O'}}$.
    $\Delta = G_c \times \frac{f_O f_{O'}}{D}$.
    Unités en mm : $G_c = 600$, $f_O = 4,0$, $f_{O'} = 25$, $D = 250$.
    $\Delta = 600 \times \frac{4,0 \times 25}{250} = 600 \times \frac{100}{250} = 600 \times 0,4 = \mathbf{\SI{240}{mm}}$.
    \item \textbf{Encombrement $L$ :} $L = \Delta + f_O + f_{O'} = 240 + 4 + 25 = \mathbf{\SI{269}{mm}}$.
    \item \textbf{Solution technique (standard $\Delta = \SI{160}{mm}$) :}
    On impose $\Delta = \SI{160}{mm}$. Il faut changer les focales pour obtenir $G_c = 600$.
    $G_c = 600 = \frac{160}{f_O} \times \frac{250}{f_{O'}} = \frac{40000}{f_O f_{O'}}$ (avec $f$ en mm).
    $f_O f_{O'} = \frac{40000}{600} \approx 66,7 \text{ mm}^2$.
    \begin{itemize}
        \item Option A : Garder oculaire $\times 10$ ($f_{O'}=25$ mm) $\Rightarrow f_O = \frac{66,7}{25} \approx \SI{2,67}{mm}$ (Objectif $\times 60$ environ, distance de travail très courte).
        \item Option B : Garder objectif $\times 40$ ($f_O=4$ mm) $\Rightarrow f_{O'} = \frac{66,7}{4} \approx \SI{16,7}{mm}$ (Oculaire $\times 15$, standard).
        \item Option C : Objectif $\times 60$ ($f_O \approx 2,67$ mm) + Oculaire $\times 10$ ($f_{O'}=25$ mm) $\Rightarrow G_c = \frac{160}{2,67} \times 10 \approx 600$.
    \end{itemize}
    \textbf{Choix recommandé :} Objectif $\times 40$ ($f_O=4$ mm) + Oculaire $\times 15$ ($f_{O'} \approx 16,7$ mm) ou Objectif $\times 60$ + Oculaire $\times 10$. L'option Objectif $\times 60$ / Oculaire $\times 10$ est courante en microscopie professionnelle (Plan Achromat 60x/0.85).
\end{enumerate}
\cle{Longueur mécanique standard} \cle{Conception d'instrument} \cle{Compromis focales}
\end{exoresolu}

\begin{methode}[Résoudre un problème complet type Probatoire (Synthèse)]
\textbf{Démarche type "Bac/Probatoire" (4-6 points) :}
\begin{enumerate}
    \item \textbf{Lecture active :} Surligner les données numériques ($f_O, f_{O'}, \Delta, D, p_O$, mode observation). Identifier la question finale.
    \item \textbf{Schéma annoté :} Dessiner l'axe, les lentilles, les foyers, l'objet. Indiquer les distances connues ($p_O, \Delta, f_O, f_{O'}$). \textbf{Indiquer le sens de la lumière.}
    \item \textbf{Étape Objectif :} Calculer $p'_O$ (Descartes) et $\gamma_O$. Vérifier : $p'_O > 0$ (image réelle), $\gamma_O < 0$ (inversée).
    \item \textbf{Étape Oculaire :} Déterminer $p_{O'} = \overline{OO'} - p'_O$ (ou $\Delta + f_{O'} - (p'_O - f_O)$).
        \begin{itemize}
            \item Si "Au repos" : Vérifier/Imposer $p_{O'} = f_{O'}$. Sinon calculer $p'_{O'}$.
            \item Si "Au point proche" : Imposer $p'_{O'} = -D$. Calculer $p_{O'}$ requis.
        \end{itemize}
        Calculer $\gamma_{O'}$.
    \item \textbf{Bilan :} $\gamma = \gamma_O \gamma_{O'}$. $G_c = |\gamma|$ (au repos) ou $G_c = \frac{\Delta}{f_O}\frac{D}{f_{O'}}$ (valeur nominale).
    \item \textbf{Puissance :} $P_c = \frac{G_c}{D}$ (en m$^{-1}$).
    \item \textbf{Phrase de conclusion :} "L'image finale est virtuelle, inversée, agrandie de $|\gamma|$ fois. La puissance commerciale est $P_c = \dots\,\delta$."
\end{enumerate}
\cle{Résolution systématique} \cle{Schéma annoté obligatoire} \cle{Phrase de conclusion}
\end{methode}

\begin{exoresolu}[Problème de synthèse : Microscope corrigé à l'infini (Moderne)]
\textbf{Énoncé :} Un microscope moderne "corrigé à l'infini" (objectif plan-apochromat) est utilisé avec une lentille tube (focale $f_T = \SI{200}{mm}$) et un oculaire $f_{O'} = \SI{25}{mm}$. L'objectif a un grossissement inscrit $\times 40$ (référence tube $\SI{200}{mm}$). L'observation se fait au repos ($D = \SI{250}{mm}$).
\begin{enumerate}
    \item Quelle est la focale $f_O$ de cet objectif ? (Rappel : pour un objectif corrigé à l'infini, le grossissement inscrit $M_{obj} = \frac{f_T}{f_O}$).
    \item Calculer le grossissement commercial $G_c$ de l'ensemble (Objectif + Lentille tube + Oculaire).
    \item Calculer la puissance commerciale $P_c$.
    \item L'image intermédiaire formée par l'objectif est-elle réelle ou virtuelle ? Où se forme-t-elle par rapport à la lentille tube ?
\end{enumerate}

\tcblower
\textbf{Solution détaillée :}
\begin{enumerate}
    \item \textbf{Focale de l'objectif :} Pour un objectif corrigé à l'infini (standard RMS / DIN), le grossissement nominal $M_{obj}$ est défini pour une lentille tube de focale $f_T$ (ici $\SI{200}{mm}$).
    $M_{obj} = \frac{f_T}{f_O} \Rightarrow f_O = \frac{f_T}{M_{obj}} = \frac{200}{40} = \mathbf{\SI{5,0}{mm}}$.
    \item \textbf{Grossissement commercial $G_c$ :}
    Le système "Objectif + Lentille tube" forme une image intermédiaire réelle au foyer image de la lentille tube (distance $f_T$ de la lentille tube). Cette image sert d'objet à l'oculaire.
    Grossissement de l'objectif (réel) : $\gamma_{obj} = -M_{obj} = -40$ (image inversée).
    Grossissement de l'oculaire (au repos) : $\gamma_{O'} = \frac{D}{f_{O'}} = \frac{250}{25} = 10$.
    Grossissement total (commercial) : $G_c = |\gamma_{obj}| \times \gamma_{O'} = 40 \times 10 = \mathbf{400}$.
    (Note : On utilise la valeur absolue pour le grossissement commercial affiché, ex: 400x).
    \item \textbf{Puissance commerciale $P_c$ :} $P_c = \frac{G_c}{D} = \frac{400}{0,250} = \mathbf{1600\,\delta}$.
    \item \textbf{Nature de l'image intermédiaire :}
    L'objectif est convergent ($f_O > 0$). L'objet est placé légèrement au-delà de $F_O$ (distance de travail $\approx \SI{0,1}{mm}$ à $\SI{1}{mm}$). L'objectif forme une image \textbf{réelle}, inversée, \textbf{à l'infini} (faisceaux émergents parallèles).
    La lentille tube (convergente, $f_T = \SI{200}{mm}$) reçoit ces faisceaux parallèles. Elle forme l'image intermédiaire réelle \textbf{dans son plan focal image}, c'est-à-dire à \SI{200}{mm} derrière elle.
    Cette image intermédiaire réelle est ensuite grossie par l'oculaire.
\end{enumerate}
\cle{Microscope corrigé à l'infini} \cle{Lentille tube} \cle{Objectif Plan-Apochromat}
\end{exoresolu}

\begin{methode}[Choisir l'objectif et l'oculaire pour un grossissement visé (Application terrain)]
\textbf{Contexte :} Au laboratoire (ex: IRAD, hôpital), on dispose d'un jeu d'objectifs et d'oculaires. On doit configurer le microscope.
\begin{enumerate}
    \item \textbf{Lister les composants disponibles :} Objectifs (ex: $\times 4, \times 10, \times 40, \times 100$ huile) $\to$ focales $f_O$. Oculaires (ex: $\times 10, \times 15$) $\to$ focales $f_{O'}$.
    \item \textbf{Calculer les couples $G_c$ possibles :} $G_c = M_{obj} \times M_{oc}$ (valeurs inscrites sur les barrels). Ex: $\times 40 \times \times 10 = \times 400$.
    \item \textbf{Choisir selon la résolution nécessaire :} La limite de résolution (Abbe) $\delta = \frac{\lambda}{2 NA}$. $NA$ (Ouverture Numérique) est inscrit sur l'objectif (ex: 0.10, 0.25, 0.65, 1.25).
    \begin{itemize}
        \item Grossissement utile : $500 \times NA < G_c < 1000 \times NA$ (règle empirique pour voir le détail résolu sans "vide grossissement").
        \item Ex: Objectif $\times 40 / NA 0.65 \Rightarrow 325 < G_c < 650$. Oculaire $\times 10 (G_c=400)$ OK. Oculaire $\times 15 (G_c=600)$ OK.
    \end{itemize}
    \item \textbf{Vérifier la distance de travail (WD) :} Objectif $\times 100$ (huile) a WD $\approx \SI{0,1}{mm}$ $\to$ risque de casser la lame. Objectif $\times 4$ a WD $\approx \SI{10}{mm}$ $\to$ facile pour manipuler.
    \item \textbf{Procédure :} Commencer $\times 4$ ou $\times 10$ (champ large, WD grande) pour trouver la zone. Monter progressivement.
\end{enumerate}
\cle{Ouverture numérique $NA$} \cle{Grossissement utile} \cle{Vide grossissement} \cle{Distance de travail}
\end{methode}

\begin{exoresolu}[Choix d'objectif selon la résolution (Application biologie)]
\textbf{Énoncé :} On observe une coupe de tissu végétal (cellules de $\approx \SI{20}{\mu m}$). Le microscope a un éclairage $\lambda = \SI{550}{nm}$ (vert). Objectifs disponibles :
\begin{center}
\begin{tabular}{|c|c|c|}\hline
\textbf{Objectif} & \textbf{Grossissement $M_{obj}$} & \textbf{Ouverture Numérique $NA$} \\ \hline
Plan Achromat & $\times 4$ & 0,10 \\
Plan Achromat & $\times 10$ & 0,25 \\
Plan Fluorite & $\times 40$ & 0,75 \\
Plan Apochromat & $\times 100$ (Huile) & 1,30 \\ \hline
\end{tabular}
\end{center}
Oculaires : $\times 10$ ($f_{O'}=25$ mm) et $\times 15$ ($f_{O'}=16,7$ mm). $\Delta = \SI{160}{mm}$ (standard).
\begin{enumerate}
    \item Calculer la limite de résolution $\delta$ (critère d'Abbe $\delta = \lambda / (2 NA)$) pour chaque objectif.
    \item Pour chaque objectif, calculer $G_c$ avec les deux oculaires.
    \item Appliquer la règle du grossissement utile ($500 NA < G_c < 1000 NA$) pour identifier les couples (Objectif, Oculaire) "utiles".
    \item Recommander la meilleure configuration pour voir les détails internes des cellules ($\sim \SI{1}{\mu m}$).
\end{enumerate}

\tcblower
\textbf{Solution détaillée :}
\begin{enumerate}
    \item \textbf{Limite de résolution $\delta = \frac{\lambda}{2 NA}$ ($\lambda = \SI{550e-9}{m} = \SI{0,55}{\mu m}$) :}
    \begin{itemize}
        \item $\times 4 / NA 0,10 : \delta = \frac{0,55}{0,20} = \SI{2,75}{\mu m}$.
        \item $\times 10 / NA 0,25 : \delta = \frac{0,55}{0,50} = \SI{1,10}{\mu m}$.
        \item $\times 40 / NA 0,75 : \delta = \frac{0,55}{1,50} = \SI{0,367}{\mu m}$.
        \item $\times 100 / NA 1,30 : \delta = \frac{0,55}{2,60} = \SI{0,212}{\mu m}$.
    \end{itemize}
    \item \textbf{Grossissements commerciaux $G_c = M_{obj} \times M_{oc}$ :}
    \begin{center}
    \begin{tabular}{|c|c|c|}\hline
    Objectif & Oculaire $\times 10$ & Oculaire $\times 15$ \\ \hline
    $\times 4$ & 40 & 60 \\
    $\times 10$ & 100 & 150 \\
    $\times 40$ & 400 & 600 \\
    $\times 100$ & 1000 & 1500 \\ \hline
    \end{tabular}
    \end{center}
    \item \textbf{Règle du grossissement utile ($500 NA < G_c < 1000 NA$) :}
    \begin{itemize}
        \item $\times 4 (NA 0,10)$ : $50 < G_c < 100$. Couples utiles : $\mathbf{\times 4 / \times 15 (G_c=60)}$ et $\mathbf{\times 4 / \times 10 (G_c=40)}$ limite bas. $\times 4/\times 10$ un peu faible.
        \item $\times 10 (NA 0,25)$ : $125 < G_c < 250$. Couples utiles : \textbf{Aucun} ($100$ trop faible, $150$ limite bas). $\times 10/\times 15$ acceptable.
        \item $\times 40 (NA 0,75)$ : $375 < G_c < 750$. Couples utiles : $\mathbf{\times 40 / \times 10 (400)}$ et $\mathbf{\times 40 / \times 15 (600)}$. \textbf{Les deux excellents}.
        \item $\times 100 (NA 1,30)$ : $650 < G_c < 1300$. Couples utiles : $\mathbf{\times 100 / \times 10 (1000)}$ et $\mathbf{\times 100 / \times 15 (1500)}$ (vide grossissement).
    \end{itemize}
    \item \textbf{Recommandation :} Pour voir des détails de $\SI{1}{\mu m}$, il faut $\delta < 1 \mu m$. Objectifs $\times 40$ ($\delta=0,37$) et $\times 100$ ($\delta=0,21$) conviennent.
    \textbf{Meilleur choix :} \textbf{Objectif $\times 40$ (NA 0,75) + Oculaire $\times 10$ ($G_c=400$)}.
    \begin{itemize}
        \item Résolution suffisante ($0,37 \mu m$).
        \item Grossissement utile ($400 \in [375, 750]$).
        \item Distance de travail confortable ($\sim \SI{0,5}{mm}$ vs $\SI{0,1}{mm}$ pour $\times 100$ huile).
        \item Pas d'huile nécessaire (plus rapide, propre).
    \end{itemize}
    L'oculaire $\times 15$ ($G_c=600$) est aussi bon pour plus de confort visuel (image plus grande), mais champ plus petit.
\end{enumerate}
\cle{Critère d'Abbe} \cle{Grossissement utile} \cle{Vide grossissement} \cle{Distance de travail}
\end{exoresolu}

\begin{attention}[Les 5 erreurs qui coûtent des points au Probatoire]
\begin{enumerate}
    \item \textbf{Confondre Grossissement $\gamma$ et Grossissement Commercial $G_c$.}
    $\gamma = \gamma_O \gamma_{O'}$ (algébrique, signe compte : image inversée $\gamma < 0$). $G_c = |\gamma|$ (valeur absolue, sans unité, notation $\times 400$). \textbf{Piège :} Écrire $G_c = -400$ ou oublier le $\times$ devant la valeur.
    \item \textbf{Unités incohérentes dans les formules ($P_c$, $\Delta$).}
    $P_c = \frac{\Delta}{f_O f_{O'}}$ exige $\Delta, f_O, f_{O'}$ en \textbf{mètres} pour avoir $P_c$ en $\delta$ ($\text{m}^{-1}$). $G_c = \frac{\Delta}{f_O}\frac{D}{f_{O'}}$ exige \textbf{même unité} pour toutes les longueurs (mm ou m). \textbf{Réflexe :} Convertir TOUT en mètres pour $P_c$, TOUT en mm pour $G_c$ (plus simple).
    \item \textbf{Oublier le signe de la distance image pour l'oculaire (vision au point proche).}
    Image finale virtuelle $\Rightarrow p'_{O'} < 0$. Au point proche : $p'_{O'} = -D = -\SI{250}{mm}$. Au repos : $p'_{O'} = -\infty$. Oublier le moins conduit à $p_{O'}$ faux et $\gamma_{O'}$ faux.
    \item \textbf{Confondre Intervalle Optique $\Delta$ et Encombrement $L$ (ou distance inter-lentilles $\overline{OO'}$).}
    $\Delta = \overline{F'_O F_{O'}}$. $L = \overline{OO'} = \Delta + f_O + f_{O'}$. L'énoncé donne souvent $L$ (ex: tube de $\SI{160}{mm}$) et l'élève utilise $L$ à la place de $\Delta$ dans $G_c = \frac{\Delta}{f_O}\frac{D}{f_{O'}}$. \textbf{Vérifier :} "Longueur du tube" = $\Delta$ (standard) ou $L$ ? Lire attentivement.
    \item \textbf{Schéma de principe bâclé ou non annoté.}
    \textbf{Exigences Probatoire :} Axe optique, Lentilles (formes), Foyers $F, F'$ étiquetés, Objet $AB$ (flèche), Image intermédiaire $A'B'$, Image finale $A''B''$, \textbf{Au moins 2 rayons particuliers par lentille} (parallèle $\to$ foyer, centre optique), Sens des flèches (lumière), Échelle respectée (qualitative). \textbf{Zéro point} si pas de schéma ou schéma sans rayons.
\end{enumerate}
\cle{Erreurs classiques} \cle{Probatoire} \cle{Rigueur notationnelle}
\end{attention}

\newpage

\coursec{Exercices}

\courssub{Je vérifie mes connaissances}

\begin{exercice}[QCM : Rôle des lentilles]
Dans un microscope composé fonctionnant en vision normale (image finale à l'infini), quelle affirmation est exacte ?
\begin{enumerate}
    \item L'objectif forme une image réelle, renversée et grossissante de l'objet.
    \item L'oculaire forme une image réelle, renversée et grossissante de l'image intermédiaire.
    \item L'objectif forme une image virtuelle, droite et grossissante de l'objet.
    \item L'oculaire fonctionne comme une loupe : il forme une image virtuelle à l'infini de l'image intermédiaire placée à son foyer image.
\end{enumerate}
\end{exercice}

\begin{exercice}[Vrai / Faux justifié : Grossissement et Puissance]
On considère un microscope dont l'objectif a une puissance $P_{\text{obj}} = 250\,\delta$ et l'oculaire une puissance $P_{\text{oc}} = 50\,\delta$. L'intervalle inter-foculaire est $\Delta = 16\,\text{cm}$.
\begin{enumerate}
    \item Le grossissement commercial $G_c$ vaut $2000$.
    \item La puissance commerciale $P_c$ vaut $300\,\delta$.
    \item Si on double $\Delta$ en gardant les mêmes lentilles, $G_c$ est doublé.
    \item L'unité du grossissement commercial est le dioptrie ($\delta$).
\end{enumerate}
Justifie chaque réponse par un calcul ou une référence à la formule.
\end{exercice}

\begin{exercice}[Définitions et notations]
Donne la définition précise (avec formule le cas échéant) des grandeurs suivantes :
\begin{enumerate}
    \item Intervalle inter-foculaire $\Delta$ (ou $L$).
    \item Grossissement commercial $G_c$ (vision normale).
    \item Puissance commerciale $P_c$.
    \item Condition de mise au point en vision au confort (distance de vision distincte $D = 25\,\text{cm}$).
\end{enumerate}
\end{exercice}

\begin{exercice}[Calculs directs : Objectif et Oculaire]
Un microscope est constitué d'un objectif de focale $f'_{\text{obj}} = 4\,\text{mm}$ et d'un oculaire de focale $f'_{\text{oc}} = 25\,\text{mm}$. L'intervalle inter-foculaire est réglé à $\Delta = 160\,\text{mm}$.
\begin{enumerate}
    \item Calcule la puissance de l'objectif $P_{\text{obj}}$ et celle de l'oculaire $P_{\text{oc}}$ en dioptries ($\delta$).
    \item Calcule le grossissement commercial $G_c$ (vision normale).
    \item Calcule la puissance commerciale $P_c$ en $\delta$.
    \item Quel est le grossissement $G_{\text{confort}}$ si l'observateur met au point pour une image finale à $25\,\text{cm}$ ?
\end{enumerate}
\end{exercice}

\courssub{Je m'entraîne}

\begin{exercice}[Choix de l'objectif pour un grossissement donné]
On souhaite réaliser un microscope fonctionnant en vision normale avec un grossissement commercial $G_c = 600$. L'intervalle inter-foculaire disponible sur le microscope est $\Delta = 16\,\text{cm}$. L'oculaire disponible a une puissance $P_{\text{oc}} = 50\,\delta$.
\begin{enumerate}
    \item Détermine la puissance $P_{\text{obj}}$ que doit avoir l'objectif.
    \item En déduire sa focale $f'_{\text{obj}}$ en millimètres.
    \item Les objectifs standards de microscopie ont des grossissements inscrits : $4\times$, $10\times$, $40\times$, $100\times$ (pour $\Delta = 160\,\text{mm}$). Lequel choisir ? Justifie.
\end{enumerate}
\end{exercice}

\begin{exercice}[Passage vision normale $\rightarrow$ vision au confort]
Un élève utilise un microscope réglé en vision normale (image finale à l'infini). L'oculaire a une focale $f'_{\text{oc}} = 25\,\text{mm}$. L'objectif utilisé est un $40\times$ (pour $\Delta = 160\,\text{mm}$).
\begin{enumerate}
    \item Rappelle la position de l'image intermédiaire par rapport à l'oculaire en vision normale.
    \item Pour passer en vision au confort (image finale nette à $D = 25\,\text{cm}$ de l'oculaire), l'image intermédiaire doit-elle se trouver avant ou après le foyer objet de l'oculaire ?
    \item En agissant sur la vis micrométrique (déplacement de l'objectif), dans quel sens faut-il déplacer l'objectif (vers l'objet ou vers l'oculaire) ? Justifie par la conjugaison de l'objectif.
    \item Estime la valeur approximative du déplacement de l'objectif (en mm) en supposant que le grossissement de l'objectif ne change pas significativement.
\end{enumerate}
\end{exercice}

\begin{exercice}[Analyse des inscriptions sur un objectif]
Un objectif de microscope porte la gravure : \textbf{Plan 40/0,65 $\infty$/0,17}.
\begin{enumerate}
    \item Explique la signification de chaque inscription : \texttt{Plan}, \texttt{40}, \texttt{0,65}, \texttt{$\infty$}, \texttt{0,17}.
    \item La mention \texttt{$\infty$} indique que l'objectif est corrigé à l'infini. Quelle est la conséquence sur la définition de l'intervalle inter-foculaire $\Delta$ pour cet objectif ? Comment se fait la mise au point dans ce cas ?
    \item Sachant que la lentille tube (intégrée au corps du microscope) a une focale image $f'_{\text{tube}} = 200\,\text{mm}$, calcule la focale image $f'_{\text{obj}}$ de cet objectif.
\end{enumerate}
\end{exercice}

\begin{exercice}[Pouvoir séparateur et immersion]
On utilise un objectif $100\times$ à immersion d'huile (indice $n = 1,515$) d'ouverture numérique $\text{NA} = 1,3$. La longueur d'onde utilisée est $\lambda = 500\,\text{nm}$.
\begin{enumerate}
    \item Calcule le pouvoir séparateur $\delta_{\text{min}}$ (distance minimale résolue) selon le critère de Rayleigh ($\delta_{\text{min}} = 0,61 \frac{\lambda}{\text{NA}}$).
    \item Une bactérie de longueur $0,5\,\mu\text{m}$ est-elle résolue (ses détails internes visibles) ?
    \item Deux points distants de $200\,\text{nm}$ sont-ils distingués ?
    \item Si on utilisait le même objectif \emph{sans} huile (air, $n=1$), l'ouverture numérique maximale serait $\text{NA}_{\text{air}} \approx 0,95$. Calcule le nouveau pouvoir séparateur et commente l'intérêt de l'immersion.
\end{enumerate}
\end{exercice}

\courssub{Je me prépare au Probatoire}

\begin{exercice}[Probatoire Type 2022 : Microscope en vision normale]
\textbf{Énoncé :} Un microscope fonctionne en vision normale. Son objectif a une puissance de $250\,\delta$, son oculaire une focale de $2,5\,\text{cm}$. L'intervalle inter-foculaire $\Delta = 16\,\text{cm}$.
\begin{enumerate}
    \item Calcule la focale de l'objectif $f'_{\text{obj}}$ en mm.
    \item Calcule la puissance de l'oculaire $P_{\text{oc}}$ en $\delta$.
    \item Détermine le grossissement commercial $G_c$ de cet instrument.
    \item Détermine sa puissance commerciale $P_c$ en $\delta$.
    \item Trace le schéma de principe complet de ce microscope. Indique : les deux lentilles (objectif et oculaire), leurs foyers objets et images ($F_{\text{obj}}, F'_{\text{obj}}, F_{\text{oc}}, F'_{\text{oc}}$), l'intervalle inter-foculaire $\Delta$, l'objet $AB$, l'image intermédiaire $A'B'$ et l'image finale à l'infini. Trace la construction géométrique de $A'B'$ (deux rayons sortants de $A$) et indique les rayons émergents parallèles sortant de l'oculaire.
\end{enumerate}
\end{exercice}

\begin{exercice}[Probatoire Type 2020 : Conception d'un microscope]
\textbf{Énoncé :} On souhaite réaliser un microscope de grossissement commercial $G_c = 600$ (vision normale) avec un intervalle inter-foculaire $\Delta = 16\,\text{cm}$. L'oculaire disponible a une puissance de $50\,\delta$.
\begin{enumerate}
    \item Détermine la puissance $P_{\text{obj}}$ que doit avoir l'objectif.
    \item En déduire sa focale image $f'_{\text{obj}}$ en mm.
    \item Les objectifs standards du laboratoire ont les grossissements inscrits suivants (pour $\Delta = 160\,\text{mm}$) : $4\times$, $10\times$, $40\times$, $100\times$. L'objectif calculé correspond-il à l'un de ces standards ? Justifie par le calcul du grossissement objectif $G_{\text{obj}} = \frac{\Delta}{f'_{\text{obj}}}$.
    \item Si l'observateur myope (non corrigé) veut utiliser ce microscope en vision normale, doit-il porter ses lunettes ? Justifie.
\end{enumerate}
\end{exercice}

\begin{exercice}[Probatoire Type : Mise au point, Schéma et Résolution]
\textbf{Énoncé :} On observe une lame de sang au microscope. L'instrument est équipé d'un objectif $40\times$ (focale $4\,\text{mm}$) et d'un oculaire $10\times$ (focale $25\,\text{mm}$). L'intervalle inter-foculaire effectif est $\Delta = 160\,\text{mm}$. L'ouverture numérique de l'objectif est $\text{NA} = 0,65$. La longueur d'onde moyenne est $\lambda = 550\,\text{nm}$.
\begin{enumerate}
    \item Calcule le grossissement commercial $G_c$ et la puissance commerciale $P_c$.
    \item L'observateur met au point en vision normale. Il souhaite passer en vision au confort ($D = 25\,\text{cm}$) sans changer d'objectif. 
    \begin{enumerate}
        \item Où se trouve l'image intermédiaire par rapport à l'oculaire en vision normale ?
        \item Où doit-elle se trouver pour la vision au confort ? (Donne la distance algébrique $\overline{F_{\text{oc}}A'B'}$).
        \item Dans quel sens (rapprocher ou éloigner l'objectif de la lame) doit-il tourner la vis macrométrique ? Justifie par la relation de conjugaison de l'objectif.
    \end{enumerate}
    \item Calcule le pouvoir séparateur $\delta_{\text{min}}$ de cet objectif (critère de Rayleigh).
    \item Deux globules rouges sont espacés de $300\,\text{nm}$. Sont-ils résolus ? Et s'ils sont espacés de $800\,\text{nm}$ ?
    \item Sur la figure ci-dessous (schéma incomplet), l'image intermédiaire $A'B'$ est tracée \emph{avant} le foyer objet de l'oculaire $F_{\text{oc}}$.
    \begin{enumerate}
        \item L'image finale est-elle à l'infini ou à distance finie ? Précise si elle est réelle ou virtuelle.
        \item Indique par une flèche le sens de rotation de la vis micrométrique pour ramener l'image intermédiaire au foyer de l'oculaire (vision normale).
        \item Trace les rayons émergents sortant de l'oculaire correspondant au point $B'$.
    \end{enumerate}
\end{enumerate}
\end{exercice}

\newpage

\coursec{Corrigés des exercices}

\coursec{Résolution de problèmes types (Probatoire)}

\begin{corrige}[Exercice 1 — QCM : Rôle des lentilles]
L'affirmation exacte est la \textbf{proposition 4}.

\emph{Justifications :}
\begin{enumerate}
    \item \textbf{Fausse.} L'objectif forme bien une image \textbf{réelle}, \textbf{renversée} et \textbf{grossissante} de l'objet (c'est un grossissement transverse $G_{\text{obj}} = -\frac{\Delta}{f'_{\text{obj}}} < -1$). L'énoncé dit « image virtuelle », ce qui est incorrect pour l'objectif.
    \item \textbf{Fausse.} L'oculaire fonctionne comme une loupe : il forme une image \textbf{virtuelle}, \textbf{droite} par rapport à l'image intermédiaire (donc renversée par rapport à l'objet initial) et \textbf{grossissante} angulairement. Il ne forme pas d'image réelle.
    \item \textbf{Fausse.} L'image formée par l'objectif est réelle (elle se forme sur le plan focal image de l'oculaire en vision normale), pas virtuelle.
    \item \textbf{Vraie.} En vision normale, l'image intermédiaire $A'B'$ est placée au foyer objet $F_{\text{oc}}$ de l'oculaire. Ce dernier fonctionne alors comme une loupe utilisée en vision normale : il transforme le faisceau convergent issu de l'objectif en faisceau émergent \textbf{parallèle} (image finale à l'infini), virtuelle, droite par rapport à $A'B'$ et grossissante angulairement.
\end{enumerate}
\end{corrige}

\begin{corrige}[Exercice 2 — Vrai / Faux justifié : Grossissement et Puissance]
Données : $P_{\text{obj}} = 250\,\delta$, $P_{\text{oc}} = 50\,\delta$, $\Delta = 16\,\text{cm} = 0,16\,\text{m}$.
Rappels : $f'_{\text{obj}} = 1/P_{\text{obj}} = 4\,\text{mm}$, $f'_{\text{oc}} = 1/P_{\text{oc}} = 20\,\text{mm}$.
Formules (vision normale, $D = 0,25\,\text{m}$) :
\[
G_c = \frac{\Delta}{f'_{\text{obj}}} \times \frac{D}{f'_{\text{oc}}} = \Delta \cdot P_{\text{obj}} \cdot D \cdot P_{\text{oc}}, \qquad
P_c = \frac{G_c}{D} = \Delta \cdot P_{\text{obj}} \cdot P_{\text{oc}}.
\]

\begin{enumerate}
    \item \textbf{Faux.}
    Calcul : $G_c = 0,16 \times 250 \times 0,25 \times 50 = 0,16 \times 3125 = 500$.
    Le grossissement commercial vaut $500$ (sans unité), pas $2000$.
    \item \textbf{Faux.}
    Calcul : $P_c = \Delta \cdot P_{\text{obj}} \cdot P_{\text{oc}} = 0,16 \times 250 \times 50 = 2000\,\delta$.
    La puissance commerciale vaut $2000\,\delta$, pas $300\,\delta$.
    \item \textbf{Vrai.}
    $G_c = \Delta \cdot P_{\text{obj}} \cdot D \cdot P_{\text{oc}}$. Si $\Delta$ double (autres paramètres constants), $G_c$ double. Relation de proportionnalité directe.
    \item \textbf{Faux.}
    Le grossissement commercial $G_c$ est un \textbf{nombre sans unité} (rapport de deux angles ou de deux longueurs). L'unité $\delta$ (dioptrie) est celle de la \textbf{puissance commerciale} $P_c$ (ou des puissances des lentilles).
\end{enumerate}
\end{corrige}

\begin{corrige}[Exercice 3 — Définitions et notations]
\begin{enumerate}
    \item \textbf{Intervalle inter-foculaire $\Delta$ (ou $L$) :} Distance algébrique $\overline{F'_{\text{obj}}F_{\text{oc}}}$ entre le foyer image de l'objectif ($F'_{\text{obj}}$) et le foyer objet de l'oculaire ($F_{\text{oc}}$). $\Delta > 0$ pour un microscope standard. C'est une longueur caractéristique de l'instrument (souvent $160\,\text{mm}$ pour les microscopes à tube fini).
    \item \textbf{Grossissement commercial $G_c$ (vision normale) :} Produit du grossissement transverse de l'objectif $G_{\text{obj}}$ par le grossissement angulaire commercial de l'oculaire $G_{\text{oc}}$.
    \[
    G_c = G_{\text{obj}} \times G_{\text{oc}} = \left(-\frac{\Delta}{f'_{\text{obj}}}\right) \times \left(\frac{D}{f'_{\text{oc}}}\right) = -\frac{\Delta \cdot D}{f'_{\text{obj}} f'_{\text{oc}}}.
    \]
    On retient souvent la valeur absolue $|G_c| = \frac{\Delta}{f'_{\text{obj}}} \cdot \frac{0,25}{f'_{\text{oc}}}$ (avec $\Delta$ et $f'$ en mètres, $D=0,25\,\text{m}$). $G_c$ est sans unité.
    \item \textbf{Puissance commerciale $P_c$ :} Quotient du grossissement commercial par la distance de vision distincte $D = 0,25\,\text{m}$.
    \[
    P_c = \frac{G_c}{D} = \frac{\Delta}{f'_{\text{obj}} f'_{\text{oc}}} = \Delta \cdot P_{\text{obj}} \cdot P_{\text{oc}}.
    \]
    Son unité est la \textbf{dioptrie} ($\delta = \text{m}^{-1}$).
    \item \textbf{Condition de mise au point en vision au confort ($D = 25\,\text{cm}$) :} L'image finale $A''B''$ doit être nette pour l'œil placé au plus près de l'oculaire, à la distance de vision distincte $D = -25\,\text{cm}$ (image virtuelle, même côté que l'objet pour l'oculaire).
    Relation de conjugaison de l'oculaire (formule de Descartes, origine à l'oculaire, sens de la lumière $+x$) :
    \[
    \frac{1}{f'_{\text{oc}}} = \frac{1}{\overline{\text{Oc}A''B''}} - \frac{1}{\overline{\text{Oc}A'B'}} \quad \text{avec} \quad \overline{\text{Oc}A''B''} = -D = -0,25\,\text{m}.
    \]
    L'image intermédiaire $A'B'$ doit donc se trouver \textbf{entre le foyer objet $F_{\text{oc}}$ et l'oculaire} (distance $\overline{F_{\text{oc}}A'B'} = \frac{{f'_{\text{oc}}}^2}{D+f'_{\text{oc}}} > 0$).
\end{enumerate}
\end{corrige}

\begin{corrige}[Exercice 4 — Calculs directs : Objectif et Oculaire]
Données : $f'_{\text{obj}} = 4\,\text{mm} = 4\times10^{-3}\,\text{m}$, $f'_{\text{oc}} = 25\,\text{mm} = 2,5\times10^{-2}\,\text{m}$, $\Delta = 160\,\text{mm} = 0,16\,\text{m}$, $D = 0,25\,\text{m}$.

\begin{enumerate}
    \item Puissances :
    \[
    P_{\text{obj}} = \frac{1}{f'_{\text{obj}}} = \frac{1}{4\times10^{-3}} = 250\,\delta, \qquad
    P_{\text{oc}} = \frac{1}{f'_{\text{oc}}} = \frac{1}{2,5\times10^{-2}} = 40\,\delta.
    \]
    \item Grossissement commercial (vision normale) :
    \[
    G_c = \frac{\Delta}{f'_{\text{obj}}} \cdot \frac{D}{f'_{\text{oc}}} = \frac{0,16}{4\times10^{-3}} \times \frac{0,25}{2,5\times10^{-2}} = 40 \times 10 = \mathbf{400}.
    \]
    \item Puissance commerciale :
    \[
    P_c = \frac{G_c}{D} = \frac{400}{0,25} = \mathbf{1600\,\delta}.
    \]
    (Vérification : $P_c = \Delta \cdot P_{\text{obj}} \cdot P_{\text{oc}} = 0,16 \times 250 \times 40 = 1600\,\delta$.)
    \item Grossissement en vision au confort ($D=0,25\,\text{m}$) :
    Grossissement de l'oculaire en vision au confort : $G_{\text{oc,confort}} = 1 + \frac{D}{f'_{\text{oc}}} = 1 + \frac{0,25}{0,025} = 11$.
    Grossissement de l'objectif (inchangé) : $G_{\text{obj}} = \frac{\Delta}{f'_{\text{obj}}} = 40$.
    \[
    G_{\text{confort}} = G_{\text{obj}} \times G_{\text{oc,confort}} = 40 \times 11 = \mathbf{440}.
    \]
    (Alternative : $G_{\text{confort}} = G_c \left(1 + \frac{f'_{\text{oc}}}{D}\right) = 400 \times (1 + 0,1) = 440$.)
\end{enumerate}
\end{corrige}

\begin{corrige}[Exercice 5 — Choix de l'objectif pour un grossissement donné]
Données : $G_c = 600$, $\Delta = 0,16\,\text{m}$, $P_{\text{oc}} = 50\,\delta$, $D = 0,25\,\text{m}$.

\begin{enumerate}
    \item Puissance de l'objectif :
    \[
    G_c = \Delta \cdot P_{\text{obj}} \cdot D \cdot P_{\text{oc}} \implies P_{\text{obj}} = \frac{G_c}{\Delta \cdot D \cdot P_{\text{oc}}} = \frac{600}{0,16 \times 0,25 \times 50} = \frac{600}{2} = \mathbf{300\,\delta}.
    \]
    \item Focale image de l'objectif :
    \[
    f'_{\text{obj}} = \frac{1}{P_{\text{obj}}} = \frac{1}{300} \approx 3,33\times10^{-3}\,\text{m} = \mathbf{3,33\,\text{mm}}.
    \]
    \item Objectifs standards (pour $\Delta = 160\,\text{mm}$) : $4\times$, $10\times$, $40\times$, $100\times$.
    Le grossissement de l'objectif seul est $G_{\text{obj}} = \frac{\Delta}{f'_{\text{obj}}} = \frac{160}{3,33} \approx 48$.
    Les standards correspondent à $G_{\text{obj}} = 4, 10, 40, 100$.
    La valeur $48$ est la plus proche de $40$ (écart $8$) que de $100$ (écart $52$).
    \textbf{Choix : l'objectif $40\times$ (focale $4\,\text{mm}$, puissance $250\,\delta$).}
    \emph{Justification :} C'est l'objectif standard dont le grossissement ($40\times$) est le plus proche du grossissement requis ($48\times$). L'utilisation de cet objectif donnera un grossissement commercial $G_c = 40 \times \frac{0,25}{0,02} = 500$, légèrement inférieur à la cible $600$, mais c'est le choix standard le plus pertinent. Un objectif $100\times$ donnerait $G_c = 1250$, trop élevé.
\end{enumerate}
\end{corrige}

\begin{corrige}[Exercice 6 — Passage vision normale $\rightarrow$ vision au confort]
Données : $f'_{\text{oc}} = 25\,\text{mm}$, Objectif $40\times$ ($\Delta = 160\,\text{mm} \Rightarrow f'_{\text{obj}} = 4\,\text{mm}$).

\begin{enumerate}
    \item En vision normale, l'image intermédiaire $A'B'$ se trouve \textbf{au foyer objet $F_{\text{oc}}$ de l'oculaire} (distance $\overline{\text{Oc}A'B'} = -f'_{\text{oc}} = -25\,\text{mm}$).
    \item Pour la vision au confort (image finale à $D = -250\,\text{mm}$), l'objet de l'oculaire ($A'B'$) doit se trouver \textbf{entre $F_{\text{oc}}$ et l'oculaire} (donc \textbf{après} $F_{\text{oc}}$ dans le sens de propagation de la lumière).
    Calcul de la position : $\frac{1}{f'_{\text{oc}}} = \frac{1}{-D} - \frac{1}{p'_{\text{oc}}} \Rightarrow p'_{\text{oc}} = -\frac{D f'_{\text{oc}}}{D+f'_{\text{oc}}} = -\frac{250 \times 25}{275} \approx -22,73\,\text{mm}$.
    Distance par rapport à $F_{\text{oc}}$ : $\overline{F_{\text{oc}}A'B'} = p'_{\text{oc}} + f'_{\text{oc}} \approx +2,27\,\text{mm} > 0$.
    \item Pour déplacer $A'B'$ vers la droite (vers l'oculaire, $\Delta v_{\text{obj}} > 0$), il faut \textbf{éloigner l'objectif de la lame (vers l'oculaire)}.
    \emph{Justification (conjugaison objectif) :} L'objet (lame) est fixe. La relation de Newton $x_{\text{obj}} x'_{\text{obj}} = {f'_{\text{obj}}}^2$ avec $x'_{\text{obj}} = \overline{F'_{\text{obj}}A'B'}$.
    En vision normale : $x'_{\text{obj,n}} = \Delta = 160\,\text{mm}$.
    En vision confort : $x'_{\text{obj,c}} = \Delta + 2,27 = 162,27\,\text{mm}$.
    $x'_{\text{obj}}$ augmente. Comme $f'_{\text{obj}}$ constant, $x_{\text{obj}} = \frac{{f'_{\text{obj}}}^2}{x'_{\text{obj}}}$ \textbf{diminue}.
    $x_{\text{obj}} = \overline{F_{\text{obj}}\text{Objet}}$ diminue $\Rightarrow$ l'objet se rapproche du foyer objet $F_{\text{obj}}$.
    Comme l'objet est fixe, c'est l'objectif qui doit se rapprocher de l'objet $\Rightarrow$ \textbf{déplacement vers l'objet (vers le haut)}.
    \emph{Note :} L'approximation « grossissement constant » ($G_{\text{obj}} \approx \text{cste}$) conduit à $|u| = v/G$ ; $v$ augmente $\Rightarrow |u|$ augmente $\Rightarrow$ objectif vers l'oculaire. Cette approximation est fausse ici car $G_{\text{obj}}$ varie de $40$ à $40,6$. Le calcul exact (ci-dessus) montre le sens inverse.
    \item Déplacement approximatif de l'objectif (hypothèse $G_{\text{obj}} \approx \text{cste}$ demandée) :
    Déplacement de $A'B'$ : $\delta x' \approx \frac{{f'_{\text{oc}}}^2}{D} = \frac{25^2}{250} = 2,5\,\text{mm}$.
    Déplacement de l'objectif (ou de l'objet) : $\delta x \approx \frac{\delta x'}{G_{\text{obj}}} = \frac{2,5}{40} = 0,0625\,\text{mm} = \mathbf{62,5\,\mu\text{m}}$.
    (Sens : vers l'oculaire selon cette approximation).
\end{enumerate}
\end{corrige}

\begin{corrige}[Exercice 7 — Analyse des inscriptions sur un objectif]
Gravure : \textbf{Plan 40/0,65 $\infty$/0,17}

\begin{enumerate}
    \item \texttt{Plan} : Objectif \textbf{planachromatique} (correction de la planéité du champ : image nette sur tout le champ, pas seulement au centre).
    \texttt{40} : Grossissement transverse de l'objectif $G_{\text{obj}} = 40\times$ (pour la longueur de tube de référence).
    \texttt{0,65} : \textbf{Ouverture numérique} $\text{NA} = n \sin\alpha = 0,65$ (objectif « sec », sans immersion, $n=1$).
    \texttt{$\infty$} : Objectif \textbf{corrigé à l'infini} (infinity-corrected). L'objectif forme une image à l'infini (faisceaux émergents parallèles) ; une \textbf{lentille tube} (intégrée au corps du microscope) forme l'image intermédiaire réelle.
    \texttt{0,17} : Épaisseur de la lame de couvre-objet pour laquelle la correction sphérique est optimale ($0,17\,\text{mm}$).
    \item \textbf{Conséquence sur $\Delta$ :} La notion d'intervalle inter-foculaire $\Delta = \overline{F'_{\text{obj}}F_{\text{oc}}}$ \textbf{n'est plus définie} de la même manière car l'objectif n'a plus de foyer image réel accessible (image à l'infini). La distance entre l'objectif et l'oculaire est fixée par la lentille tube (souvent $f'_{\text{tube}} = 200\,\text{mm}$).
    \textbf{Mise au point :} Elle se fait généralement en déplaçant la \textbf{platine (l'objet)} par rapport à l'objectif (mise au point macrométrique/micrométrique), l'objectif et la lentille tube restant fixes par rapport au corps du microscope.
    \item Calcul de $f'_{\text{obj}}$ :
    Pour un objectif corrigé à l'infini, $G_{\text{obj}} = \frac{f'_{\text{tube}}}{f'_{\text{obj}}}$.
    $40 = \frac{200\,\text{mm}}{f'_{\text{obj}}} \implies f'_{\text{obj}} = \frac{200}{40} = \mathbf{5\,\text{mm}}$.
    (Puissance $P_{\text{obj}} = 200\,\delta$).
\end{enumerate}
\end{corrige}

\begin{corrige}[Exercice 8 — Pouvoir séparateur et immersion]
Données : $\text{NA} = 1,3$ (huile $n=1,515$), $\lambda = 500\,\text{nm} = 5\times10^{-7}\,\text{m}$.
Critère de Rayleigh : $\delta_{\text{min}} = 0,61 \frac{\lambda}{\text{NA}}$.

\begin{enumerate}
    \item $\delta_{\text{min}} = 0,61 \times \frac{500\,\text{nm}}{1,3} = \frac{305}{1,3} \approx \mathbf{234,6\,\text{nm}} \approx \mathbf{0,235\,\mu\text{m}}$.
    \item Bactérie $0,5\,\mu\text{m} = 500\,\text{nm}$. Comme $500\,\text{nm} > \delta_{\text{min}} (235\,\text{nm})$, \textbf{oui}, la bactérie est résolue (ses détails internes de taille $\sim 500\,\text{nm}$ sont visibles).
    \item Deux points à $200\,\text{nm}$. Comme $200\,\text{nm} < \delta_{\text{min}} (235\,\text{nm})$, \textbf{non}, ils ne sont pas distingués (image unique floue).
    \item Sans huile (air, $n=1$), $\text{NA}_{\text{air}} \approx 0,95$ (maximum théorique $\sin\alpha \le 1$).
    $\delta_{\text{min, air}} = 0,61 \times \frac{500}{0,95} \approx \mathbf{321\,\text{nm}}$.
    \textbf{Commentaire :} L'immersion à l'huile ($n=1,515$) permet d'augmenter l'ouverture numérique de $0,95$ à $1,3$ (facteur $1,37$), ce qui \textbf{réduit le pouvoir séparateur} (meilleure résolution) d'un facteur $1,37$ ($321 \to 235\,\text{nm}$). C'est crucial pour observer des détails submicroniques (bactéries, organites).
\end{enumerate}
\end{corrige}

\begin{corrige}[Exercice 9 — Probatoire Type 2022 : Microscope en vision normale]
\textbf{Barème indicatif :} 6 points au total.
\begin{itemize}
    \item Q1, Q2, Q3, Q4 : 0,5 pt chacun (calcul + unité + résultat).
    \item Q5 (Schéma) : 4 pts (lentilles, foyers, $\Delta$, objets/images, construction géométrique 2 rayons, rayons émergents parallèles, légendes).
\end{itemize}
\textbf{Points de vigilance :} Unités en mm pour focales, $\delta$ pour puissances. $G_c$ sans unité. Schéma : échelles respectées (qualitativement), flèches sur les rayons, sens de la lumière (gauche $\to$ droite), image intermédiaire renversée.

\begin{enumerate}
    \item $f'_{\text{obj}} = \frac{1}{P_{\text{obj}}} = \frac{1}{250} = 0,004\,\text{m} = \mathbf{4\,\text{mm}}$.
    \item $P_{\text{oc}} = \frac{1}{f'_{\text{oc}}} = \frac{1}{0,025} = \mathbf{40\,\delta}$.
    \item $G_c = \frac{\Delta}{f'_{\text{obj}}} \cdot \frac{D}{f'_{\text{oc}}} = \frac{0,16}{0,004} \times \frac{0,25}{0,025} = 40 \times 10 = \mathbf{400}$.
    \item $P_c = \frac{G_c}{D} = \frac{400}{0,25} = \mathbf{1600\,\delta}$.
    \item \textbf{Schéma de principe} (voir figure ci-dessous).
\end{enumerate}

\begin{popfigure}
\begin{tikzpicture}[scale=0.8, >=Stealth, thick]
% Parameters (drawing units: cm)
\def\fobjd{1.0}    % f'_obj = 1cm
\def\focd{1.0}     % f'_oc = 1cm
\def\Deltad{4.0}   % Delta = 4cm
% Positions
\coordinate (Oobj) at (0,0);
\coordinate (Fobj) at (-\fobjd,0);
\coordinate (Fpobj) at (\fobjd,0);
\coordinate (Foc) at (\fobjd+\Deltad,0);
\coordinate (Ooc) at (\fobjd+\Deltad+\focd,0);
\coordinate (Fpoc) at (\fobjd+\Deltad+2*\focd,0);

% Object AB (height 0.8cm, placed at x = -1.3cm, i.e. before F_obj at -1.0)
\coordinate (A) at (-1.3, 0.8);
\coordinate (B) at (-1.3, 0);

% Magnification Obj: G_obj = -Delta/f'_obj = -4.0
% Image A'B' at Foc (x=5), height = -3.2
\coordinate (Ap) at (\fobjd+\Deltad, -3.2);
\coordinate (Bp) at (\fobjd+\Deltad, 0);

% Draw lenses
\draw[popblue, thick] (Oobj) ++(0,-1.8) -- ++(0,3.6) node[midway, above left] {Objectif};
\draw[poporange, thick] (Ooc) ++(0,-1.8) -- ++(0,3.6) node[midway, above right] {Oculaire};

% Draw focal points
\foreach \pos/\name/\col in {Fobj/$F_{\text{obj}}$/popblue, Fpobj/$F'_{\text{obj}}$/popblue, Foc/$F_{\text{oc}}$/poporange, Fpoc/$F'_{\text{oc}}$/poporange} {
    \fill[\col] (\pos) circle (2pt);
    \node[\col, below=2pt] at (\pos) {\name};
}

% Delta
\draw[<->, popdark] (\fobjd, -2.2) -- (\fobjd+\Deltad, -2.2) node[midway, below] {$\Delta$};

% Rays from A (top of object) - Objective
% Ray 1: Parallel to axis -> through F'_obj
\draw[popblue, ->] (A) -- (0,0.8) -- (Ap);
% Ray 2: Through center of objective (undeviated)
\draw[popblue, ->] (A) -- (Oobj) -- (Ap);
% Ray 3: Through F_obj -> exits parallel
\draw[popblue, ->] (A) -- (Fobj) -- ++(7,0);

% Intermediate image
\draw[popgreen, thick, ->] (Bp) -- (Ap) node[midway, right] {$A'B'$};

% Ocular rays (Vision normale: A' at F_oc)
% Ray from A' through center of ocular -> straight
\draw[poporange, ->] (Ap) -- (Ooc) -- ++(2.5, 3.2);
% Ray from A' parallel to axis -> through F'_oc
\draw[poporange, ->] (Ap) -- (\fobjd+\Deltad+\focd, -3.2) -- (Fpoc) -- ++(2.5, 3.2);

% Rays from B' (on axis) -> emerge parallel to axis
\draw[poporange, ->] (Bp) -- (Ooc) -- ++(2.5, 0);
\draw[poporange, ->] (\fobjd+\Deltad+\focd, 0.5) -- ++(2.5, 0); % representative parallel ray

% Final image at infinity label
\node[popdark, align=left, anchor=west] at (10.5, 2.5) {Image finale à l'infini\\(faisceaux émergents parallèles)};

% Axis
\draw[popline, ->] (-2,0) -- (11,0) node[right] {Axe optique};
\end{tikzpicture}
\legende{Schéma de principe du microscope en vision normale. Construction de l'image intermédiaire $A'B'$ (rayons bleus) et rayons émergents parallèles de l'oculaire (rayons orange).}
\end{popfigure}
\end{corrige}

\begin{corrige}[Exercice 10 — Probatoire Type 2020 : Conception d'un microscope]
\textbf{Barème indicatif :} 5 points.
\begin{itemize}
    \item Q1, Q2 : 1 pt chacun (formule, calcul, unité).
    \item Q3 : 1,5 pt (calcul $G_{\text{obj}}$, comparaison, choix justifié).
    \item Q4 : 1,5 pt (réponse + justification optique : myope voit net au point lointain fini, image à l'infini floue sans correction).
\end{itemize}

\begin{enumerate}
    \item $G_c = \Delta \cdot P_{\text{obj}} \cdot D \cdot P_{\text{oc}} \implies P_{\text{obj}} = \frac{G_c}{\Delta \cdot D \cdot P_{\text{oc}}} = \frac{600}{0,16 \times 0,25 \times 50} = \frac{600}{2} = \mathbf{300\,\delta}$.
    \item $f'_{\text{obj}} = \frac{1}{P_{\text{obj}}} = \frac{1}{300}\,\text{m} \approx \mathbf{3,33\,\text{mm}}$.
    \item Grossissement objectif standard pour $\Delta = 160\,\text{mm}$ : $G_{\text{obj}} = \frac{\Delta}{f'_{\text{obj}}} = \frac{160}{3,33} \approx 48$.
    Standards : $4\times\ (f'=40\,\text{mm})$, $10\times\ (f'=16\,\text{mm})$, $40\times\ (f'=4\,\text{mm})$, $100\times\ (f'=1,6\,\text{mm})$.
    La valeur $48$ est la plus proche de $40$.
    \textbf{Choix : Objectif $40\times$.}
    \emph{Justification :} C'est l'objectif standard dont le grossissement ($40\times$) est le plus proche du grossissement requis ($48\times$). Il donnera $G_c = 40 \times 12,5 = 500$ (au lieu de $600$), ce qui est acceptable. L'objectif $100\times$ donnerait $1250$, trop fort.
    \item \textbf{Oui, il doit porter ses lunettes (verres divergents).}
    \emph{Justification :} Un myope non corrigé a un point lointain à distance finie (il ne voit net que les objets proches). En vision normale, le microscope fournit une image finale \textbf{à l'infini}. Sans correction, l'œil myope ne peut pas mettre au point une image à l'infini (elle se formerait en avant de la rétine). Les lunettes divergentes décalent le point lointain de l'œil à l'infini, permettant la vision nette de l'image finale.
\end{enumerate}
\end{corrige}

\begin{corrige}[Exercice 11 — Probatoire Type : Mise au point, Schéma et Résolution]
\textbf{Barème indicatif :} 8 points.
\begin{itemize}
    \item Q1 : 1 pt ($G_c$, $P_c$).
    \item Q2 : 2,5 pts (positions, sens déplacement, justification conjugaison).
    \item Q3 : 1 pt ($\delta_{\text{min}}$).
    \item Q4 : 1 pt (résolution 300 nm et 800 nm).
    \item Q5 : 2,5 pts (nature image, sens vis, rayons émergents).
\end{itemize}
\textbf{Points de vigilance :} Distinction vision normale/confort (position $A'B'$ par rapport à $F_{\text{oc}}$). Sens du déplacement : justification rigoureuse via relation de conjugaison (Newton ou Descartes) et non approximation. Unités (nm, $\delta$). Schéma : rayons émergents convergents (image réelle).

Données : $f'_{\text{obj}} = 4\,\text{mm}$, $f'_{\text{oc}} = 25\,\text{mm}$, $\Delta = 160\,\text{mm}$, $\text{NA}=0,65$, $\lambda=550\,\text{nm}$, $D=250\,\text{mm}$.

\begin{enumerate}
    \item $G_c = \frac{\Delta}{f'_{\text{obj}}} \cdot \frac{D}{f'_{\text{oc}}} = \frac{160}{4} \times \frac{250}{25} = 40 \times 10 = \mathbf{400}$.
    $P_c = \frac{G_c}{D} = \frac{400}{0,25} = \mathbf{1600\,\delta}$.
    \item \textbf{Mise au point vision confort :}
    \begin{enumerate}
        \item Vision normale : $A'B'$ au \textbf{foyer objet $F_{\text{oc}}$} de l'oculaire ($\overline{\text{Oc}A'B'} = -f'_{\text{oc}} = -25\,\text{mm}$).
        \item Vision au confort : Image finale à $\overline{\text{Oc}A''B''} = -D = -250\,\text{mm}$.
        Relation de Descartes (oculaire) : $\frac{1}{f'_{\text{oc}}} = \frac{1}{-D} - \frac{1}{\overline{\text{Oc}A'B'}}$.
        $\frac{1}{25} = -\frac{1}{250} - \frac{1}{p'} \implies \frac{1}{p'} = -\frac{1}{250} - \frac{1}{25} = -\frac{11}{250}$.
        $p' = \overline{\text{Oc}A'B'} = -\frac{250}{11} \approx \mathbf{-22,73\,\text{mm}}$.
        Distance algébrique par rapport à $F_{\text{oc}}$ : $\overline{F_{\text{oc}}A'B'} = p' - (-f'_{\text{oc}}) = -22,73 + 25 = \mathbf{+2,27\,\text{mm}}$.
        $A'B'$ est \textbf{après $F_{\text{oc}}$} (entre $F_{\text{oc}}$ et l'oculaire).
        \item Pour amener $A'B'$ de $F_{\text{oc}}$ (pos $-25\,\text{mm}$) vers $-22,73\,\text{mm}$ (vers la droite, vers l'oculaire), il faut \textbf{augmenter $v_{\text{obj}}$} (distance objectif $\to$ image intermédiaire).
        Objet (lame) fixe. Relation de Newton objectif : $x_{\text{obj}} x'_{\text{obj}} = {f'_{\text{obj}}}^2$ avec $x'_{\text{obj}} = \overline{F'_{\text{obj}}A'B'} = v_{\text{obj}} - f'_{\text{obj}}$.
        $x'_{\text{obj}}$ augmente ($\Delta \to \Delta+2,27$). $f'_{\text{obj}}$ constant $\Rightarrow x_{\text{obj}} = \overline{F_{\text{obj}}\text{Objet}}$ \textbf{diminue}.
        L'objet se rapproche du foyer objet $F_{\text{obj}}$. Comme l'objet est fixe, l'objectif doit se rapprocher de l'objet.
        \textbf{Sens : vers l'objet (vers le haut / rapprocher l'objectif de la lame).}
    \end{enumerate}
    \item Pouvoir séparateur (Rayleigh) :
    $\delta_{\text{min}} = 0,61 \frac{\lambda}{\text{NA}} = 0,61 \times \frac{550\,\text{nm}}{0,65} = \frac{335,5}{0,65} \approx \mathbf{516\,\text{nm}} \approx \mathbf{0,52\,\mu\text{m}}$.
    \item \textbf{Résolution :}
    \begin{itemize}
        \item Espacement $300\,\text{nm} < \delta_{\text{min}} (516\,\text{nm})$ : \textbf{Non}, les deux globules ne sont pas résolus (image confondue).
        \item Espacement $800\,\text{nm} > \delta_{\text{min}} (516\,\text{nm})$ : \textbf{Oui}, les deux globules sont résolus (détails distincts).
    \end{itemize}
    \item \textbf{Analyse du schéma incomplet (image intermédiaire avant $F_{\text{oc}}$) :}
    \begin{enumerate}
        \item L'image intermédiaire $A'B'$ est située \textbf{avant le foyer objet $F_{\text{oc}}$} (côté objectif, $|p'| > f'_{\text{oc}}$). L'oculaire fonctionne alors comme une loupe en \textbf{vision au confort (ou distance finie)} : il forme une image finale \textbf{virtuelle}, \textbf{droite} par rapport à $A'B'$, située \textbf{à distance finie} (entre $F'_{\text{oc}}$ et l'oculaire, ou au-delà de $F'_{\text{oc}}$ selon la position exacte de $A'B'$). Ici, comme $A'B'$ est avant $F_{\text{oc}}$, l'image finale est virtuelle du même côté que $A'B'$.
        \item Pour ramener $A'B'$ \textbf{au foyer $F_{\text{oc}}$} (vision normale), il faut déplacer $A'B'$ vers la \textbf{droite} (vers l'oculaire). Cela nécessite d'\textbf{augmenter $v_{\text{obj}}$}. Comme l'objet est fixe, d'après la relation de Newton ($x_{\text{obj}} x'_{\text{obj}} = {f'_{\text{obj}}}^2$), augmenter $x'_{\text{obj}}$ impose de diminuer $x_{\text{obj}}$, donc de \textbf{rapprocher l'objectif de l'objet (lame)}.
        \textbf{Sens de rotation de la vis macrométrique :} Dans le sens qui \textbf{rapproche l'objectif de la platine} (généralement le sens horaire vu du côté de la vis, mais dépend du montage ; on précise : « tourner pour faire descendre le tube / monter la platine / rapprocher objectif et objet »).
        \item \textbf{Rayons émergents pour $B'$ :} $B'$ est sur l'axe optique. Tous les rayons issus de $B'$ et passant par l'oculaire émergent \textbf{parallèles à l'axe optique} (faisceau cylindrique centré sur l'axe).
    \end{enumerate}
\end{enumerate}
\end{corrige}

\newpage

\coursec{Fiche bilan}

\begin{popfigure}
\begin{tikzpicture}[mindmap, grow cyclic, every node/.style={concept, text=white, font=\small\bfseries},
  root concept/.append style={concept color=popblue, font=\large\bfseries, minimum size=3.2cm},
  level 1/.append style={level distance=4.4cm, sibling angle=51.4},
  level 1 concept/.append style={minimum size=2.6cm, font=\footnotesize\bfseries}]
\node[root concept]{Le microscope}
  child[concept color=poporange]{ node {Description et r\^ole} }
  child[concept color=popgreen]{ node {Sch\'ema de principe} }
  child[concept color=poppurple]{ node {Puissance $P=\dfrac{\alpha_i}{AB}$} }
  child[concept color=poppink]{ node {Grossissement $G=\dfrac{\alpha_i}{\alpha_0}$} }
  child[concept color=popgold]{ node {Mise au point} }
  child[concept color=popblue!70]{ node {Tube optique $\Delta$} }
  child[concept color=popmuted]{ node {Limites et erreurs} };
\end{tikzpicture}
\legende{Carte mentale de la le\c{c}on : les sept id\'ees forces du microscope.}
\end{popfigure}

\begin{aretenir}[L'essentiel de la le\c{c}on : le microscope]
\textbf{D\'efinitions cl\'es}
\begin{itemize}
  \item \cle{Microscope} : instrument d'optique form\'e de deux syst\`emes convergents coaxiaux --- l'\cle{objectif} (c\^ot\'e objet, tr\`es courte focale $f'_1$ de l'ordre du millim\`etre) et l'\cle{oculaire} (c\^ot\'e \oe il, focale $f'_2$ de quelques centim\`etres) --- qui donne d'un objet proche et tr\`es petit une image fortement agrandie.
  \item \cle{Image interm\'ediaire} $A_1B_1$ : image r\'eelle, renvers\'ee et agrandie, donn\'ee par l'objectif ; elle se forme au voisinage du foyer objet $F_2$ de l'oculaire et sert d'objet pour celui-ci.
  \item \cle{Intervalle optique} (ou longueur du tube optique) : $\Delta = \overline{F'_1F_2}$, de l'ordre de \SI{16}{cm} pour un microscope standard.
  \item \cle{Puissance} $P$ : quotient de l'angle $\alpha_i$ sous lequel on voit l'image \`a travers l'instrument par la dimension de l'objet : $P = \dfrac{\alpha_i}{AB}$ ; unit\'e SI : la \cle{dioptrie} ($\delta$), $1\,\delta = \SI{1}{m^{-1}}$.
  \item \cle{Grossissement} $G$ : rapport de l'angle $\alpha_i$ (image vue \`a travers l'instrument) \`a l'angle $\alpha_0$ sous lequel l'objet est vu \`a l'\oe il nu \`a la distance minimale de vision distincte $d_m = \SI{25}{cm}$ : $G = \dfrac{\alpha_i}{\alpha_0}$ (sans unit\'e).
  \item \cle{Mise au point} : r\'eglage de la distance objectif--objet pour que l'image d\'efinitive $A'B'$ soit nette pour l'observateur (au punctum remotum, en g\'en\'eral \`a l'infini pour un \oe il emm\'etrope qui n'accommode pas).
\end{itemize}

\textbf{Formules \`a conna\^itre par c\oe ur}
\begin{center}
\begin{tabularx}{\linewidth}{|l|X|l|}
\hline
\rowcolor{popblueL} \textbf{Grandeur} & \textbf{Expression} & \textbf{Unit\'e} \\
\hline
Puissance intrins\`eque & $P_i = \dfrac{\Delta}{f'_1 f'_2}$ (image \`a l'infini) & dioptrie ($\delta$) \\
\hline
Puissance commerciale & $P_c = \dfrac{\alpha_i}{AB}$ avec $\alpha_i$ pour image \`a \SI{25}{cm} & dioptrie ($\delta$) \\
\hline
Grossissement commercial & $G_c = \dfrac{\alpha_i}{\alpha_0} = d_m \cdot P_c = \dfrac{P_c}{4}$ ($d_m = \SI{0,25}{m}$) & sans unit\'e \\
\hline
Grandissement de l'objectif & $|\gamma_1| = \left|\dfrac{\overline{O_1A_1}}{\overline{O_1A}}\right| \approx \dfrac{\Delta}{f'_1}$ & sans unit\'e \\
\hline
Relation puissance--grossissement & $G = d_m \times P$ & --- \\
\hline
\end{tabularx}
\end{center}

\textbf{M\'ethodes express}
\begin{itemize}
  \item \cle{Construction g\'eom\'etrique} : (1) tracer l'image $A_1B_1$ donn\'ee par l'objectif (rayon issu de $B$ parall\`ele \`a l'axe $\to$ \'emerge en passant par $F'_1$ ; rayon passant par $O_1$ non d\'evi\'e) ; (2) placer $A_1B_1$ au voisinage du foyer objet $F_2$ de l'oculaire ; (3) si $A_1$ est en $F_2$, l'image d\'efinitive part \`a l'infini : les rayons \'emergents de l'oculaire sont parall\`eles entre eux ; $\alpha_i$ est l'angle de ces rayons avec l'axe.
  \item \cle{Calcul type} : convertir toutes les distances en m\`etres ; calculer $P_i = \Delta/(f'_1 f'_2)$ ; en d\'eduire $G = P/4$ si la vision \`a \SI{25}{cm} est demand\'ee ; v\'erifier l'ordre de grandeur ($G$ de quelques dizaines \`a quelques centaines).
  \item \cle{Mise au point pratique} : placer la pr\'eparation, r\'egler l'\'eclairage (miroir, diaphragme), approcher l'objectif au maximum en regardant de profil, puis \'eloigner lentement le tube (vis macrom\'etrique) en observant dans l'oculaire jusqu'\`a nettet\'e ; affiner avec la vis microm\'etrique. Ne jamais descendre l'objectif l'\oe il \`a l'oculaire (risque d'\'ecraser la lame).
  \item \cle{Analyse dimensionnelle} : $[P_i] = \dfrac{\text{m}}{\text{m}\times\text{m}} = \text{m}^{-1}$ : la dioptrie est bien homog\`ene \`a l'inverse d'une longueur.
\end{itemize}
\end{aretenir}

\begin{attention}[Pi\`eges fr\'equents au Probatoire]
\begin{itemize}
  \item Ne pas confondre \cle{puissance} (en dioptries, $\text{m}^{-1}$) et \cle{grossissement} (sans unit\'e) : ce sont deux grandeurs de natures diff\'erentes, reli\'ees par $G = d_m \times P$.
  \item Dans $G_c = \dfrac{P_c}{4}$, le facteur $4$ vient de $d_m = \SI{0,25}{m}$ exprim\'e en \textbf{m\`etres} : $G_c = 0,25 \times P_c = \dfrac{P_c}{4}$. Oublier la conversion est l'erreur la plus fr\'equente.
  \item Les angles $\alpha_i$ et $\alpha_0$ doivent \^etre exprim\'es en \textbf{radians} dans les formules $P = \dfrac{\alpha_i}{AB}$ et $G = \dfrac{\alpha_i}{\alpha_0}$.
  \item L'image interm\'ediaire $A_1B_1$ est \textbf{r\'eelle} (elle sert d'objet r\'eel pour l'oculaire) ; l'image d\'efinitive $A'B'$ est \textbf{virtuelle} lorsqu'elle est rejet\'ee \`a l'infini.
  \item L'approximation $|\gamma_1| \approx \dfrac{\Delta}{f'_1}$ suppose l'objet $A$ plac\'e au voisinage du foyer objet $F_1$ de l'objectif : citer cette condition dans la r\'edaction.
  \item Toujours pr\'eciser le signe dans les relations de conjugaison : le grandissement $\gamma_1$ de l'objectif est \textbf{n\'egatif} (image renvers\'ee), d'o\`u l'usage de sa valeur absolue.
\end{itemize}
\end{attention}

\begin{exemplebox}[Ordres de grandeur \`a retenir]
Un microscope de laboratoire scolaire typique : $f'_1 = \SI{4}{mm}$, $f'_2 = \SI{25}{mm}$, $\Delta = \SI{16}{cm} = \SI{0,16}{m}$.
\begin{align*}
P_i &= \dfrac{\Delta}{f'_1 f'_2} = \dfrac{0,16}{4\times 10^{-3} \times 25\times 10^{-3}} = \dfrac{0,16}{10^{-4}} = \SI{1600}{\delta} \\
G_c &= \dfrac{P_c}{4} \approx \dfrac{1600}{4} = 400
\end{align*}
On retrouve bien l'ordre de grandeur annonc\'e par les fabricants (grossissement commercial de quelques centaines).
\end{exemplebox}

\begin{aretenir}[Checklist avant le Probatoire]
\begin{itemize}
  \item[$\square$] Je sais nommer et situer sur un sch\'ema : objectif, oculaire, foyers $F_1$, $F'_1$, $F_2$, $F'_2$, intervalle optique $\Delta$.
  \item[$\square$] Je sais que l'objectif donne une image r\'eelle, renvers\'ee, agrandie, et que l'oculaire joue le r\^ole de loupe.
  \item[$\square$] Je sais construire proprement l'image interm\'ediaire $A_1B_1$ puis l'image d\'efinitive \`a l'infini.
  \item[$\square$] Je connais la condition de mise au point : $A_1$ au voisinage de $F_2$ pour une observation sans accommodation.
  \item[$\square$] Je connais par c\oe ur $P_i = \dfrac{\Delta}{f'_1 f'_2}$ et je sais la justifier par l'analyse dimensionnelle.
  \item[$\square$] Je sais que la puissance s'exprime en dioptries ($1\,\delta = \SI{1}{m^{-1}}$).
  \item[$\square$] Je connais la d\'efinition du grossissement $G = \dfrac{\alpha_i}{\alpha_0}$ avec $\alpha_0 = \dfrac{AB}{d_m}$.
  \item[$\square$] Je sais passer de la puissance au grossissement : $G_c = d_m \times P_c = \dfrac{P_c}{4}$ avec $d_m = \SI{0,25}{m}$.
  \item[$\square$] Je convertis syst\'ematiquement les focales et $\Delta$ en m\`etres avant tout calcul num\'erique.
  \item[$\square$] Je sais d\'ecrire le protocole de mise au point au microscope (vis macrom\'etrique puis microm\'etrique, jamais en descendant l'objectif l'\oe il \`a l'oculaire).
  \item[$\square$] Je v\'erifie mes r\'esultats : unit\'e SI pr\'esente, ordre de grandeur r\'ealiste, chiffres significatifs raisonnables.
  \item[$\square$] Je sais citer une application concr\`ete : laboratoires hospitaliers (diagnostic du paludisme au Cameroun), contr\^ole qualit\'e, biologie.
\end{itemize}
\end{aretenir}

\findecours

\end{document}
